\documentclass[11pt]{article}
\usepackage{amsmath,amssymb,amsthm,hyperref}
\title{Toward an Unconditional Navier--Stokes Closure:\\A Leverage-Layer Research Program}
\author{Ioannis Tsiokos}
\date{Research draft, 4 October 2026}
\newtheorem{proposition}{Proposition}
\begin{document}
\maketitle
\begin{abstract}
We organize a fresh constructive program for the unforced whole-space
Navier--Stokes existence and smoothness problem. The existing spectral
continuation criterion is a translated conditional layer. We isolate its
unresolved original-equation producer, independently prove elementary
absorption and weighted-tail laws, and construct written conditional
continuation and full initial-face/pressure return adapters. The all-data
source estimates remain open. This draft does not establish unconditional closure.
\end{abstract}
\section{Original target}
For every viscosity $\nu>0$ and divergence-free Schwartz datum on
$\mathbb R^3$, seek globally smooth velocity and pressure on the closed
half-space, satisfying the original unforced equation and bounded kinetic
energy. This is statement (A) of Fefferman's
\href{https://www.claymath.org/wp-content/uploads/2022/06/navierstokes.pdf}{official Clay problem statement}.
\section{Translated criterion}
For $\gamma>5/2$, write $g(t)=\langle\xi\rangle^\gamma\widehat u(t)$ and
$P_N=\mathbf1_{|\xi|\le N}$. The unresolved condition is
\[
\begin{gathered}
\forall T>0\;\exists N\in\mathbb N,\;0\le q_T<1\;\forall t\in\mathcal T,\\
t\le T\Longrightarrow\|(I-P_N)g(t)\|_2^2\le q_T\|g(t)\|_2^2.
\end{gathered}
\]
The cutoff and fraction must be fixed over each entire preterminal horizon.
\section{Independent laws and unpaid applicability}
\begin{proposition}
If $X=L+H$, $L\le K$, $H\le qX$ and $q<1$, then
$X\le K/(1-q)$.
\end{proposition}
\begin{proof}
$(1-q)X\le L\le K$ and $1-q>0$.
\end{proof}
For $\delta>0$, multiplication by the original frequency weight gives
\[
\|(I-P_N)g\|_2^2\le(1+N^2)^{-\delta}
\int (1+|\xi|^2)^\delta|g(\xi)|^2\,d\xi.
\]
These laws do not produce terminal-uniform original Navier source bounds.
Assuming a bounded higher norm already assumes continuation-strength content.
\section{Recovered kinetic baseline and weighted gap}
For the original projected convection, define signed kinetic transfer
$\Phi_R=\operatorname{Re}\langle P_R\widehat u,
\widehat{\mathbb P((u\cdot\nabla)u)}\rangle$.
The earlier low-band energy identity and its high-band consequence give
\[
-E_{>R}(0)\le 2\int_0^t\Phi_R(s)\,ds\le E_{\le R}(0)\le E_0.
\]
This bounds net transfer, not its positive part or total variation.
With $w(r)=(1+r^2)^\gamma$ and finite $N<M$, Tonelli gives
\[
w(N)E_{>N}(t)+\int_N^M w'(r)E_{>r}(t)\,dr
=\int_{|\xi|>N}w(\min\{|\xi|,M\})|\widehat u(t,\xi)|^2\,d\xi.
\]
The nonnegative spectral expression converges monotonically to the strong
tail. Combining actual-path balances introduces the signed weighted
flux $w(N)\Phi_N+\int_N^M w'(r)\Phi_r\,dr$. The energy-only upper bound
grows as $E_0w(M)$; it cannot pay the infinite-cutoff weighted budget.
Real-radius continuity and finite-rectangle Fubini justify the finite-cutoff
balance. A uniform upper bound on cumulative signed flux, independent of
outer cutoff and existing time over each full finite preterminal horizon,
implies a uniform strong bound. Monotone convergence is needed only for
nonnegative spectral weights; no infinite signed-flux limit is required
for this inequality return. The uniform cumulative budget remains unproved.

\section{Conditional critical laws and original return}
For the same maximal strong solution set
$Z=\|\nabla u\|_2^2$, $Q=\|\Delta u\|_2^2$ and
$W=\|\widehat u\|_1$. A whole-space Fourier interpolation gives
\[
 W^2\le \frac{2}{\pi^2}\sqrt{ZQ}.
\]
Bounded frequency cutoffs, the original full convection product and
nonnegative monotone convergence yield
\[
 X_s(t)+\nu\int_0^tY_s
 \le X_s(0)+\frac{4c_s^2}{\nu}\int_0^t W^2X_s,
 \qquad c_s=2^{\max(s-1,0)}.
\]
Consequently a supplied finite full-endpoint $\int W^2$ implies continuation.
Enstrophy estimates also show that supplied finite doses
$\int\|u\|_5^5$ or $\int Z^2$ pay this Wiener dose. These are conditional
laws, not independently established arbitrary-data integrability results.

Writing $S=\operatorname{sym}\nabla u$, the positive middle strain eigenvalue
gives a further conditional interface. For ordered $a\le b\le c$ with
$a+b+c=0$, $-4abc\le2\max(b,0)(a^2+b^2+c^2)$.
The actual-path strain identity therefore bounds production by
$2\|\lambda_2^+(S)\|_\infty\|S\|_2^2$. A supplied finite time integral
of that coefficient pays bounded enstrophy and integrated $Q$, hence
continuation. Its all-data supplier remains open.

For original Schwartz input, high-order local mild construction, exact
weight lowering, same-path uniqueness and bounded-cutoff estimates prove
all Sobolev orders on each existing compact interval. Temporal induction
and Sobolev embedding give all mixed velocity derivatives through $t=0$.
The pressure multiplier
$\widehat p=-\xi_i\xi_j\widehat{u_i u_j}/|\xi|^2$ gives the same closed-face
smoothness in a coherent $L^2$ gauge. Globality, when supplied, then returns
all original Clay properties including the native kinetic bound.
Detailed written arguments are preserved in the round02 research dossiers;
the full PDE proof is not claimed to be mechanized.

\section{A signed heat-scale source gate}
Introduce an independent heat coordinate $\theta>0$, with
$H_\theta=e^{\theta\Delta}$, $u_\theta=H_\theta u$,
$S_\theta=\operatorname{sym}\nabla u_\theta$ and
$\tau_\theta=H_\theta(u\otimes u)-u_\theta\otimes u_\theta$.
The covariance is positive semidefinite and
$\int\operatorname{tr}\tau_\theta=E-E_\theta$.
Its exact scale equation and the original physical-time equation give
\[
\begin{split}
\tau_\theta&=2\int_0^\theta H_{\theta-s}
    [\nabla u_s(\nabla u_s)^T],ds,\\
\Pi_\theta&=-\int S_\theta:\tau_\theta
 =\langle H_{2\theta}u,(u\cdot\nabla)u\rangle,\\
E_\theta'+2\nu Z_\theta&=-2\Pi_\theta.
\end{split}
\]
The Gaussian stress-scale construction also appears in
\href{https://arxiv.org/pdf/1912.00293v1}{Johnson (2019), Section II};
here the derivation retains the original unforced whole-space carrier.
Covariance positivity does not sign its contraction with trace-free strain.
An explicit real solenoidal Schwartz field
\[
U=(4xz^2-2x,-2yz^2+y,-4x^2z+2y^2z+z)e^{-|x|^2}
\]
has exact transfer
\[
\Pi_\theta(U)=\frac{64\pi^{3/2}\theta(10\theta+3)}
 {(16\theta+3)^{11/2}}>0\qquad(\theta>0).
\]
The opposite field has the same quadratic observations and opposite
transfer. For $U_R=R^{3/2}U(Rx)$ and $\theta_R=c/R^2$, kinetic energy
is fixed, transfer grows as $R^{5/2}$ and the viscous complement as $R^2$.
Hence a universal instantaneous inequality
$\Pi_\theta\le a\nu(Z-Z_\theta)+C(E_0,\nu)(E-E_\theta)$ with finite
coefficients independent of heat scale and high-order initial shape fails
even at fixed viscosity. This is an initial-state control, not a singular
trajectory or a rejection of time-integrated phase information.

For integer $m>\gamma$, a bounded multiplier
\[
M_\epsilon(r)=\kappa_{\gamma m}\int_\epsilon^\infty
 \theta^{-\gamma-1}(1-e^{-8\pi^2\theta r^2})^m,d\theta
 \uparrow r^{2\gamma}
\]
has $\kappa_{\gamma m}=(8\pi^2)^{-\gamma}/
\int_0^\infty z^{-\gamma-1}(1-e^{-z})^m dz$.
The original energy balance reads
\[
\begin{split}
X_\epsilon(t)+2\nu\int_0^tD_\epsilon
 &=X_\epsilon(0)-2\int_0^t\mathcal P_\epsilon,\\
\mathcal P_\epsilon&=\kappa_{\gamma m}\int_\epsilon^\infty
 \theta^{-\gamma-1}\sum_{j=1}^m(-1)^j\binom mj\Pi_{j\theta}\,d\theta.
\end{split}
\]
Here $X_\epsilon=\int M_\epsilon|\widehat u|^2$ and
$D_\epsilon=\int4\pi^2r^2M_\epsilon|\widehat u|^2$.
The large-heat-scale part is independently paid by
$|\Pi_\theta|\le C_H\theta^{-5/4}E_0^{3/2}$.
A terminal-uniform one-sided cumulative bound for the remaining signed
small-scale part would give strong continuation by nonnegative monotone
convergence. That supplier remains open; the exact balance alone contains
the unknown terminal strong norm.

\section{Physical-time windows and the remaining cone source}
For a fixed positive width $h$ inside the initial local existence interval,
observe the same path along the heat characteristic
$\theta+\nu(t-s)$. Its original energy identity is
\[
\frac{d}{ds}E_{\theta+\nu(t-s)}(s)
 =-2\Pi_{\theta+\nu(t-s)}(s).
\]
Every shifted binomial term, including $j=0$, therefore gives
\[
\begin{split}
X_\epsilon(t)&=Y_{\epsilon,h}(t-h)-2\int_{t-h}^tW_\epsilon(s,t)\,ds,\\
W_\epsilon(s,t)&=\langle M_\epsilon(D)H_{2\nu(t-s)}u(s),B(s)\rangle,\\
0\le Y_{\epsilon,h}&\le
 \left(\frac{\gamma}{8\pi^2e\nu h}\right)^\gamma E_0.
\end{split}
\]
The past boundary is thus paid independently of the terminal strong norm.
Large base heat scales are paid by the native energy law and smoothing;
for small base scales, any fixed positive distance from $s=t$ is also
energy-paid. The remaining supplier is a uniform upper bound on the signed
source at both small heat scale and recent physical lag. It is open.

Converting heat differences into time integrals leaves bulk derivatives
and terminal production terms. For the explicit Schwartz
control, the cubic state derivative in the direction $\Delta U$ is
\[
T_\theta[\Delta U]
 =-\frac{2304\pi^{3/2}\theta(80\theta^2+52\theta+9)}
 {(16\theta+3)^{13/2}}<0\qquad(\theta>0).
\]
At original initial datum $bU$, the actual physical derivative is
$\nu b^3T_\theta+b^4Q_\theta$, where $Q_\theta$ retains full pressure
and is finite. Small positive $b$ makes that derivative nonzero, refuting
its omission. This initial-state control supplies no integrated source bound.

On a physical lag band $I=[t-2d,t-d]$, a whole-space Sobolev/stress estimate
gives, uniformly in cutoff,
\[
\left|-2\int_I W_\epsilon^{\rm small}\right|
 \le C_\gamma E_0^{3/4}\nu^{-\gamma-1/2}
       d^{-\gamma-1/4}\left(\int_I Z\right)^{3/4}.
\]
Its necessary dissipation price for positive source work $A$ is proportional
to $A^{4/3}d^{(4\gamma+1)/3}$. The algebraic pattern
$A=R^{2\gamma}$, $d=R^{-2}$ has summable price $R^{-2/3}$.
This tests the strength of the estimate; no realizing trajectory or
disjoint-event sequence is asserted. A stronger signed original source
constraint remains necessary for this route.

\section{Retained source structure and scoped estimate controls}
The fixed-window cone source is uniformly bounded if and only if the
strong norm is uniformly bounded, after the paid boundary and source
adapters. This equivalence supplies no independent estimate.
There is, however, an exact original quartic structure to retain.
Set $F=\mathbb P((u\cdot\nabla)u)$ and $H=H_{2\theta}$.
The quartic part of the characteristic heat-transfer derivative is
\[
Q=-\langle(H-I)u,(F\cdot\nabla)u\rangle
  +\langle[u\cdot\nabla,H]u,F\rangle.
\]
Pressure remains in $F$. Its first increment bound contains unpaid
$\|u\|_\infty$ and $\|\nabla u\|_\infty$ factors.

Native absolute interpolation gives lag weights $\tau^{-b}Z^a$ with
$a=3/2-3/q$, $b=-1/4+3/q$, $2\le q\le4$.
Since $a+b=5/4$, time H\"older against the native $L^1$ budget for $Z$
cannot pay this family at lag zero. Scalar budget concentrations test
this majorant's strength; they are not original solution trajectories.
At each finite cutoff, a large terminal excess produces a positive
source band by exact dyadic partition. Its selected lag index and overlap
with bands from other terminals remain uncontrolled. Cyclic triad
cancellation supplies an interaction gain on separated frequencies,
with a short reflected multiplier comparison, but no full source sum.
The signed finite-cutoff cap bound is finite and grows as
$\epsilon^{-\gamma}$; it cannot discharge the cutoff-uniform supplier.

\section{Original vorticity and two positive classes}
With $V_\epsilon=M_\epsilon^{\rm small}$, set
$L_\epsilon(r)=V_\epsilon(r)/(4\pi^2r^2)$, extended by zero at the origin.
It is bounded for each positive cutoff and gives exactly
\[
 \int L_\epsilon|\widehat\omega|^2
   =\int V_\epsilon|\widehat u|^2,\qquad \omega=\operatorname{curl}u.
\]
The original heat-window transport and stretching pairings sum to minus
the velocity source. This is an equivalent carrier, not an estimate.
For the original Gaussian $U$, the nonlinear contributions to
$\tfrac12\partial_t\|\Delta\omega\|_2^2$ are respectively
$\frac{23936}{81}(\pi/3)^{3/2}$ and
$\frac{35200}{81}(\pi/3)^{3/2}$. Their sum is
$-\langle(-\Delta)^3U,B\rangle$; the signed velocity Fourier pairing
has the opposite sign and a factor $(2\pi)^{-6}$.

For axisymmetric no-swirl Schwartz input, original uniqueness preserves
the symmetry and $\eta=\omega_\theta/r$ satisfies
\[
 D\eta=\nu(\partial_{rr}+3r^{-1}\partial_r+\partial_{zz})\eta.
\]
Maximum and physical $L^1$ Kato estimates retain the favorable axis sink
and pay $A=\|\eta_0\|_\infty$, $B_1=\|\eta_0\|_1$ independently.
An actual $\mathbb R^3$ Biot--Savart angular calculation gives
\[
 \|u_r/r\|_\infty\le
 \frac9{2(2\pi)^{1/3}}A^{2/3}B_1^{1/3}.
\]
The isolated azimuthal strain eigenvalue bounds the positive middle
eigenvalue by $|u_r/r|$, paying the existing continuation return on this
strict class. Compare the scalar budgets and exponents with
\href{https://cml.centre-mersenne.org/item/10.5802/cml.25.pdf}{Gallay--\v{S}ver\'ak (2015)}.
The relevant statements are Proposition 2.6 and Lemma 5.1; their global
theorem is not a premise of this project-native proof.

A complementary viscosity-dependent class follows directly from the
native enstrophy law. With $E_0=\|u_0\|_2^2$, $Z_0=\|\nabla u_0\|_2^2$
and $S=4/\sqrt3$, the spatial estimate is
\[
 |\langle\Delta u,B\rangle|
 \le S^{3/2}(E_0Z)^{1/4}Q.
\]
If $E_0Z_0\le27\nu^4/65536$, a first-hit bootstrap gives
$Z+\nu\int Q\le Z_0$. The resulting finite Wiener dose proves original
continuation and the existing full closed-face return. This is a small
initial-product class, not an arbitrary-data theorem. Both positive
classes are known calibration mechanisms; no novelty is asserted.

\section{Directional source, actual events and bending}
Define $\alpha=\omega\cdot S\omega/|\omega|^2$ where $\omega\ne0$,
and zero otherwise. The original curl equation gives
\[
 Z'+2\nu Q\le2\|\alpha_+\|_\infty Z.
\]
The Biot--Savart contraction retains the perpendicular vorticity through
\[
 k(x,y)=\frac{|\omega(y)\times n(x)|}{|x-y|},
 \qquad n(x)=\frac{\omega(x)}{|\omega(x)|}.
\]
At fixed physical radius $R$ its remote sector is independently bounded by
$39R^{-5/2}\sqrt{E_0}/(2\sqrt{5\pi})$.
The local integral of $3k/(4\pi|x-y|^2)$ remains unpaid in general;
in no swirl, $k\le|\eta(y)|$ recovers the class payment.

A hypothetical finite maximal endpoint forces unbounded $Z$ by the
existing squared-enstrophy criterion. First upper-level hits and last
lower-level crossings construct disjoint actual intervals $J_k$ with
$z_k\le Z\le2z_k$, $z_k=2^kZ_0$, and
\[
 \int_{J_k}\|\alpha_+\|_\infty\ge\frac{\log2}2,
 \qquad |J_k|\ge\frac{3\nu^3}{8C_Zz_k^2},
 \qquad \int_{J_k}Z\ge\frac{3\nu^3}{8C_Zz_k}.
\]
The native lower prices are summable. This proves actual conditional
event accounting, not an independent upper stretching dose or blowup.

For the smooth zero-safe bending carrier
$J_k=\omega\times\partial_k\omega$, $K=\sum_k|J_k|^2$,
the original equation retains co-area deformation, index mixing,
$\nabla A$ production and a viscous cross term, $A_{ij}=\partial_j u_i$.
In no swirl $K=r^2\eta^4$ is generally nonzero, and exact reduction gives
\[
\begin{split}
 (D-\nu\Delta)K={}&2(u_r/r)K
 -4\nu(\eta^2+r\eta\eta_r)^2\\
 &-8\nu r^2\eta^2\eta_r^2-12\nu r^2\eta^2\eta_z^2.
\end{split}
\]
General pointwise negativity of the viscous bending source is false:
the compactifiable solenoidal jet
$u=(0,0,y-x^2/2+x^2y/4)$ has complete viscous
$(D-\nu\Delta)K=+\nu$ at the origin. This exact original-state control
says nothing about the global integral sign or a singular trajectory.

With axisymmetric swirl, $\Gamma=ru_\theta$ retains maximum and physical
$L^1$ budgets, but $q=u_\theta/r$ has reaction $2(u_r/r)q$ and drives
$\eta$ by $\partial_z(q^2)$. The strain adapter needs the additional
swirl shear:
\[
 \lambda_2(S)_+\le |u_r/r|+\tfrac r2|\nabla_{r,z}q|.
\]
Source and shear doses remain unpaid. Neither calibration nor the new
identity/event accounting proves unconditional closure.

\section{Native magnitude dissipation and weighted direction return}
Let $\rho_\epsilon=(|\omega|^2+\epsilon^2)^{1/2}$ and
$g_\epsilon=\rho_\epsilon-\epsilon$. Direct original curl differentiation gives
\[
 (D-\nu\Delta)g_\epsilon
 =\frac{\omega\cdot S\omega}{\rho_\epsilon}-\nu G_\epsilon,
 \qquad G_\epsilon=
 \frac{|\nabla\omega|^2-|\nabla\rho_\epsilon|^2}{\rho_\epsilon}\ge0.
\]
Subtracting $\epsilon$ avoids the infinite constant-volume mass.
An exponential spatial weight proves $L^1$ legality from the existing
$L^2$ path before its radius tends to infinity. Since
$\int|S||\omega|\le Z/\sqrt2$ and $2\nu\int Z\le E_0$,
nonnegative Fatou as $\epsilon\downarrow0$ gives
\[
 \|\omega(t)\|_1+\nu\int_0^t\!\int_{\{\omega\ne0\}}
 |\omega||\nabla n|^2
 \le\|\omega_0\|_1+\frac{E_0}{2\sqrt2\nu}.
\]
This is an independently paid native direction budget, with possible
extra nonnegative zero-set dissipation discarded. It does not control
the critical stretching dose by itself.

Use the weighted near-field observable, distinct from the earlier
signed kinetic flux,
\[
 \Phi_R^{\rm dir}=\frac3{4\pi}\iint_{|x-y|<2R}
 \frac{|\omega(x)||\omega(x)\times\omega(y)|}{|x-y|^3}\,dy\,dx,
 \qquad c_R=\frac{\Phi_R^{\rm dir}}Z,
\]
with $c_R=0$ when $Z=0$. It is zero-safe and satisfies $c_R\le b_R$.
The fixed-radius remote payment gives
\[
 Z'+2\nu Q\le2(F_R+c_R)Z,
 \qquad F_R=\frac{39R^{-5/2}\sqrt{E_0}}{2\sqrt{5\pi}}.
\]
Finite full-horizon $\int c_R$ suffices for the original return.
It is enough instead to pay its mass on the union of the actual disjoint
growth intervals, since each has $\int_{J_k}(F_R+c_R)\ge\log2/2$.
This arbitrary-data supplier remains open.

There is an exact independent estimate. Double Cauchy--Schwarz,
the increment identity
\[
 \iint\frac{|\omega(x+r)-\omega(x)|^2}{|r|^4}\,dx\,dr
 =2\pi^2\|(-\Delta)^{1/4}\omega\|_2^2\le2\pi^2\sqrt{ZQ},
\]
and Sobolev interpolation give, for $Z>0$,
\[
 c_R\le3\sqrt\pi S^{3/2}\sqrt R\,\frac Q{\sqrt Z},
 \qquad S=4/\sqrt3.
\]
In the small-initial-product class, divide $Z'+\nu Q\le0$ through
$\sqrt{Z+\delta}$ and pass to a nonnegative limit. This gives
$\int Q/\sqrt Z\le2\sqrt{Z_0}/\nu$, and therefore
\[
 \int_0^t c_R\le\frac{6\sqrt\pi S^{3/2}\sqrt{RZ_0}}\nu
\]
without a positive lower bound on $Z$ or an assumed continuation.
No-swirl independently pays $b_R$ through the unsigned $\eta$ potential.
Neither class payment covers arbitrary data. Swapping the exact
Biot--Savart production also yields the absolutely convergent identity
\[
 \int\omega\cdot S\omega
 =\frac3{8\pi}\iint
 \frac{[(\omega_x-\omega_y)\cdot r]
       [r\cdot(\omega_y\times\omega_x)]}{|r|^5}\,dy\,dx.
\]
Its simultaneous increment and bending factors have not yet been
priced by the magnitude budget.

\section{Integrated bending sign and one-receiver control}
The complete integrated viscous coefficient is
\[
 \mathcal V=-2\sum_{j,k}\int|\partial_jJ_k|^2
 -4\sum_{j,k}\int J_k\cdot
       (\partial_j\omega\times\partial_j\partial_k\omega).
\]
It can be positive on actual solenoidal Schwartz states. To prove this,
take a compact smooth bump $v$ and its disjoint translate $v(\cdot-3)$.
For their one-dimensional heat evolutions define
$W_s=(0,-(H_sv)'',-(H_sv(\cdot-3))'')$.
Then $F(s)=\int|W_s\times W_s'|^2$ has $F(0)=0$ and $F(s)>0$
for $s>0$: a vanishing Wronskian would make a nonzero analytic
Schwartz function proportional to a fixed translate, contradicting
its finite $L^2$ mass. Hence some $\tau\in(0,1)$ has $F'(\tau)>0$.
The finite-scale original datum
\[
 u_L=\operatorname{curl}H_\tau^{(3)}
 \left[e^{-(y^2+z^2)/L^2}(0,v(x),v(x-3))\right]
\]
is solenoidal Schwartz. Exact separable heat and rescaled Schwartz
convergence give $L^{-2}\mathcal V(\operatorname{curl}u_L)
\to(\pi/4)F'(\tau)>0$. Thus a sufficiently large finite $L$ gives
the positive coefficient on the original state space.
Its full initial integrated derivative at datum $b u_L$ is
$b^4(\nu\mathcal V+b\mathcal N)$, where $\mathcal N$ is finite.
Small positive $b$ makes this positive and also satisfies the
small-initial-product threshold. Universal monotonicity of integrated
$K$ therefore fails even inside that globally regular calibration class.
This provides no later-time sign, singular trajectory or source payment.

For the separate one-receiver functional $\Pi_\beta$, its exact
variational gradient is
$G_\beta=H_{2\beta}B+(\nabla u)^TH_{2\beta}u
 -(u\cdot\nabla)H_{2\beta}u$ and $Q_\beta=-\langle G_\beta,\mathbb PB\rangle$.
At the earlier original Gaussian $U$ and $\beta=3/8$, exact full-pressure
Fourier/Gaussian moments give
\[
 \frac{Q_{3/8}(U)}{\pi^{3/2}}
 =-\frac{196957853}{335879864320}
 +\frac{2762863258\sqrt{13}}{39257040948125}<0.
\]
Both longitudinal pressure slots are nonzero; independent physical
raw-pairing checks agree, and $\beta=0$ cancels exactly.
The sign is certified by rational squaring, not floating quadrature.
It concerns one initial quartic coefficient and is not a trajectory law.

\section{Convex entropy and the fixed-parameter carrier}
The unregularized half-amplitude $h=\omega/\sqrt{|\omega|}$ need not
belong to $H^1$ at the initial face. For the admissible potential
$P=(0,0,x^3e^{-|X|^2})$, its vorticity has the form
$\omega_0=xe^{-|X|^2}W$, with $W_3=-6+O(x^2)$ near the plane $x=0$.
Thus $l=|x|a$ for smooth positive $a$, and
$\int |\nabla l|^2/l=\infty$ on a transverse patch. This is an initial
applicability defect, not a rejection of positive-time entropy laws.

Fix $\kappa>0$ and put $f(l)=l\log(1+l/\kappa)$. Direct differentiation gives
\[
 f'(l)=\log(1+l/\kappa)+\frac l{l+\kappa},\qquad
 f''(l)=\frac{l+2\kappa}{(l+\kappa)^2}.
\]
On each existing compact interval, use the zero-safe approximation
\[
 g_\epsilon=\sqrt{l^2+\epsilon^2}-\epsilon.
\]
It satisfies
\[
 (D-\nu\Delta)f(g_\epsilon)
 =\frac{f'(g_\epsilon)}{\rho_\epsilon}\omega\cdot S\omega
 -\nu\bigl[f'(g_\epsilon)G_\epsilon
                  +f''(g_\epsilon)|\nabla\rho_\epsilon|^2\bigr].
\]
Here the bracket lies between zero and $4|\nabla\omega|^2/\kappa$.
The mass is bounded by $l^2/\kappa$, and the source is dominated by
$f'(l)l|S|$. Exponential outer cutoffs disappear using compact-path bounds;
dominated convergence in $\epsilon$ then proves the exact integrated law
\[
 \int f(l(t))+\nu\int_0^t\!\int
 \bigl[f''(l)|\nabla l|^2+f'(l)l|\nabla n|^2\bigr]
 =\int f(l_0)+\int_0^t\!\int f'(l)l\alpha .
\]
Zero-set densities are assigned zero, justified by the almost-everywhere
vanishing of the smooth vorticity gradient on its zero set. Define the unpaid price
\[
 J_\kappa(t)=\int_0^t\!\int l|S|\log(1+l/\kappa).
\]
Native kinetic energy pays the remaining source term, yielding
\[
 \int f(l(t))+\nu\int_0^t\!\int
 \left[\frac{|\nabla l|^2}{l+\kappa}
          +l\log(1+l/\kappa)|\nabla n|^2\right]
 \le\int f(l_0)+\frac{E_0}{2\sqrt2\nu}+J_\kappa(t).
\]
This conditional dissipation law does not pay $J_\kappa$.

The fixed-parameter vector $H_\kappa=\omega/\sqrt{l+\kappa}$ has a valid
Sobolev chain rule at zeros. Its exact gradient split gives
\[
 |\nabla H_\kappa|^2
 =\frac{(l+2\kappa)^2}{4(l+\kappa)^3}|\nabla l|^2
     +\frac{l^2}{l+\kappa}|\nabla n|^2
 \le\frac{|\nabla l|^2}{l+\kappa}+l|\nabla n|^2.
\]
The native magnitude bound pays its $L^2$ mass; a finite $J_\kappa$ pays its
gradient time integral. Since $l\le |H_\kappa|^2+\sqrt\kappa|H_\kappa|$,
Sobolev and interpolation return an $L^1_tL^3_x$ vorticity bound on finite
horizons. This does not give the squared $L^3$ time dose or the remaining
weighted event charges.

The original Gaussian $U$ has $l^2=144e^{-10}$ and
$\omega\cdot S\omega=720e^{-15}$ at $(1,2,0)$.
Fixed-energy concentration $U_L=L^{3/2}U(LX)$ therefore gives
\[
 \int|f'(l_L)l_L\alpha_L|\ge cL^2\log L
\]
on a scaled open neighborhood,
whereas $Z_L=L^2Z_U$. This rejects a universal instantaneous absolute
entropy-source bound by $C(E_0,\nu,\kappa)Z$; it does not reject a signed
time-integrated supplier.

\section{Grouped pair diffusion and explicit event debts}
Let $V=\omega_x\times\omega_y$, $r=y-x$, $d=|r|$,
$g=\sqrt{l_x^2+\epsilon^2}-\epsilon$ and
$s=\sqrt{|V|^2+\delta^2}-\delta$. For a nonnegative fixed-radius cutoff
$\theta_R$, let $K_\lambda=\theta_R(d)(d^2+\lambda^2)^{-3/2}$ and
$I=(3/(4\pi))\iint K_\lambda gs$. Write
$\rho=\sqrt{l_x^2+\epsilon^2}$, $h=\sqrt{|V|^2+\delta^2}$ and
\[
 \mathcal R=-A_x^TV+\omega_x\times[(A_y-A_x)\omega_y].
\]
At positive regularizers, integration by parts groups the kernel Laplacian,
both coarea defects and the mixed product gradient exactly. The complete law is
\begin{align*}
 I'={}&\frac3{4\pi}\iint K_\lambda
 \left[s\frac{\omega_x\cdot S_x\omega_x}{\rho}
       +g\frac{V\cdot\mathcal R}{h}
       +\nu s\frac{\omega_x\cdot\Delta\omega_x}{\rho}
       +\nu g\frac{V\cdot\Delta_{xy}V}{h}\right]\,dx\,dy\\
 &+\frac3{4\pi}\iint gs(u_y-u_x)\cdot\widehat r\,K_\lambda'\,dx\,dy.
\end{align*}
The fixed-radius transition is included. The cancellations
\[
 V=\omega_x\times(\omega_y-\omega_x),\qquad
 \Delta_{xy}V=\Delta\omega_x\times(\omega_y-\omega_x)
             +\omega_x\times(\Delta\omega_y-\Delta\omega_x)
\]
give $O(d)$ vanishing for the grouped reaction and diffusion brackets.
The transport has two $O(d)$ factors against $K'\sim d^{-4}$.
All sources and the mass consequently have regularizer-uniform integrable
diagonal majorants on existing compact slabs. Sequential removal of the
regularizers returns an absolutely continuous integrated law for
$I_\theta=(3/(4\pi))\iint\theta_R l_x|V|/d^3\ge\Phi_R$, with zero
normalized factors assigned their regularized limits. Separate $d^{-5}$
kernel/coarea limits are not taken. The required majorants use higher
derivatives and are not endpoint-native prices.

For the admissible compactifiable local jet
$\omega=(0,bx,a_0+c(x^2+y^2)/4)$, the leading grouped viscous coefficient is
$2c\operatorname{sgn}(a_0)|a_0br_x|$. It has either sign. This is a local
initial control, not the sign of the full integrated functional.

The corrected increment bridge, with $M_p=\|\omega\|_p$ and
$D_3=\|\nabla|\omega|\|_3$, is
\[
 \Phi_R\le6R\left[M_5^{5/2}\sqrt{G_1}+M_3^2D_3\right].
\]
Its angular Holder exponents are $(5/4,5)$, not an interpolation from
$L^1\cap L^2$ to $L^3$. On the actual disjoint doubling intervals $J_k$,
put $z_k=2^kZ_0$, $A_k=z_k^{-2}\int_{J_k}M_5^5$,
$B_k=z_k^{-1}\int_{J_k}M_3^2D_3$ and $G_k=\int_{J_k}G_1$. Then
\[
 \sum_k\int_{J_k}c_R\le6R
 \left[\sqrt{\Bigl(\sum_k A_k\Bigr)\Bigl(\sum_kG_k\Bigr)}
                         +\sum_kB_k\right].
\]
Only $\sum G_k$ is independently paid. If both other sums were finite at
a hypothesized finite endpoint, the original weighted event return would
contradict its infinitely many required growth doses. Finally,
$P=\int\omega\cdot S\omega=Z'/2+\nu Q$ gives
$L_kd_k\ge z_k/\sqrt2$, where $L_k=\sup_{J_k,x}l$ and
$d_k=\int_{J_k}Z$. This lower amplitude debt cannot replace the upper
amplitude input needed to pay $A_k$.

\section{Native reciprocal weights and a capped tensor return}
The original enstrophy inequality and kinetic law give
\[
 Z'+\nu Q\le256\nu^{-3}Z^3,\qquad
 \int_0^tZ\le\frac{E_0}{2\nu}.
\]
For each fixed $\theta>0$, multiplication by $(\theta+Z)^{-2}$ and
the signed primitive $-1/(\theta+Z)$ give the native budget
\[
 \nu\int_0^t\frac{Q}{(\theta+Z)^2}
 \le\frac{128E_0}{\nu^4}+\frac1\theta-\frac1{\theta+Z_0}=:B_\theta.
\]
It also pays $\int Q^{1/3}$ by time Holder. At $\theta=Z_0>0$, the actual
event charge $\sum z_k^{-2}\int_{J_k}Q$ is at most $9B_\theta/\nu$.
The old direct $c_R$ majorant requires an additional $z_k^{3/2}$ factor;
that return remains unpaid.

There is a separate parameter-independent native source component. For the
uncentered cube maximal operator, a greedy triple-cube covering gives weak
constant $27$, and layer cake gives $\|Ml\|_2\le6\sqrt6\|l\|_2$.
Writing $G=\log(1+Ml/\kappa)$ and
$R_m=\log((\kappa+Ml)/(\kappa+l))$, the exact entropy source splits into
the maximal-log source, its ratio correction and a bounded remainder. Since
\[
 lR_m\le l\log(Ml/l)\le Ml/e,
 \qquad \left|\int l\alpha R_m\right|
                    \le\frac{6\sqrt6}{e\sqrt2}Z,
\]
the ratio is paid by native energy, including zeros. The fixed-threshold
high maximal-log source remains unpaid. No scalar BMO assertion pays its
direction product.

Projectors $P=n\otimes n$ represent the readout exactly:
\[
 |\omega_x\times\omega_y|
    =\frac{l_xl_y}{\sqrt2}|P_x-P_y|_F.
\]
An actual compact-curl flat zero-plane datum has orthogonal one-sided
projector limits, so a raw projective line increment cannot be assumed.
Instead use the globally $C^2$ tensor
$T=\omega\otimes\omega/(l+\kappa)$ and its trace
$b(l)=l^2/(l+\kappa)$. The full tensor law retains signed mixed and
projector diffusion terms; tracing gives the native positive law
\[
 \int b(l(t))+\nu\int_0^t\!\int
   \left[b''(l)|\nabla l|^2+b'(l)l|\nabla n|^2\right]
 \le\int b(l_0)+\frac{E_0}{2\sqrt2\nu}=:B_\kappa.
\]
For each fixed $M>0$, set $T_M=b(\min(l,M))P$, with value zero at
vorticity zeros. The weak cap derivative and the tensor zero adapter give
\[
 \|T_M(t)\|_2^2\le b(M)B_\kappa,\qquad
 \int_0^t\|\nabla T_M\|_2^2\le\frac{C_{\kappa,M}B_\kappa}{\nu}.
\]
\[
 C_{\kappa,M}=\max\left\{
 \frac{M^2(M+2\kappa)^2}{2\kappa^3},2M\right\}.
\]
Restrict $\Phi_{\rm low}$ to pairs with both $l_x,l_y\le M$. Exact
tensor difference and the global $H^1$ translation estimate imply
\[
 \Phi_{\rm low}\le3\sqrt2R(M+\kappa)\sqrt Z\|\nabla T_M\|_2.
\]
Thus on actual disjoint growth events at a hypothesized finite endpoint,
\[
 \sum_k\int_{J_k}\frac{\Phi_{\rm low}}Z
 \le3\sqrt2R(M+\kappa)
       \sqrt{\frac{T_*C_{\kappa,M}B_\kappa}{\nu Z_0}}.
\]
The zero path is handled undivided. This genuinely pays the capped low-pair
sector across arbitrary zeros, without assuming an entropy log price.
Pairs with at least one endpoint above the fixed cap remain unpaid;
there is no cap-to-infinity passage or all-data closure.

An actual growing cap can also be priced while retaining its time derivative.
For any positive nondecreasing locally absolutely continuous $m(t)$, use
the untruncated trace $b(l,m)=l^2/(l+m)$. Its parameter derivative is
$\partial_m b=-l^2/(l+m)^2$, giving the positive parameter work
\[
 H_m(t)+\nu\int_0^tD_m
       +\int_0^t m'(s)\int\frac{l^2}{(l+m)^2}
 \le H_m(0)+\frac{E_0}{2\sqrt2\nu}=:B_{\rm run}\le A_*.
\]
At each time the matched spatial cap $C_m=b(\min(l,m),m)P$ has
\[
 \int_0^t\frac{\|\nabla C_m\|_2^2}{m}\le\frac{9B_{\rm run}}{2\nu}.
\]
No time evolution of the capped trace is asserted. Choose, with fixed
$\delta>0$,
\[
 m(t)=\delta\left(\max_{s\le t}Z(s)\right)^{1/3}.
\]
On the actual growth intervals the running maximum is at most $2z_k$,
so $m^3/Z\le2\delta^3$. The same tensor translation and one event-union
Cauchy estimate give
\[
 \sum_k\int_{J_k}c_{\rm low,m(t)}
 \le18\sqrt2R\delta^{3/2}\sqrt{\frac{B_{\rm run}T_*}{\nu}}.
\]
This consumes one native budget. It leaves interactions with at least one
endpoint above the running cap as an explicit unpaid source. No repeated
per-level budget or omitted parameter flux is used.

The entropy increment also sharpens the conditional route to
$\Phi_R\le6\sqrt2R\sqrt{(\kappa M_4^4+M_5^5)(r_\kappa+g_\kappa)}$.
Actual tail interpolation compares its amplitude charge to
$\sum z_k^{-2}\int_{J_k}M_5^5$. Finite log source and finite fifth-moment
event sum suffice for the original event return, without a separate radial
supplier in that route. Both inputs remain open. Its signed entropy-gap
identity supplies only a one-sided residual price; high events require the
opposite-sign debt, which is not paid by that inequality.

\section{Native radial weights and amplitude tails}
Let $w$ be a fixed smooth nonnegative bounded integrable weight,
$W=\int_0^\infty w$, and define $\Psi''=w$, $\Psi(0)=\Psi'(0)=0$.
The radial composition is $C^2$ at zero with Hessian $w(0)I$;
$\Psi'/l\le\|w\|_\infty$ controls its angular Hessian.
Original chain calculus and compact-slab dominated convergence give
\[
 H_\Psi(t)+\nu\int_0^t D_w
 =H_\Psi(0)+\int_0^t\!\int\Psi'(l)l\alpha.
\]
Write $D_w=\int\mathcal D_w$, where
\[
 \mathcal D_w=w(l)|\nabla l|^2+\Psi'(l)l|\nabla n|^2.
\]
Since $\Psi'\le W$, native energy pays its absolute source, so
$H_\Psi(t)+\nu\int_0^tD_w\le B_\Psi\le WA_*$.
The separate supremum mass and total diffusion bounds each have price
$B_\Psi$; their sum with the supremum has price at most $2B_\Psi$.

For $w(l)=(l+\kappa)^{-1-\varepsilon}$, $\kappa>0$,
$0<\varepsilon\le1$, put
$h=\omega/(l+\kappa)^{(1+\varepsilon)/2}$.
Decreasing $w$ gives $\Psi'\ge lw$ and $\Psi\ge l^2w/2$.
The exact radial and tangential Jacobian coefficients then imply
\[
 |h|^2\le2\Psi(l),\qquad |\nabla h|^2\le\mathcal D_w,
 \qquad \int_0^t\|\nabla h\|_2^2\le B_\Psi/\nu.
\]
This is a native regularized carrier, without a logarithmic stretching
premise or an $\varepsilon\downarrow0$ limit.

Let $r_w=\int w(l)|\nabla l|^2$ and
\[
 V=M_5^5+2^\varepsilon
       (\kappa^{1+\varepsilon}M_4^4+M_{5+\varepsilon}^{5+\varepsilon}).
\]
The vector-increment Cauchy factors give
$\Phi_{d<\rho}\le3\sqrt2\rho\sqrt{V(G_1+r_w)}$.
For $\rho=\min(2R,\lambda Z^{3/2}/\sqrt V)$, native budgets pay
\[
 \int_{\cup J_k}c_{\rm near}
 \le3\sqrt2\lambda\sqrt{\frac{E_0}{2\nu}}
                         \sqrt{\frac{(1+W)A_*}{\nu}}.
\]
The remaining annulus has the unpaid upper source
\[
 3\frac{M_3^3}{Z}\log_+\!\left(\frac{2R\sqrt V}{\lambda Z^{3/2}}\right).
\]
The radius is selected in a fixed-time estimate and is never differentiated.
No weakest-criterion or reduced-difficulty claim follows.

The actual fifth-moment charge has a simpler independent return.
Native interpolation $Z^4\le M_1^3M_5^5$ and actual event duration give
\[
 a_k=z_k^{-2}\int_{J_k}M_5^5
       \ge\frac{3\nu^3}{8C_ZA_*^3},\qquad C_Z=256.
\]
Thus independently finite $\sum a_k$ already excludes a finite maximal
endpoint, without an entropy or direction input. Its arbitrary-data
upper supplier remains open; the comparison to $\int_{J_k}Z^2$ does not
control the gaps between events.

There is also a native amplitude-tail law. For $U_M=\int(l-M)_+$,
convex mollification away from zero and a direct change-of-level argument
give, for almost every $M>0$ and every existing time,
\[
 U_M(t)+\nu\int_0^t\!\int_{l=M}|\nabla l|\,d\mathcal H^2
       +\nu\int_0^t\!\int_{l>M}l|\nabla n|^2
 =U_M(0)+\int_0^t\!\int_{l>M}l\alpha.
\]
For every $M$ there is the separate inequality obtained by dropping
the radial flux and bounding the source on $l\ge M$.
On a hypothetical finite endpoint let
$C_M=U_M(0)+\int_0^{T_*}\int_{l\ge M}l|S|$.
Native space-time integrability of $l|S|$ gives $C_M\to0$, hence
\[
 \beta(M):=\sup_{t<T_*}\int_{l>2M}l\le2C_M\longrightarrow0.
\]
For $z_k\ge4MA_*$, restricted interpolation yields
$a_k\ge3\nu^3/(128C_Z\beta(M)^3)$; a hypothetical endpoint therefore
forces $a_k\to\infty$. This is an enhanced required lower charge.
Uniform first-moment amplitude tails give no upper critical price.

\section{Asymmetric caps and optimized distance estimates}
The matched running tensor also pays every low readout endpoint,
without restricting the other endpoint. For $m=\delta\bar Z^{1/3}$,
$\bar Z=\max_{s\le t}Z(s)$, write $C_m=b(\min(l,m),m)n\otimes n$,
$b(l,m)=l^2/(l+m)$. Contraction with $I-n_y\otimes n_y$ gives
$|C_x-C_y|\ge b(l_x,m)\sin\angle(n_x,n_y)$ when $l_x\le m$.
The original asymmetric pair carrier is therefore at most
$2m l_y|C_x-C_y|$. The single native cap-gradient budget yields
\[
 \int_{\cup J_k}\frac{\Phi_{l_x\le m}}Z
 \le36R\delta^{3/2}\sqrt{\frac{B_{\rm run}T_*}{\nu}}.
\]
For the reverse mixed sector, exact rank-one geometry gives
\[
 \int_{\cup J_k}\frac{\Phi_{l_x>m,l_y\le m}}Z
 \le18R\sqrt{\frac{B_{\rm run}H_4}{\nu}},\qquad
 H_4=\sum_kz_k^{-2}\int_{J_k}m\int_{l>m}l^4.
\]
Both $H_4$ and the high--high pair dose remain unpaid. The native lower
floor proportional to $z_k^{-2/3}$ is summable and supplies no upper price.

Put $D=G_1+r_w$, $\int_0^{T_*}D\le D_*=(1+W)A_*/\nu$.
The full-space positive carrier is legitimate on every compact original
interval and dominates $|P|$. For every $\rho>0$,
\[
 \Phi_{d<\rho}\le3\sqrt2\rho\sqrt{VD},\qquad
 \Phi_{d>\rho}\le\frac{\sqrt3}{2\sqrt\pi}Z^{3/2}\rho^{-3/2}.
\]
The second estimate is ordinary positive-kernel Cauchy and uses no signed
cancellation. Minimizing their sum divided by $Z$ gives
\[
 c_\infty\le C_*V^{3/10}D^{3/10}Z^{-2/5},\qquad
 C_*=\frac5{\sqrt2}(3/\pi)^{1/5},
\]
and consequently
\[
 \int_{\cup J_k}c_\infty\le C_*D_*^{3/10}W_Z^{7/10},\qquad
 W_Z=\int_{\cup J_k}V^{3/7}Z^{-4/7}.
\]
Independently finite $W_Z$ would contradict the $\log2/2$ required
per actual event. Its high-moment supplier is open.

The original velocity-energy tail gives a distinct optimized law.
For $C_{
m far}=39/(2\sqrt{5\pi})$,
\[
 (\log Z)'\le12\sqrt2\frac{\sqrt{VD}}Z\rho
                      +2C_{
m far}\sqrt{E_0}\rho^{-5/2}.
\]
Fixed-time minimization and Holder on the disjoint event union give
\[
 N\log2\le C_EE_0^{1/7}D_*^{5/14}W_E^{9/14},\qquad
 W_E=\int_{\cup J_k}V^{5/9}Z^{-10/9},
\]
where $C_E=(7/5)(12\sqrt2)^{5/7}(5C_{
m far})^{2/7}$.
No radius is differentiated; selectors are used where $D$ is finite,
almost everywhere. $D=0$ uses the limiting-radius bound. Zero paths
avoid logarithms; an interior zero implies future zero by forward
uniqueness. Neither $W_E$ nor $W_Z$ is natively supplied. Their density
ratio $V^{8/63}Z^{-34/63}$ gives no uniform comparison from the native
lower bound on $V$ alone. An elementary sibling criterion,
$\int_{\cup J_k}M_4^2/\sqrt Z<\infty$, follows directly from
$P\le M_4^2\sqrt{Z/2}$ and is also unpaid for arbitrary data.

\section{Moving levels and signed excess production}
For positive nondecreasing locally AC $m(t)$, one-sided convex
approximation gives the exact integrated positive-part law
\[
 U_m(t)+\nu\int_0^t\!\int_{l>m}l|\nabla n|^2+\nu\Lambda_m(t)
       +\int_0^tm'|\{l>m\}|
 =U_m(0)+\int_0^t\!\int_{l>m}l\alpha.
\]
Here $\Lambda_m$ is a continuous nondecreasing radial Stieltjes defect,
not an asserted classical moving-surface flux. Each separate nonnegative
term is bounded by $B_{\rm level}=U_m(0)+E_0/(2\sqrt2\nu)\le A_*$;
this does not replace the tensor-gradient price $B_{\rm run}$.

For $a=(l-m)_+$ let $\sigma_e=\int a^2\alpha$, using the original strain.
Since $0\le l^2-a^2\le2ml$, the original production splits as
\[
 P=\sigma_e+P_{\rm low},\qquad |P_{\rm low}|\le\sqrt2mZ.
\]
For $m=\theta Z/A_*$ the total low contribution to event log growth is
at most $4\theta$. For the monotone $m=\theta\bar Z/A_*$ it is at most
$8\theta$, since $\bar Z\le2Z$ on actual events. Hence
\[
 N\log2\le8\theta+2\sum_{k<N}\int_{J_k}\sigma_e/Z
\]
for the latter choice. An independently uniform upper bound on these
\emph{signed} prefixes would exclude finite maximality. It is open;
finite positive-part charge is a stronger sufficient input.
The squared-excess companion retains both signed sources and work:
\begin{align*}
 \int a(t)^2&+2\nu\int_0^t\!\int_{l>m}
 (|\nabla l|^2+al|\nabla n|^2)+2\int_0^tm'U_m\\
 &=\int a(0)^2+2\int_0^t\sigma_e+2\int_0^tm\int a\alpha.
\end{align*}
Instantaneous $m=\theta Z/A_*$ need not be monotone and its work is signed.
Both thresholds satisfy the paid subterm
$\int_0^{T_*}Z\int_{l>m}|\nabla n|^2\le A_*^2/(\theta\nu)$.
Neither this price nor the companion supplies the signed prefix.

\section{Moving radial gaps and projected curls}
Define $f(l,m)=l-m-m\log(l/m)$ for $l>m$, zero otherwise.
Its joint first derivatives vanish at the threshold. Direct moving-parameter
chain calculus, with $m\ge m_0>0$, gives
\begin{align*}
 H_f(t)&+\nu\int_0^t\!\int_{l>m}
 \left(\frac m{l^2}|\nabla l|^2+(l-m)|\nabla n|^2\right)
 +\int_0^tm'\int_{l>m}\log(l/m)\\
 &=H_f(0)+\int_0^t\!\int_{l>m}(l-m)\alpha.
\end{align*}
Dominated convergence gives the explicit radial density on compact intervals.
Each separate nonnegative term is bounded by
\[
 B_f=H_f(0)+E_0/(2\sqrt2\nu)\le A_*.
\]

The excess $e=(l-m)_+n$ satisfies
$\operatorname{div}e=\mathbf1_{l>m}(m/l)n\cdot\nabla l$.
Its actual solenoidal lift $v_e=\operatorname{curl}(-\Delta)^{-1}e$
has curl $Pe=e+G$, with $G=\nabla(-\Delta)^{-1}\operatorname{div}e$.
Explicit Fourier multipliers and Plancherel give
\[
 \int_0^{T_*}\frac{\|\nabla G\|_2^2}{m}
 =\int_0^{T_*}\frac{\|\operatorname{div}e\|_2^2}{m}
 \le B_f/\nu.
\]
For $c=\omega-e$ and its lift strain $S_c$, the Piola/Betchov identity
for the solenoidal $H^2$ lift yields the exact original excess source
\[
 \sigma_e=-4\int\det S_e-2\int e\cdot S_eG
                      -\int G\cdot S_eG+\int e\cdot S_ce.
\]
All longitudinal and cap-strain corrections remain. The weighted
projection price does not bound this complete signed cubic source.

Finally smooth radial weights supported in a fixed band $(M,2M)$ give
\[
 \int_0^{T_*}\!\int_{M<l<2M}|\nabla l|^2\le MC_M/\nu,
 \qquad \int_0^{T_*}|\{l>2M\}|^{1/3}\le\frac{16C_M}{3\nu M}.
\]
For $M_j=2^jM_0$, the sum of occupancy roots is bounded by
\[
 \sum_j\int_0^{T_*}|\{l>2M_j\}|^{1/3}
 \le\frac{32C_{M_0}}{3\nu M_0}.
\]
Unweighted $\sum_j C_{M_j}$ introduces the exact source count
\[
 1+\lfloor\log_2(l/M_0)\rfloor\quad\hbox{on }l\ge M_0.
\]
Native integrability of $l|S|$ does not supply it. Every critical
supplier and the original unconditional Clay(A) target remain open.

\section{Native vector forcing and a stretch-only return}
For a positive nondecreasing threshold, write $e=(l-m)_+n$ and
$T=\mathbf1_{l>m}(Se+m\alpha n)=DF_m(\omega)S\omega$.
The radial map is a contraction, so $|T|\le l|S|$. Its actual weak equation is
\begin{align*}
 (\partial_t+u\cdot\nabla-\nu\Delta)e
 &=T-m'\mathbf1_{l>m}n-\nu m\mathbf1_{l>m}|\nabla n|^2n\\
 &\quad-2\nu(m/l)\mathbf1_{l>m}\nabla l\cdot\nabla n-\nu n\mu.
\end{align*}
Here $\mu$ is the nonnegative spacetime radial level measure supported on
$l=m>0$, with mass at most $\Lambda_m$. Direction has unit length there;
local total variation of $n\mu$ is exactly $\mu$. No classical moving
coarea identity is presumed. Threshold, level and angular terms share
one $B_{\rm level}$ budget, while the gap pays the radial factor. Thus
\[
 |F|((0,T_*)\times\mathbb R^3)
 \le E_0/(2\sqrt2\nu)+B_{\rm level}
                 +2\sqrt{B_fB_{\rm level}}\le4A_*.
\]
For fixed $m(0)\ge M>0$, both budgets and the stretch source are bounded
by $C_M$, giving $\sup\|e\|_1\le C_M$ and $|F|\le4C_M\to0$.
This is a forcing-mass price, without a critical-norm return.

Pair the smooth ramp equation with its own excess. The mixed vector term
is orthogonal, and the level-energy term tends to zero before the measure
limit. Hence, for $Y=\|e\|_2^2$,
\[
 \tfrac12Y(t)+\nu\int_0^t\!\int_{l>m}
 (|\nabla l|^2+(l-m)l|\nabla n|^2)+\int_0^tm'U_m
 =\tfrac12Y(0)+\int_0^t\!\int e\cdot T.
\]
It follows, including at zeros, that
$\sqrt{Y(t)}-\sqrt{Y(s)}\le\int_s^t\|T\|_2$.
Choose fixed $m=M+\theta\bar Z/A_*$ with $0<\theta\le1/16$.
At sufficiently high actual upper-hit endpoints, $z_k\ge8MA_*$,
$Y(\tau_k)\ge3z_k/2$, whereas $Y(\sigma_k)\le z_k$. Consequently
\[
 \int_{J_k}\frac{\|T\|_2}{\sqrt Z}
 \ge\frac{\sqrt3}{2}-\frac1{\sqrt2}>0.
\]
An independently finite sum of these charges excludes finite maximality.
The stronger full-horizon $\int\|T\|_2<\infty$ also suffices. Both
arbitrary-data suppliers remain open. Globally smooth Gaussian forced-heat
controls have forcing mass and supremum $L^1$ tending to zero, fixed
$\int\|v\|_2^2$, and unbounded tested critical charges. They reject only a
uniform transfer from those budget sizes in that control class; they are
not original Navier--Stokes paths or actual-event counterexamples.

\section{Native radial carriers and longitudinal applicability}
For fixed $s>0$ and positive monotone $m$, set
\[
 f_s(l,m)=\int_m^l(l-r)m^sr^{-1-s}\,dr\quad(l>m),
 \qquad h_s=\int_m^l\sqrt{m^sr^{-1-s}}\,dr,
\]
zero below threshold. Direct moving-parameter calculus gives
$(f_s)_l=[1-(m/l)^s]/s$ and $(f_s)_m=[sf_s-l+m]/m\le0$.
The original law prices radial $m^sl^{-1-s}|\nabla l|^2$,
angular $l(f_s)_l|\nabla n|^2$ and favorable parameter work,
with each separate cost at most $C_M/s$ for $m(0)\ge M$.
Since $h_s^2\le2f_s$ and $h_s^2\le(l-m)(f_s)_l$, the vector carrier obeys
\[
 \sup\|h_sn\|_2^2\le2C_M/s,\qquad
 \nu\int\|\nabla(h_sn)\|_2^2\le C_M/s.
\]
At $s=1$, $h_1=\sqrt m\log(l/m)$. For $g=\log(l/m)_+$,
support H\"older and the original Sobolev law give
\[
 \int m^{5/3}[2C_{M/2}]^{-2/3}\|g\|_2^2
 \le\frac{16C_M}{3\nu},
\]
with the zero-tail case handled before division. The support refinement uses
$l\le2(l-M/2)$ on $l>m\ge M$. No critical stretch norm or free $s\downarrow0$
price follows.

The actual positive longitudinal correction is $G=\Pi c$, whose kernel is
\[
 G=c/3+\operatorname{PV}\frac{I-3\widehat x\otimes\widehat x}
                                      {4\pi|x|^3}*c.
\]
For $L=\operatorname{Lip}c$ a three-region split yields
\[
 \|G\|_\infty\le m\left[\frac73+\frac1{2\pi}
       +2\log_+\!\left(\frac{A_*^{1/3}L}{m^{4/3}}\right)\right].
\]
Original all-order persistence supplies compact-interval $L$; base
$H^\gamma$ alone does not supply $\nabla\omega\in L^\infty$.
Uniform spatial $C_0$ tails of the original compact-time path pass to the
regularized $G$. They justify the global weighted level-measure limit
even if unweighted defect mass escapes. The exact projected energy retains
transport $-\int(u\cdot\nabla e)\cdot G$, equivalently
$+\int e\cdot(u\cdot\nabla G)$, and level $-\nu\int n\cdot G\,d\mu$.
Terminal logarithmic and transport prices remain open.

\section{Positive radial continuation and its original dynamics}
The unprojected excess energy needs only
$r_+=\mathbf1_{l>m}l\alpha_+$ and $f_+=\|r_+\|_2$:
\[
 Y'/2\le\int(l-m)_+l\alpha_+\le\sqrt Y f_+.
\]
Negative radial production is favorable, and transverse stretch is
orthogonal in this pairing. The exact vector decomposition is
$T=\mathbf1_{l>m}[l\alpha n+(l-m)P_{n^\perp}Sn]$, so
$f_+\le\|T\|_2$. The actual-event floor and both conditional returns
above hold with $f_+$ in place of $\|T\|_2$. Also
$(\sigma_e)_+/Z\le f_+/\sqrt Z$; this supplies a one-way upper
signed-prefix comparison. Native $\int\|r_+\|_1\le C_M$ does not
pay its required spatial $L^2$ norm.

Let $B=P_{n^\perp}Sn$ and $H=\nabla^2p$. On $l>m>0$ the exact
pressure-complete original law is
\begin{align*}
 L\alpha={}&|B|^2-\alpha^2-n\cdot Hn
 +2\nu\sum_j\partial_jn\cdot(\alpha I-S)\partial_jn\\
 &-4\nu\sum_j\partial_jn\cdot(\partial_jS)n
 +(4\nu/l)\sum_j(\partial_jl)\partial_jn\cdot Sn.
\end{align*}
Native energy pays the unweighted transverse source
$\int|B|^2\le E_0/(4\nu)$. Its amplitude-weighted role, the full
directional Hessian and viscous couplings remain open.

The actual longitudinal contact gives projected level work
\[
 -(\nu/3)\int m\,d\mu-\nu\int n\cdot G_{\rm nl}\,d\mu.
\]
On the union $U$ of actual events with $z_k\ge z_K$,
\[
 \frac\nu3\int_U\frac mZ\,d\mu
 \le\left(\frac M{z_K}+\frac{2\theta}{A_*}\right)\frac{C_M}{3}.
\]
This native price also pays the constant terms of the nonlocal kernel
envelope. Its remaining absolute logarithmic child is
\[
 2\nu\int_U\frac mZ
       \log_+\!\left(\frac{A_*^{1/3}L}{m^{4/3}}\right)d\mu.
\]
It is open, together with the other projected sources and transport.
No normalized energy return follows by dividing an identity by $Z$
without its denominator derivatives. The original root remains open.

\section{Original acceleration and a threshold-vanishing radial source}
Let $\zeta=S\omega=A\omega$. The original product equation cancels
quadratic strain exactly:
\[
 L\zeta=-\nabla^2p\,\omega
       -2\nu\sum_j(\partial_jS)(\partial_j\omega).
\]
Here $\Omega\omega=0$ for the actual antisymmetric velocity gradient.
Pressure and derivative viscosity remain as sources. Native energy gives
$\int_0^T\|\zeta\|_1\le E_0/(2\sqrt2\nu)$. Moreover
$\zeta=\operatorname{div}(u\otimes\omega)$ and
\[
 \int_0^T\|u\otimes\omega\|_1
       \le E_0\sqrt{T/(2\nu)}.
\]
This prices ordinary pairings with bounded tests having bounded gradients.
Unbounded Lipschitz tests have only a flux-defined extension; nonlinear
threshold and direction test costs remain open.

For $r=l\alpha$ the exact original equation is
\begin{align*}
 Lr={}&l|B|^2-l\,n\cdot Hn+\nu r|\nabla n|^2
        -2\nu l\sum_j\partial_jn\cdot S\partial_jn\\
 &-4\nu l\sum_j\partial_jn\cdot(\partial_jS)n
        -2\nu\sum_j(\partial_jl)n\cdot(\partial_jS)n.
\end{align*}
Its $\alpha^2$ terms cancel, but it does not vanish at $l=m$.
Instead set $a=(l-m)_+$, $q_+=a\alpha_+$ and $g_a=\|q_+\|_2$.
The actual signed excess satisfies
\[
 (\sigma_e)_+\le\sqrt Y\,g_a\le\sqrt Z\,g_a,
 \qquad
 N\log2\le2\sqrt2MT_*+8\theta
                   +2\sum_{k<N}\int_{J_k}\frac{g_a}{\sqrt Z}.
\]
Finite actual-event charge would therefore give the same original global
return. This is a total prefix demand, without an individual-event floor.
The comparison $g_a\le f_+$ gives no strict separation of actual
trajectory finiteness classes; at large $l$, $a$ is asymptotic to $l$.
It also cannot replace $f_+$ in the unprojected excess-energy increment.

On the high set, write $R_+=\mathbf1_{\alpha>0}(R_0+R_1+R_2)$, where
\begin{align*}
 R_0={}&a|B|^2+m\alpha^2-a\,n\cdot Hn
                  +\nu(a-m)\alpha|\nabla n|^2,\\
 R_1={}&-2\nu a\sum_j\partial_jn\cdot S\partial_jn
         -4\nu a\sum_j\partial_jn\cdot(\partial_jS)n,\\
 R_2={}&-4\nu(m/l)\sum_j(\partial_jl)\partial_jn\cdot Sn
         -2\nu\sum_j(\partial_jl)n\cdot(\partial_jS)n.
\end{align*}
The exact positive field law retains two distinct nonnegative defects:
\[
 Lq_+=R_+-m'\mathbf1_{l>m}\alpha_+
             -\nu\alpha_+\mu-\nu a\lambda_\alpha.
\]
Both energy pairings vanish through separate pre-limit small factors,
because $q_+$ vanishes on the amplitude and $\alpha$ zero sets. This
removes neither defect from the field law. The compact energy identity is
\[
 \frac12\|q_+(t)\|_2^2+\nu\int_0^t\|\nabla q_+\|_2^2
 +\int_0^t m'\!\int a\alpha_+^2
 =\frac12\|q_+(0)\|_2^2+\int_0^t\!\int q_+R_+.
\]
Native energy separately pays the unweighted normalized reaction on
actual events $U_K$:
\[
 \int_{U_K}\frac mZ\int_{l>m}\alpha_+^2
 \le\left(\frac M{z_K}+\frac{2\theta}{A_*}\right)\frac{E_0}{4\nu}.
\]
It does not pay the energy source $ma\alpha_+^3$, pressure or derivative
viscosity. The critical supplier and original unconditional root remain open.

\section{Native time densities, fractional gradients and a paid far sector}
For a positive moving level $m\ge m_0$, the native mass bound gives
\[
 |U_m(t)-U_m(s)|\le\sqrt{2A_*/m_0}\|\omega(t)-\omega(s)\|_2
                         +(A_*/m_0)|m(t)-m(s)|.
\]
Actual compact-slab $C^1_tL^2_x$ vorticity regularity makes $U_m$ absolutely
continuous. The exact moving law therefore yields an absolutely continuous
cumulative radial defect, with $\lambda_m=\Lambda_m'\ge0$ and
$\nu\int\lambda_m\le C_M$ at the fixed offset. Uniform convergence of the
whole-space ramp cumulatives, followed by compact spacetime tests and spatial
exhaustion, proves the actual marginal domination
$\mu_{\rm time}\le\lambda_m\,dt$. A cumulative upper bound alone would not
justify this step.

Let $J_k=\|\nabla^ku\|_2^2$ use full ordered tensors, and
$c_p=\sqrt3+p-2$. Explicit ordered Leibniz contractions and elementary
integration by parts give, for $k=3,4$,
\[
 J_k'+2\nu J_{k+1}\le D_k\|\nabla u\|_\infty J_k,
 \quad D_3=2(4+3c_4),\quad D_4=2(5+10c_6c_3).
\]
A physical Fourier split, then multiplication by
$(\kappa+J_k)^{-2k/(2k-1)}$, uses the bounded primitive
$-(2k-1)(\kappa+J_k)^{-1/(2k-1)}$. Young absorption is paid by native
$\int Z$. It gives finite $\int J_3^{1/5}$ and $\int J_4^{1/7}$ through
any proposed finite maximal horizon. A second split yields
\[
 \int\|\nabla\omega\|_\infty^{1/3}
 \le(\sqrt2G_{2,4})^{1/3}(\int Z)^{1/36}
                                (\int J_4^{1/7})^{35/36},
 \quad G_{2,4}=\frac1{\sqrt2\pi\sqrt5}+\frac1{\sqrt2\pi}.
\]
These fractional time powers give no critical continuation by themselves.

Set $h=\|\nabla\omega\|_\infty^{1/3}$,
$K=A_*^{1/9}M^{-4/9}$ and $c_K=M/z_K+2\theta/A_*$. The logarithmic level
component has the conditional estimate
\[
 C_{\log}\le6\nu c_K
 \left[H_m+\lambda_0\left(T_*+K\int_0^{T_*}h\right)\right],
 \quad H_m=\int_0^{T_*}\lambda_m\log_+(\lambda_m/\lambda_0).
\]
The reference $\lambda_0>0$ has density units. The factor three from the
cube inside the logarithm is retained. Native $L^1$ mass does not pay the
open entropy $H_m$; the other projected sources remain.

The original fixed-$R$ smooth shell also pays both radial far contributions.
Its curl estimate $4+6\sqrt2<13$ yields
$F_R=39R^{-5/2}\sqrt{E_0}/(2\sqrt{5\pi})$ and
\[
 \frac{g_a}{\sqrt Z}\le\frac{g_{a,R}^{\rm near}}{\sqrt Z}+F_R,
 \qquad
 \frac{f_+}{\sqrt Z}\le\frac{f_{+,R}^{\rm near}}{\sqrt Z}+F_R.
\]
The disjoint event union costs at most $F_RT_*$ once. The near radial
supplier remains open; $c_R\le b_R$ does not bound its spatial $L^2$ norm
by the weighted scalar $c_R$. The original unconditional root is open.

\section{Actual level diffusion, band budget and production controls}
Positive-level uniform $C_0$ tails confine every level ramp to one finite
ball on each existing compact slab. This identifies the full temporal
marginal $\mu_m(dt,\mathbb R^3)=\lambda_m(t)dt$. Comparing the original
positive-part time chain with its exact moving law gives
\[
 \lambda_m+G_m=-\int_{l>m}n\cdot\Delta\omega\ge0,
 \qquad
 \int m^{1/5}(\lambda_m+G_m)^{2/5}
       \le C_{\rm support}^{1/5}\int J_3^{1/5}.
\]
Here $C_{\rm support}=A_*$, or the tail refinement $2C_{M/2}$.
This native fractional price does not give time-density entropy.
For fixed $0<\theta_0<\theta_1\le1/16$ and
$m_\theta=M+\theta\bar Z/A_*$, the actual joint ramp limit is
\[
 \int_{\theta_0}^{\theta_1}\mu_{m_\theta}\,d\theta
  =\frac{A_*}{\bar Z}\mathbf1_{m_{\theta_0}<l<m_{\theta_1}}
                       |\nabla l|^2\,dxdt.
\]
Tonelli and each native level law pay its total mass by
$(\theta_1-\theta_0)C_M/\nu$. The restricted band is essential;
its logarithmic weight and the all-space $Q/Z$ price remain open.

For $\mathcal L=\partial_t+u\cdot\nabla-\nu\Delta$,
$d=|\nabla l|^2$ and $b=\mathcal Ll$, the original gradient law is
\[
 \mathcal Ld=2\nabla l\cdot\nabla b
       -2\nabla l\cdot S\nabla l-2\nu|\nabla^2l|^2.
\]
The exact law for $k_\delta(l-m)d$ retains the ramp derivatives,
moving work and mixed diffusion. Parameter-band differentiation also
retains $-\bar Z'B_\delta/\bar Z$. Neither native mass nor ramp scaling
alone pays the complete time-entropy production.

Tracefree algebra bounds the local positive-radial block by
$a^2\alpha(|S|^2-l^2/6)+a(m-a)\alpha^3$ on $\alpha>0$.
The positive cubic coefficient lies in $m<l<2m$; compact solenoidal
Schwartz jets realize both signs of this local block. The full nonlocal
pressure and nonlinear test derivatives remain.
A separate actual Gaussian solenoidal datum has a persistent regular
positive level. Direct coordinates produce a nonzero singular graph
measure as the ramp-density limit. Uniform spacetime entropy of the
spatial ramp densities would imply uniform integrability and contradict
that graph mass. This rejects only that auxiliary estimate; temporal
entropy of $\lambda_m$ and grouped cancellations remain viable.

\section{Projection and native corrections}
The actual threshold fields have $e\in AC_{\mathrm{loc}}(L^2)$;
the threshold derivative is locally $L^1$ in time. With
$w=Pe=e+G$ and $G=-\Pi e$, orthogonality gives
$\|e\|_2^2=\|w\|_2^2+\|G\|_2^2$. For the original vector forcing
$F=(\partial_t+u\cdot\nabla-\nu\Delta)e$, the complete projected
energy correction groups into longitudinal storage and dissipation:
\[
 \tfrac12(\|w\|_2^2)' +\nu\|\nabla w\|_2^2
 =\langle F,e\rangle-\tfrac12(\|G\|_2^2)'
                         -\nu\|\nabla G\|_2^2.
\]
This bypasses termwise trace envelopes in this energy identity; it does
not supply its remaining production or retire other consumers.
A fixed amplitude band $m<l<Km$ has native normalized time price through
$B_m=MT_*+\theta E_0/(\nu A_*)$. The complementary above-band source
remains open.

A native cube $L^2$ atom calculation bounds the longitudinal projection
by $C_{\mathrm{at}}=\sqrt{32\pi\sqrt3}+256/27$ in $L^1$.
The inspected singular-integral proof supplies the original projection's
$L^4$ law. Its quadratic correction has a native actual-event price.
For the mixed correction, a direct non-tangential grand-maximal estimate
uses $E=w$, $\pi=u_i$ and $\nabla\pi\in L^2$:
\[
 \|M_6(E\cdot\nabla\pi)\|_1
       \le C_{\mathrm{gm}}\|E\|_2\|\nabla\pi\|_2.
\]
Subtracting nested ball means, critical Sobolev--Poincare and annular
kernel decay give the summable shell factor $2^{-2j}$. Replaying the
inspected constructive atomic proof directly with this grand-maximal
norm yields the exact cube-atom input, without assuming a recalled
Poisson-norm equivalence. Consequently
\[
 \|\Pi((w\cdot\nabla)u)\|_1
       \le C_H\|w\|_2\|\nabla u\|_2,
 \qquad \left|\int G\cdot Sw\right|\le C_HmZ.
\]
The cap $c$ is bounded by $m$, and the actual events pay $\int_Um\le B_m$.
Writing $\sigma_e=\int e\cdot Se$ and $\sigma_w=\int w\cdot Sw$, the exact
mixed and quadratic expansion therefore bounds their normalized prefix
difference by a finite native constant. Their signed suppliers remain open.

\section{Actual self--cap decomposition and a scoped kernel control}
Native vorticity $L^1$ and compact-interval $H^1$ justify the low-frequency
Biot--Savart velocities $v=BS(w)$ and $v_c=BS(c)$ in $H^2$.
The original $u=v+v_c$, and signed Piola integration followed by $H^2$
approximation gives
\[
 \sigma_w=-4\int\det S_w+R_{\mathrm{cap}},
 \qquad R_{\mathrm{cap}}=\int w\cdot S_cw,
 \quad S_w=\operatorname{sym}\nabla v.
\]
This supplies no sign. For $T=w\otimes w$, the actual quadratic transpose is
\[
 R_{\mathrm{cap}}=\int c\cdot\Psi_w,
 \qquad\Psi_w=-\operatorname{curl}(-\Delta)^{-1}\operatorname{div}T.
\]
Symmetry cancels the Hessian contact term; its off-diagonal kernel is
$-3\hat r\times(T\hat r)/(4\pi|r|^3)$. A fixed-radius split of this
already differentiated kernel pays
\[
 \int_U\frac{|R_{\mathrm{cap,far}}|}{Z}
                      \le\frac{3A_*T_*}{4\pi R^3}.
\]
The combined signed self-determinant and near-cap production remains
an independent unpaid supplier.

For the admissible Gaussian velocity
$u_0=(\partial_y e^{-|x|^2},-\partial_x e^{-|x|^2},0)$,
vorticity covariance is $(\pi/2)^{3/2}\operatorname{diag}(1,1,8)$.
Every sufficiently small positive excess threshold retains nonscalar
projected covariance $C_m=\int w_m\otimes w_m$. Its actual transpose has
\[
 \Psi_{w_m}(r n)=-\frac{3}{4\pi r^3}n\times(C_m n)+O(r^{-4}).
\]
A nonzero angular cone forces $\|\Psi_{w_m}\|_1=\infty$.
This rejects only the generic $m\|\Psi_w\|_1$ envelope. The coupled
$L^2$ pairing and native remote price remain valid; no singular trajectory
or failure of an original continuation supplier follows.

\section{A sublinear cap loss and the remaining self-source}
A quantitative native Calderon--Zygmund argument for each strain component
and its adjoint gives a weak $(1,1)$ constant $W=1024$.
Distribution-function interpolation and duality yield
$\|S_c\|_p\le K_0p\|c\|_p$ for $p\ge4$.
For $h=e$ or $w$, clipping, projection, Sobolev and Young give
\[
 \frac{|\int h\cdot S_ch|}{Z}
 \le\frac{\nu Q}{4Z}
   +C p m\left(\frac{64A_*^2m}{\nu^3}\right)^{1/(2p-3)}.
\]
Choose $p_j=j+4$ on the actual event with $z_k=\zeta2^j$.
Its remainder is at most $C_{\mathrm{fixed}}(j+4)m$.
The native sequence $b_j=\int_{J_{k_0+j}}m$ is summable, hence
$\sum_{j<n}(j+4)b_j=o(n)$. This does not assert a finite full weighted sum.
The original enstrophy identity, including its native low-production and
projection errors, therefore gives
\[
 \frac n2\log2+\frac{3\nu}{4}\int_{U_n}\frac QZ
 \le\sum_{j<n}\int_{J_{k_0+j}}\frac{-4\int\det S_w}{Z}+o(n)+C.
\]
An independent upper self-source slope below $\log2/2$ would suffice.
That producer is open.

The signed excess self-core has the exact absolutely convergent identity
\[
 \int e\cdot S_we=\frac3{8\pi}\iint
 \frac{[e_x-e_y,\hat r,e_y]((e_x-e_y)\cdot\hat r)}{|r|^3}\,dx\,dy,
 \quad r=x-y.
\]
Both endpoints are high; signed and absolute pair criteria remain distinct.
Native excess $L^1$ pays its fixed-radius remote sector. Applying the
reviewed Hardy law anew with potential $v_i=BS(w)_i$ pays the self-strain
mixed correction; its quadratic correction has the native $L^4$ price.
Thus the excess core and projected Betchov self-source have a bounded
normalized prefix difference, rather than being exactly equal.

The pressure source $f=\sum_{ij}\partial_i u_j\partial_j u_i$ is a sum
of actual div--curl products. The direct grand-maximal and atomic laws,
with the scalar Hessian multiplier equal to the negative projection
component, give
\[
 \|\nabla^2p\|_{L^1}\le C_pZ,
 \qquad\int_0^{T_*}\|\nabla^2p\|_1\le C_pE_0/(2\nu).
\]
This pays bare or fixed bounded tests. Amplitude-weighted pressure
production remains open. The same native singular-integral argument gives
$\|S\|_3\le C_3\|\omega\|_3$ and the actual-event lower demand
$\int_J\|\omega\|_3^3/Z\ge\log2/(2C_3)$; its supplier is also open.

\section{Native matrix direction and complete joint geometry}
For $X=\nabla u$ or $S$, the original matrix equation and the pressure
Hessian law bound its magnitude source in $L^1$ by $(1+C_p)Z$.
A zero-safe Kato identity and whole-space cutoff argument prove
\[
 \|X(t)\|_1+\nu\int_0^t\!\int_{X\ne0}
 |X|\,|\nabla(X/|X|)|^2
 \le\|X(0)\|_1+(1+C_p)E_0/(2\nu).
\]
These one-amplitude matrix-direction budgets do not pay critical strain
or nonlinear pressure weights.
The complete integrated pressure contraction is
\[
 -\int\omega\cdot\nabla^2p\,\omega
 =\int l\,\partial_n l\,\partial_n p+\int l^2\kappa\cdot\nabla p,
 \qquad\kappa=(n\cdot\nabla)n.
\]
The excess identity retains the extra threshold coefficient $ma/l$;
the nonlinear positive-radial test retains axial amplitude and strain
derivatives and curvature costs.

For $s=|S|$, $q_0=\sqrt{2/3}$ and $\alpha=n\cdot Sn$, let
$K_q=l(qs-\alpha)$, with $q=1$ or $q_0$.
Tracefree algebra makes $K_1$ comparable to $ls$, so it is not a
vanishing depletion measure. The sharper $K_{q_0}$ can vanish on actual
local curl jets. The smooth nonnegative carrier
$q(\rho\sigma-\epsilon\eta)-(\omega\cdot S\omega)/\rho$
has an exact original evolution retaining all pressure and mixed diffusion.
For $D=nn^T-I/3-q_0S/s$,
\[
 |D|^2=2q_0K_{q_0}/(ls),\qquad
 l\nabla^2p:(nn^T-q_0S/s)=l\nabla^2p:D+(l/3)\Delta p.
\]
The trace term remains. Neither the native carrier mass nor the separate
direction budgets supplies this complete weighted production.

\section{Actual self absorption and adaptive pair frontier}
The native order-zero inverse-curl gradient estimate at exponent $3/2$,
combined with elementary $W^{1,1}$ Sobolev applied to $|v|^2$, gives
$\|v\|_3\le C_v\|e\|_{3/2}\le C_v(UY)^{1/3}$, where
$v=BS(e)$, $U=\|e\|_1$ and $Y=\|e\|_2^2$.
Solenoidal integration by parts then gives
\[
 |C_w|\le C_S\|v\|_3Q,\qquad
 \frac n2\log2\le\int_{U_n}
 [C_S\|v\|_3-3\nu/4]_+\frac QZ+o(n)+C.
\]
This absorbs good times. The weighted bad-time remainder still needs an
independent upper price. The exact critical product law retains the
sources $Y\int_{l>m}l\alpha+2U\int_{l>m}al\alpha$, moving-threshold
work, radial level density and all diffusion. In particular the scalar
$U$ angular coefficient is $l|\nabla n|^2$; replacing it by $m|\nabla n|^2$
would omit a favorable term from an equality.

The instantaneous radius $\rho_j^3=A_*/((j+4)m)$ pays the remote
self-source by $(3/(4\pi))\sum_{j<n}(j+4)\int_{J_j}m=o(n)$.
Its complementary near region includes $0<|x-y|<2\rho_j$ and the cutoff
transition; there is no established lower core radius. A source split
introduces no radius derivative, whereas an evolving pair functional would.
The original matrix Kato laws also give uniform amplitude tails of
$\nabla u$ and $S$, without a quantitative critical tail rate.

For fixed smooth compact radial $\phi$, the exact pair mass
$\iint\phi|e_x\times e_y|^2$ retains both complete vector force pairings,
both gradient dissipations, drift from $u_x-u_y$ and $2\nu\Delta\phi$.
The force contains smooth chain diffusion; its presence prevents inferring
a sign for the complete diffusion from the explicit gradient squares alone.
With $D=nn^T-I/3-q_0S/s$, the native direction laws give
\[
 \nu\int\min(l,s)|\nabla D|^2\le4A_*+2q_0^2A_S.
\]
Neither this one-amplitude budget nor the smooth pair identity supplies a
cubic source or a singular-kernel return.

\section{Complete pair diffusion and an unclipped convex law}
The explicit negative gradient squares in the clipped quartic pair law
do not fix the sign of its full diffusion. For one endpoint with magnitude
$l$, excess $a=l-m$, direction $n$, and the other excess $bN$, set
$c=n\cdot N$, $d=\partial l$, $h=\partial n$. The complete local block is
\[
 J=d^2(1-c^2)+[a^2+am(1-c^2)]|h|^2-a^2(h\cdot N)^2
       -2a(1+m/l)dc(h\cdot N).
\]
Its contribution is $-2\nu b^2J$. The admissible local jet
$l=2$, $m=a=b=1$, $n=e_3$, $N=(3/5,0,4/5)$, with only
$\partial_2\omega=2e_1+2e_3$, has $J=-11/25$.
A two-stage polynomial potential realizes it in compact solenoidal
original velocity data. This tests a local block, not a global pair sign.
The absolute complete diffusion admits a bound by
$12\nu\|\phi\|_\infty YH_m$, where
$H_m=\int_{l>m}(|\nabla l|^2+al|\nabla n|^2)\le Q$;
neither $\int Q$ nor $\int H_m$ is a native endpoint budget.

For original $V=\omega_x\times\omega_y$, the convex regularized norm
$g_\epsilon=(|V|^2+\epsilon^2)^{1/2}-\epsilon$ has a complete fixed
smooth-kernel law. Its mass, both stretching terms, relative drift and
kernel Laplacian have native prices from $A_*$, kinetic energy and
$\int Z\le E_0/(2\nu)$. The zero-set defect remains as part of the
regularized identity; Fatou's lemma alone does not identify it.
Actual vorticity is continuous in $L^1$ through the initial face by the
original curl mild equation and integrable source bounds.

For the actual near core, the angular and radial moments
\[
 \mathrm{Ang}=\iint\psi_\rho\frac{a_xa_y^2|n_x-n_y|^2}{r^3},\qquad
 \mathrm{Rad}=\iint\psi_\rho\frac{a_x|a_x-a_y|^2}{r^3}
\]
are restricted to both high endpoints, and obey
\[
 |C_{e,\mathrm{near}}|\le\frac3{8\pi}
 [\sqrt{\mathrm{Ang}\,\mathrm{Rad}}+\mathrm{Ang}].
\]
Their fallback $|C_{e,\mathrm{near}}|\le6L_e\rho^2Q$ retains the
unpaid compact-time amplitude $L_e=\|e\|_\infty$.

\section{Singular pair laws and exact diagonal compensation}
Let $\phi=r^{-3}\eta(r/R)$, with a full smooth inner hole before
limits and a nonincreasing outer profile. For original $V$, compact-time
global Sobolev translations give center $L^2$ bounds of order $r$ for
$V$, its full stretching source $F_2$, and $\Delta_6V$.
For the quartic pair, the complete viscosity must be combined as
\[
 -2\nu\iint\phi|\nabla_6V|^2
 +\nu\iint(\Delta_6\phi)|V|^2
 =2\nu\iint\phi V\cdot\Delta_6V.
\]
Here $\Delta_6\phi=2\Delta_{\rm rel}\phi$.
Writing $K=\int\sum_j|\omega\times\partial_j\omega|^2$, the two
separated leading logarithms have coefficients $-16\pi\nu K$ and
$+16\pi\nu K$. The identity
$\int r^4\Delta_{\rm rel}\phi=6\int r^2\phi$ includes every hole
and outer transition. Individual terms are not finite singular budgets.
For the norm carrier, global center $L^1$ translation bounds give the
integrable radial majorant $dr$ for its combined physical generator.
These are legal local identities, not terminal source estimates.

The actual shrinking radius $\rho_j^3=A_*/((j+4)m)$ gives favorable
radius work for a nonnegative pair mass because
$\rho'\le0$ and $\partial_\rho\phi\ge0$.
Downward index resets give nonpositive jumps. Gap contributions remain,
and normalization by $Z$ retains $-BZ'/Z^2$ with the original signed
$Z'=2P-2\nu Q$. Favorable radius work supplies no critical event charge.

\section{Cubic degree and the missing readout}
For $q=3/2$, $g_\epsilon=(|V|^2+\epsilon^2)^{q/2}-\epsilon^q$
is convex and has a complete singular local pair law. Global
$L^6$--$L^2$ translations give $L^{3/2}$ increments of order $r$;
the radial majorant is $r^{q-1}dr$. For fixed smooth kernels,
\[
 B_q\le\|\phi\|_\infty A_*Z,\qquad
 |\mathrm{reaction}|\le\frac3{\sqrt2}\|\phi\|_\infty
                  \sqrt{A_*}\,Z\,\|\omega\|_3^{3/2}.
\]
The drift contains $2\|\nabla\phi\|_\infty
\sqrt{A_*E_0Z}\,\|\omega\|_3^{3/2}$. These nonlinear source
weights are unpaid; bounded-kernel prices do not extend to $r^{-3}$.

Matching cubic degree does not give the core readout. With
$X=e_3$, $Y=\sin\theta e_2-\cos\theta e_3$ and
$\widehat r=(e_1+e_3)/\sqrt2$, the ratio of the actual core numerator
to $|X\times Y|^{3/2}$ is
$(1+\cos\theta)/(2\sqrt{\sin\theta})$, which diverges.
This refutes only a uniform pointwise adapter, not integrated cancellation.
The valid envelope $F=|X\times Y||X-Y|$ has an indefinite separate
Hessian and a nonzero high-side threshold derivative
$b^2|n\times N|$. Its full source, diffusion and level measures remain
required. The separately convex envelope
$|X|^2|Y|+|X||Y|^2$ has a nonzero diagonal and therefore an infinite
absolute singular pair integral. The next construction seeks a convex
diagonal-vanishing carrier with a complete original-source return.

\section{A diagonal convex cubic and its exact law}
For endpoint values $X,Y$, let $a=|X|$, $b=|Y|$, $d=X-Y$ and
\[
 G(X,Y)=(a+b)|d|^2+\frac23|d|^3.
\]
This finite symmetric cubic is separately convex, vanishes on the
diagonal, and obeys
\[
 \sup_{|q|=1}|(q\cdot(X\times Y))(q\cdot d)|
 =\frac12|X\times Y||d|\le\frac14G(X,Y).
\]
The constant $1/4$ is approached by equal amplitudes and a small angle.
Radial excess clipping decreases both amplitudes and the increment,
so $G(e_x,e_y)\le G(\omega_x,\omega_y)$. Its original singular
pair moment satisfies
\[
 B_G=\iint\phi_\rho G(\omega_x,\omega_y),\qquad
 |C_{e,\mathrm{near}}|\le\frac{3B_G}{32\pi}.
\]
A strict event-prefix slope below $(16\pi/3)\log2$ for $B_G/Z$
is sufficient; its supplier is open.

The full original common-time generator retains both stretching forces,
relative drift and combined physical viscosity. With regularized norms
$\rho_X=(|X|^2+\epsilon^2)^{1/2}$ and
$\rho_d=(|d|^2+\epsilon^2)^{1/2}$, its diagonal angular mass contains
\[
 \frac{4\pi}{3}(2K_\epsilon+\epsilon Q)r^2,
 \qquad K_\epsilon=\int\rho_\omega|\nabla\omega|^2.
\]
The complete hole logarithms cancel with opposing coefficients
$\nu(32\pi K_\epsilon+16\pi\epsilon Q)$ at fixed $\epsilon$,
before sending $\epsilon$ to zero. The regularization's
$\epsilon|d|^2$ term must not be omitted. Higher weighted gradients
remain unpaid. Also $G(0,Y)=5|Y|^3/3$; those pairs add real mass even
though the signed cross core vanishes there.

\section{Original increment and viscosity allocation}
For original vorticity define
\[
 B_\rho=\iint\phi_\rho|\omega_x\times\omega_y|^2,\qquad
 D_\rho=\iint\phi_\rho|\omega_x-\omega_y|^2,
 \qquad\phi_\rho=r^{-3}\eta(r/\rho).
\]
The exact full-sphere translation inequality is
\[
 D_\rho\le\frac{8\pi}{3}\rho^2Q.
\]
The quadratic increment law retains both
stretching forces, relative transport and the complete physical viscosity;
its separate leading hole logarithms $\pm16\pi\nu Q$ cancel.
With physical Fourier convention $2\pi$, its multiplier is
\[
 \Lambda_\rho(\xi)=8\pi\int_0^{2\rho}
 \eta(r/\rho)[1-\operatorname{sinc}(2\pi|\xi|r)]\,\frac{dr}{r}.
\]
The sharp angular factor and Cauchy give
\[
 |C_{e,\mathrm{near}}|\le\frac3{16\pi}\sqrt{B_\rho D_\rho}
 \le\frac{3\nu}{4}Q+\frac{\rho^2 B_\rho}{32\pi\nu}.
\]
Consequently a strict normalized event slope below $16\pi\nu\log2$
for $\rho^2B_\rho/Z$ is sufficient. With viscosity allocation
$\lambda\nu Q$, $0<\lambda\le3/4$, the threshold is
$(64\pi/3)\lambda\nu\log2$; the factor $1/3$ is required.

For actual event intervals $J_k$, the anticipating radius
$R_k^3=3A_*|J_k|/(4\pi\delta)$ prices the far contribution by
$\delta$ per event. Taking its minimum with the canonical radius
gives $\delta n+o(n)$, and the sufficient pair threshold becomes
$16\pi\nu\log2-32\pi\nu\delta$, where $0<\delta<\log2/2$.
This instantaneous split requires no radius derivative. Evolving the
carrier still retains upward interevent jumps, gaps and signed $Z'$.
Neither the logarithmic multiplier nor these event prices supplies the
critical quartic pair dose.

\section{Native signed-component geometry and its return gap}
Let $L=\max_i\|\omega_i\|_\infty>0$, $V_3=4\pi/3$,
$q=(1+\lambda)\delta-1>0$, and choose
$(1+\eta)^3=1+q/2$. A radial linear $H^1$ cutoff gives
\[
 C_\eta^2=\frac23 V_3\frac{(1+\eta)^3-1}{\eta^2},\qquad
 r\ge K_0E_0^{1/5}L^{-2/5},\quad
 K_0=\left(\frac{2C_\eta}{qV_3}\right)^{2/5}.
\]
At every such radius and every center, each of six signed-component
sets $\{\pm\omega_i>\lambda L\}$ occupies at most fraction $\delta$.
The proof compares the strict lower pairing $qV_3Lr^3/2$ with the
curl-test upper bound $C_\eta\sqrt{rE_0}$. It works at all radii,
including the equality threshold, and uses the actual datum energy.
This is component geometry, not full-vector sparseness or continuation.
Its primary motivation is Bradshaw--Farhat--Gruji\'c,
\url{https://arxiv.org/abs/1704.05546v4}; the direct proof imports
none of that paper's analyticity theorems.

Against $r_c=\kappa\sqrt{\nu/L}$, the native radius ratio is
$(K_0/\kappa)E_0^{1/5}\nu^{-1/2}L^{1/10}$. A sufficient localized
energy rate to reach $r_c$ is
\[
 \int_{B(x_0,(1+\eta)r_c)}|u|^2
 \le\left(\frac{qV_3}{2C_\eta}\right)^2
 \kappa^5\nu^{5/2}L^{-1/2}.
\]
That rate and a subsequent geometric or analytic return are open.
Separately, smooth mollified truncated logarithms
$[\log R-\max(0,\log|x|)]_+$ have bounded BMO seminorm but a
weighted absolute integral against $(1+|x|)^{-4}$ growing like
$\log R$. This rejects only a generic seminorm-only scalar estimate;
actual local means, Riesz cancellation and the PDE analytic branch
remain to be paid. It is not a rejection of the primary's NS theorem.

\section{Newton core and a scoped potential control}
For the actual excess $e$, set $c=e_x\times e_y$, $d=e_x-e_y$,
$\psi_\rho=\eta(r/\rho)$ and $\Phi=1/r$. The signed near core is
\[
 C_{e,\mathrm{near}}=-\frac1{8\pi}\iint
 \psi_\rho c_i(\partial_{ij}\Phi)d_j.
\]
The Hessian contact vanishes since $c\cdot d=0$; the weak inner-hole
boundary is of order $r^2$. Integration by parts retains actual
$\operatorname{curl}e$, its directional gradient and the full outer
transition. With $a=|e|$, define
\[
 \Psi_\rho[a](x)=\int\frac{\psi_\rho a_y}{r}\,dy,\qquad
 W_\rho=\iint\frac{\psi_\rho a_y|d|^2}{r^3}.
\]
The bulk and transition are bounded, respectively, by
\[
 \frac{1+\sqrt2}{8\pi}\sqrt{\|\Psi\|_\infty W_\rho Q},
 \qquad \frac{\|\eta'\|_\infty A_*Z}{4\pi\rho^3}.
\]
The latter is $o(n)$ after event normalization at the canonical radii.
Near/far optimization gives
\[
 \|\Psi\|_\infty\le3\pi^{1/3}(UY)^{1/3}.
\]
The sufficient strict normalized $\Psi W/Z$ slope is
$96\pi^2\nu\log2/(3+2\sqrt2)$ after spending $3\nu Q/4$.
Its source remains unpaid. Also $2W\le B_G$, with an inequality
because the full cubic $G$ contains an additional increment term.

The coarse static conversion $W\le C\|\Psi\|_\infty Q$ fails
on allowed compact solenoidal rescalings at an actual positive threshold:
\[
 u_\epsilon=\epsilon^{-1}U(x/\epsilon),\qquad
 A_*\sim\epsilon,\quad Z\sim\epsilon^{-1},\quad Q\sim\epsilon^{-3}.
\]
One-zero-endpoint pairs give
\[
 W\gtrsim\epsilon^{-3}\log(\rho/(C\epsilon)),\qquad
 \|\Psi\|_\infty Q\lesssim\epsilon^{-3},\qquad
 \rho\sim\epsilon^{1/3}.
\]
The radius here includes the actual initial choice.
Those pairs have zero signed cross core and inactive endpoint gradient.
Only the unrestricted absolute conversion is rejected; signed, active
and time-integrated source laws remain open.

\section{A scoped convex obstruction and balanced escape}
A finite nonnegative symmetric separately convex carrier $H$ with
$H(X,X)=H(0,Y)=0$ vanishes at $H(tY,Y)$ for every $t\ge0$.
Finite-convex continuity then gives
\[
 H(X+tY,Y)\le H(X,Y),\qquad t\ge0.
\]
For noncollinear $X,Y$, the angular supremum
$|X\times Y||X+(t-1)Y|/2$ diverges. Thus this strict class cannot
uniformly majorize the pointwise readout. Local compact velocity-curl
patches realize the pairs, but imply no integrated trajectory obstruction.

The ambient quadratic
\[
 H_\alpha(V,d)=\alpha|V|^2+\alpha^{-1}|d|^2
\]
is jointly convex in independent $(V,d)$ and separately convex in
original endpoints. Its nonzero single-zero-endpoint value escapes the
strict class. With $\alpha=\rho^2/(3\nu)$ it recovers the previous
viscosity allocation. The complete variable-parameter law retains
$\alpha'B-\alpha^{-2}\alpha'D$.
Moreover $(1-\operatorname{sinc}z)/z^2$ is nonincreasing globally.
Consequently $D_\rho/\rho^2$ decreases with radius, whereas
$\rho^2B_\rho$ increases. Shrinking the radius produces opposing
component works, with no automatic sign for their balanced sum.

\section{Favorable quadratic Fourier diffusion and logarithmic readout}
For the scalar multiplier $T$ with symbol $\Lambda_\rho$, let
$A=u\cdot\nabla$ and $F=S\omega$. The exact quadratic law is
\[
 D'=2\langle T\omega,F\rangle
   -\langle[T,A]\omega,\omega\rangle
   -2\nu\int k^2\Lambda_\rho(k)|\widehat\omega|^2.
\]
The full physical viscosity is nonpositive. Profile integration by
parts proves $k\Lambda'_\rho=8\pi\int(-\eta')(1-\operatorname{sinc})
\ge0$. For $X=D/Z$, the viscous term is
$-2\nu\operatorname{Cov}_{\mu}(k^2,\Lambda_\rho)\le0$,
where $d\mu=|\widehat\omega|^2d\xi/Z$. Both original nonlinear
terms remain unpaid. This specific quadratic sign says nothing about
an arbitrary nonlinear pair carrier.

An explicit envelope and Jensen give
\[
 \Lambda_\rho(k)\le16\pi\log(1+k^2\rho^2/3),\qquad
 \frac DZ\le16\pi\log\left(1+\frac{\rho^2Q}{3Z}\right).
\]
For $b=B/Z$, $a=9b/(16\pi)$ and $c=\rho^2/3$, a safe scalar
allocation against $\lambda\nu Q/Z$ leaves
\[
 R_\lambda=\frac{\sqrt a}{2}
 \left[1+\log_+\left(\frac{c\sqrt a}{2\lambda\nu}\right)\right],
 \qquad R_\lambda(0)=0.
\]
At $\lambda=3/4$, take its minimum with
$\rho^2B/(32\pi\nu Z)$. A strict event slope below $\log2/2$
for this residual is sufficient. The allocated viscosity cancels the
matching coercive term; native energy does not pay $\int Q/Z$.

For the original velocity increment $D^u=\langle u,Tu\rangle$,
the coherent pressure cancels by a scalar relative-hole multiplier
preserving divergence freedom. Actual compact product/Riesz pressure
regularity justifies the limit. Fourier curl gives
\[
 (D^u)'=-\langle[T,A]u,u\rangle-2\nu D^\omega.
\]
The cubic velocity source remains unpaid at a fixed radius. A logarithmic
weight is not the earlier polynomial strong-tail or heat cone-tip
consumer; its return requires a separate adapter.

\section{A genuinely native adaptive velocity source}
The original velocity drift $I_\rho$ has a finite full-profile bound.
Translation, vector Sobolev with $S=4/\sqrt3$ and interpolation give
\[
 \|\delta_hu\|_3^3\le(2S)^{3/2}|h|^{3/2}Z^{3/2}.
\]
Including every radial derivative and relative-hole transition yields
\[
 |I_\rho|\le K_\eta\sqrt\rho\,Z^{3/2},\qquad
 K_\eta=14\pi(2S)^{3/2}\int_0^2\eta(s)s^{-1/2}\,ds.
\]
The hole transition tends to zero at rate $\sqrt{\text{hole}}$ on
actual compact intervals. No terminal higher norm enters this estimate.

Fix $R,\alpha>0$ and use the global actual running maximum
$\bar Z(t)=\max_{s\le t}Z(s)$ to choose
\[
 \rho(t)=\min\{R,\alpha\nu^2/\bar Z(t)\}.
\]
The zero datum is separate. This positive, locally absolutely continuous
radius decreases on the entire existing path, including every physical
gap. Since $Z\le\bar Z$, native energy supplies the absolute source price
\[
 \int|I_{\rho(t)}|\,dt\le\frac{K_\eta\sqrt\alpha}{2}E_0.
\]
The complete original moving law therefore gives
\[
 \begin{split}
 D^u_{\rho(t)}(t)+2\nu\int_0^tD^\omega_{\rho(s)}\,ds
   +\int_0^t(-\rho')\partial_\rho D^u_\rho\,ds
 \le D^u_{\rho(0)}(0)+\frac{K_\eta\sqrt\alpha}{2}E_0.
 \end{split}
\]
The true initial term is at most $(8\pi/3)\rho(0)^2Z_0$.
Only nonnegative time integrals pass by monotone convergence at a
hypothetical finite endpoint; no terminal solution or trace is assumed.

The symbol lower bound $\Lambda_\rho(k)\ge4\pi\log_+(k\rho/2)$
pays a logarithmic enstrophy dose above the growing cutoff $2/\rho$.
This closes a specific original velocity source. Its radius decreases
like $\bar Z^{-1}$, while the canonical core radius decreases like
$((j+4)\bar Z)^{-1/3}$. The intermediate annulus and the fixed-cutoff
strong continuation consumer remain unpaid. Critical initial scaling
and a linear heat control give an unweighted dose of order $L^{-1}$
per natural scale; they are not actual nonlinear escape events.
The native estimate is a proved rung with an open original return.

\section{Construction frontier}
The full original strain equation retains pressure Hessian, vorticity and
two eigenframe diffusion contributions with opposite signs. Compact
solenoidal data with fixed energy and a fixed local jet can have arbitrarily signed initial
middle-strain production; the fixed neighborhood may shrink and global high
norms are not fixed. This rules out only a local instantaneous production
bound from those observations. For a fixed radius $r_0>0$, the remote
directional pressure Hessian is energy-paid by
$168E_0/(4\pi r_0^5)$; local pressure remains unbounded by this law.

Signed triad conservation and receiving/advected symmetrization expose
continuous weight differences. The clamped weight's fixed inner jump is
already paid cumulatively by kinetic net-flux bounds. The unresolved source
is the continuous weighted transfer uniformly over outer cutoff and the
full preterminal horizon. None of these components supplies that budget.

Seek a signed original-equation high-frequency interaction law or an
independently integrable production coefficient. Retain the proven written
conditional adapters and their separate formal implementation obligations.
Finite periodic controls and adviser suggestions remain exploratory evidence.
Every prospective rung needs exact applicability, law and return obligations,
source citations, controls, and independently reviewed proof bodies.
\section{Weighted adaptive parameter absorption}
Write $b=\bar Z$ and choose the whole-path radius
$\rho=\alpha\nu^2/b$ for nonzero initial data. The original native
strain and vector Sobolev estimates yield
\[
 Z'\le C_gZ^3/\nu^3,\qquad
 0\le b'\le C_gb^3/\nu^3,\qquad
 C_g=\frac{27}{128}(C_3S_0^{3/2})^4.
\]
The maximum is locally absolutely continuous. Its derivative is zero
off contact and agrees with $Z'$ almost everywhere on $Z=b$.
This is a comparison envelope, without an integrated $Q$ budget.

For $M=bD^u_\rho$, the full original chain rule is
\[
 M'+2\nu bD^\omega_\rho
 =bI_\rho+b'\int[\Lambda(z)-z\Lambda'(z)]|\widehat u|^2\,d\xi,
 \qquad z=k\rho.
\]
For $z\le1$, profile support gives $v=zs\le2$ and
\[
 2\operatorname{sinc}v-\cos v-1
 =2\cos(v/2)[\operatorname{sinc}(v/2)-\cos(v/2)]\ge0.
\]
Thus the parameter symbol is nonpositive there. For $z\ge1$, the
proved monotonicity of $\Lambda$ gives its positive part at most
$\Lambda\le z^2\Lambda$. Consequently the positive parameter work
is at most $C_g\alpha^2\nu bD^\omega_\rho$. For
$0<\alpha\le C_g^{-1/2}$,
\[
 M(t)+\nu\int_0^t bD^\omega_\rho
 \le M(0)+\int_0^t bI_\rho,
 \qquad 0\le M\le\frac{8\pi}{3}\alpha^2\nu^4.
\]
All physical gaps and the initial trace are retained. Only nonnegative
time integrals pass by monotone convergence to a hypothetical endpoint;
no endpoint trace or positive limiting radius is assumed.

An independent uniform upper bound $B$ for the actual signed source
on one finite preterminal horizon would therefore pay
\[
 \int_0^t b\int\log_+(k\rho/2)|\widehat\omega|^2\,d\xi\,ds
 \le\frac{M(0)+B}{4\pi\nu}.
\]
A separate uniform positive charge on all sufficiently high disjoint
actual growth events within that same horizon would exclude a finite
unbounded-enstrophy endpoint. Both suppliers remain open. Native energy
pays $\int Z$, whereas the weighted drift envelope leaves $\int bZ$.
Bounded $M$ alone gives no source bound. The critical heat calibration
makes the weighted dose scale invariant; it supplies no actual event.

\section{Native constant and actual event duration}
Physical Fourier transversality gives the exact original identity
$\|S\|_2^2=Z/2$. Vector interpolation therefore improves the native
production constant to
\[
 P\le C_0Z^{3/4}Q^{3/4},\quad
 C_0=S_0^{3/2}/\sqrt2,\quad C_0^4=1024/27.
\]
The original scalar optimization gives $Z'\le8Z^3/\nu^3$ and the
same causal maximum bound with coefficient 8. Thus the previously
proved weighted parameter absorption permits
$\alpha\le1/(2\sqrt2)$. This changes the constant without supplying
the weighted nonlinear source or the original strong-tail consumer.

Split original vorticity at physical frequency $k=1/\rho$, keeping
full original strain. Optimizing only the high-frequency dissipation
and favorably dropping the low part gives
\[
 Z'\le4C_0Z^{3/4}Q_{\rm lo}^{3/4}
       +128Z^2Z_{\rm hi}/\nu^3.
\]
The inner profile and its proved monotone symbol give
\[
 \Lambda(z)\ge\frac{19\pi}{30}\min(z^2,1),\qquad
 D^\omega_\rho\ge\frac{19\pi}{30}
                  (\rho^2Q_{\rm lo}+Z_{\rm hi}).
\]
On an actual doubling interval $J$ of level $h$, put
\[
 X=b(\rho^2Q_{\rm lo}+Z_{\rm hi}),\quad
 \xi=\nu^{-3}\int_JX,\quad \tau=h^2|J|/\nu^3,\quad
 A=4C_0 2^{5/4}.
\]
Using only $h\le Z\le2h$ and $Z\le b\le2Z$ on $J$, the exact
logarithmic endpoint difference and time H\"older give
\[
 \log2\le A\alpha^{-3/2}\xi^{3/4}\tau^{1/4}+128\xi.
\]
Consequently the actual weighted event dose is at least
\[
 \int_J bD^\omega_\rho\ge\frac{19\pi\nu^3}{30}
 \min\left\{\frac{\log2}{256},
       \left(\frac{\log2}{2A}\right)^{4/3}
                    \alpha^2\tau^{-1/3}\right\}.
\]
This is a uniform positive charge on every class with bounded critical
duration $\tau$. Long events remain open. Geometric levels with scalar
$\tau_j=2^{j/2}$ admit both a finite energy-duration sum and summable
vanishing charges. This is a scalar compatibility control, without an
actual NS realization.

\section{An actual-data control for source-alone absorption}
The periodic curl-generated seed
\[
 u=(\sin y,0,\sin x-\sin(x+y)),\qquad
 A=(0,-\cos x+\cos(x+y),-\cos y)
\]
has production density $ac(b-c)$, where $a=\cos y$, $b=\cos x$ and
$c=\cos(x+y)$, with mean $1/4$. For a nonnegative nonzero compact
smooth $\chi$, the whole-space field $U_R=\operatorname{curl}
[\chi(x/R)A]$ satisfies
\[
 P(U_R)=\frac{R^3}{4}\int\chi^3+o(R^3)>0
\]
for large fixed $R$. Cutoff derivative errors are only $O(R^2)$.
Thus this construction supplies allowed compact solenoidal Schwartz
data, without using a periodic field as whole-space data.

For a fixed such $U$, let
$c_\eta=(4\pi/3)\int_0^2\eta(s)s\,ds$. The exact initial-face law
and Fourier Taylor limit give
\[
 D^u_\rho/\rho^2\longrightarrow c_\eta Z,\quad
 D^\omega_\rho/\rho^2\longrightarrow c_\eta Q,\quad
 I_\rho/\rho^2\longrightarrow2c_\eta P.
\]
At data $LU$ and their actual radius
$\rho_L=\alpha\nu^2/(L^2Z_U)$,
\[
 \frac{I_{\rho_L}[LU]}{\nu D^\omega_{\rho_L}[LU]}
 \sim\frac{2LP_U}{\nu Q_U}\longrightarrow+\infty.
\]
This rejects only a uniform pointwise source-alone absorption constant.
Full parameter cancellation, datum-dependent bounds and cumulative signed
source laws remain available. Exact rational Fourier algebra verifies the
seed curl, divergence, mean production and phase reversal; it is not a
nonlinear PDE computation.


\section{Actual advected shell source and its time currency}
For the original compact-time divergence-free drift, the auxiliary vector
propagator $U(t,s)$ can be constructed by heat Duhamel in $L^1$ and $L^2$.
Kato's vector inequality gives an $L^1$ contraction, and energy gives
$Y'\le-2\nu\|\nabla f\|_2^2$. The independently derived Sobolev constant
$S_0=4/\sqrt3$ yields Nash's inequality
\[
 \|\nabla f\|_2^2\ge S_0^{-2}\|f\|_1^{-4/3}\|f\|_2^{10/3}.
\]
Integration, the reversed adjoint and half-interval composition give
\[
 \begin{gathered}
 \|U(t,s)\|_{1\to2},\ \|U(t,s)\|_{2\to\infty}
 \le\left(\frac4{\nu(t-s)}\right)^{3/4},\\
 \|U(t,s)\|_{1\to\infty}
 \le\left(\frac8{\nu(t-s)}\right)^{3/2}.
 \end{gathered}
\]
Constants are independent of drift bounds used to justify construction.
No Fourier support preservation or terminal coefficient bound is assumed.

For the actual shell $v_j=\Delta_j u$, $\lambda=2^j$, its exact equation is
\[
 (\partial_t+u\cdot\nabla-\nu\Delta)v_j=R_j,
 \qquad R_j=[u\cdot\nabla,\Delta_j]u-\nabla p_j.
\]
Both source terms have full two-increment representations with Schwartz
kernel gradients. The pressure kernel is the localized order-zero tensor
symbol, followed by its gradient; no $L^1$ Riesz estimate is used.
The bound $\|\delta_yu\|_4^2\le(2S_0)^{3/2}|y|^{1/2}Z$ gives
\[
 \|R_j\|_2\le C_R\lambda^{1/2}Z,\qquad
 \int\|R_j\|_2\le C_R\lambda^{1/2}\frac{E_0}{2\nu}.
\]
Here $C_R$ is the sum of the full weighted kernel-gradient moments.
Actual Duhamel and the adjoint bound therefore imply
\[
 \lambda^{-1/2}\|v_j(t)\|_\infty
 \le\frac{4^{3/4}\lambda^{-1/2}\|v_j(\sigma)\|_2}
          {[\nu(t-\sigma)]^{3/4}}
 +\frac{4^{3/4}C_R}{\nu^{3/4}}
       \int_\sigma^t(t-s)^{-3/4}Z(s)\,ds.
\]
Fixed positive lag and low-frequency Bernstein bounds pay the opening
observation. The convolution has weak $L^{4/3}$ time norm at most
$4\sqrt2\int Z$. Indeed splitting $k(t)=t^{-3/4}$ at height
$a=A/(2\int Z)$ yields
\[
 |\{k*Z>A\}|\le8\,2^{1/3}\left(\frac{\int Z}{A}\right)^{4/3}.
\]
This is not the critical strong $L^4$ currency. A scalar unit-mass pulse
$Z_\epsilon=\epsilon^{-1}\chi(t/\epsilon)$ makes the strong $L^4$
convolution norm grow like $\epsilon^{-1/2}$. This tests only that positive
convolution estimate, without an NS realization. Earlier heat constructions
already considered the $4/3$ currency; the present gain is an actual advected
adapter and pressure-complete source, rather than a new time exponent.

\section{Retained parameter work gives a one-source continuation return}
Return to $b=\max_{s\le t}Z(s)$, $\rho=\alpha\nu^2/b$,
$0<\alpha\le1/(2\sqrt2)$ and $M=bD^u_\rho$. Put
\[
 J(z)=\Lambda(z)-z\Lambda'(z),\qquad
 N=b'\int[J(k\rho)]_-|\widehat u|^2\,d\xi\ge0.
\]
The full balance retains this formerly discarded favorable work:
\[
 M'+2\nu bD^\omega_\rho+N
 =bI_\rho+b'\int[J(k\rho)]_+|\widehat u|^2\,d\xi,
 \qquad M'+\nu bD^\omega_\rho+N\le bI_\rho.
\]
For $0\le v\le2$, the decreasing alternating series gives
\[
 2\operatorname{sinc}v-\cos v-1
 \ge v^2/6-v^4/40\ge v^2/15.
\]
Since $v=zs$ with $0\le s\le2$, the inner unit profile yields
$[J(z)]_-\ge c_mz^2$ for $z\le1$, $c_m=4\pi/15$.
At almost every time with $b'>0$, contact gives $Z=b$. Split actual
vorticity into $Z_{\rm lo}+Z_{\rm hi}=b$ at $k\rho=1$; this sharp split is
not differentiated. The low negative work and high viscosity satisfy
\[
 N\ge c_m\alpha^2\nu^4\frac{Z_{\rm lo}}b\frac{b'}b,
 \qquad
 \nu bD^\omega_\rho\ge\frac{19\pi\nu^4}{240}
                      \frac{Z_{\rm hi}}b\frac{b'}b.
\]
Thus on the whole path, including noncontact gaps,
\[
 \nu bD^\omega_\rho+N\ge\kappa\nu^4\frac{b'}b,
 \qquad\kappa=\frac{4\pi\alpha^2}{15}.
\]
Positive locally absolutely continuous $b$ and $M\ge0$ now give
\[
 \kappa\nu^4\log\frac{b(t)}{b_0}
 \le M_0+\int_0^t b(s)I_{\alpha\nu^2/b(s)}(s)\,ds.
\]
An independently finite upper bound $B_T$ for the signed integral, uniform
in every existing $t<\min(T,T_*)$, would imply
\[
 b(t)\le b_0\exp\!\left(\frac{M_0+B_T}{\kappa\nu^4}\right),
 \qquad M_0\le\frac{8\pi\alpha^2\nu^4}{3}.
\]
At a hypothetical finite endpoint native energy then pays $\int Z^2$,
and the original strong continuation and full closed-face return apply.
For nonzero solenoidal $L^2$ data $Z_0>0$ by Fourier curl/divergence;
zero data are separately zero. No terminal trace is used.

This improves the conditional return to one unresolved source child.
The separate all-event duration/dose child is unnecessary for this route;
its historical alternative remains open. The signed source bound is still
\emph{unproved}, and bounded $M$ or native $\int Z$ does not supply it.
Consequently unconditional Clay closure remains open.


\section{Actual local source bursts and strict payment}
Choose the compact-curl $U$ with $P_U>0$ constructed above, and keep
$\nu,\alpha,\eta$ fixed. The original data
$V_\lambda=\lambda^{3/2}U(\lambda x)$ have fixed kinetic energy and
\[
 Z_\lambda=\lambda^2Z_U,\quad Q_\lambda=\lambda^4Q_U,\quad
 P_\lambda=\lambda^{9/2}P_U.
\]
For large $\lambda$, the actual initial derivative is positive and
$b'=Z'$. At the prescribed radius, the effective radius on $U$ is
$\epsilon=\lambda\rho=\alpha\nu^2/(\lambda Z_U)$. Since
$J(z)=-c_\eta z^2+O(z^4)$, exact spatial scaling gives
\[
 bI_\rho\sim N\sim
     2c_\eta\alpha^2\nu^4\lambda^{5/2}P_U/Z_U,
 \qquad \nu bD^\omega_\rho=O(\lambda^2).
\]
Thus a fixed $\theta<1$ in $bI\le\theta N+C\nu bD^\omega+C_EZ$
cannot hold uniformly over this fixed-energy family with finite
$\lambda$-independent constants. This is initially a pointwise control.

A genuine shrinking-time control follows using positive-viscosity paths
\[
 u_\lambda(x,t)=\lambda^{3/2}v_\mu(\lambda x,\lambda^{5/2}t),
 \qquad\mu=\nu/\sqrt\lambda.
\]
For the fixed seed $U$, the actual NS$_\mu$ paths have a common finite
$H^8$ window. Indeed differentiated transport and pressure cancel, and in
each remaining product of total order at most nine one factor has order
at most four and is controlled in $L^\infty$, the other in $L^2$.
The resulting $y'\le Cy^2$ bound and the original positive-viscosity norm
continuation alternative give a common window without an Euler limit.
Uniform time continuity in $H^6$ keeps $P>0$, $Z>0$ and $Q$ bounded;
for small $\mu$ contact persists throughout a smaller window $[0,\delta]$.

Uniform increment Taylor estimates give
$I_\epsilon[v]=2c_\eta\epsilon^2P_v+O(\epsilon^3)$.
The full kernel derivative, including its outer transition, is retained.
Consequently on physical time $[0,\delta\lambda^{-5/2}]$,
\[
 \int bI=2c_\eta\alpha^2\nu^4\int_0^\delta P_v/Z_v\,d\tau
                    +O(\lambda^{-1}),
\]
\[
 \int N=c_\eta\alpha^2\nu^4
          \int_0^\delta(2P_v-2\mu Q_v)/Z_v\,d\tau
                    +O(\lambda^{-2}).
\]
Both are uniformly bounded above and below by positive constants;
their difference tends to zero. No numerical Euler-limit value is needed.
Meanwhile $\int Z$ and $\nu\int bD^\omega$ are $O(\lambda^{-1/2})$.
This excludes the corresponding strict $\theta<1$ window payment
\emph{without an initial or additive credit}. A positive initial credit
can pay this one finite burst. No repeated-burst trajectory, divergent
full-horizon source or blowup is constructed.

At coefficient one, $R=bI-N$ obeys
$M'+\nu bD^\omega\le R$. A bound on its cumulative integral would pay
viscosity, but would remove the low-frequency log-growth charge $N$.
An independent low-contact return or a controlled credit recurrence is
therefore still required for that modified route. The full original
signed source and unconditional target remain open.

\section{Two paid full-drift source sectors}
For the actual shell $v_j=\Delta_j u$, set $\lambda=2^j\ge1$,
$m=\lambda^{3/5}/16$, $M=8\lambda^2$, and use a fixed real even smooth
cutoff $0\le\chi\le1$, equal to one on the unit ball and zero outside
the radius-two ball. Write $a=S_m u$, $v=S_Mu$, $b_j=v-a$, $w=u-v$.
The algebraic pressure-complete source is
\[
 R[q]=q\cdot\nabla\Delta_jq-\mathbb P\Delta_j\operatorname{div}(q\otimes q),
 \qquad R[u]=L_j+F_j+R[b_j],
\]
where $L_j=R[v]-R[b_j]$ and $F_j=R[u]-R[v]$.
Every response is propagated by the same full-$u$ drift $U(t,s)$.
The filtered $b_j$ is not asserted to solve autonomous NS.

Let $z$ be the sharp restriction of original $u$ to
$\lambda/4\le|\xi|\le4\lambda$. Support of the low cross products makes
their other input equal to $z$; the unprojected term uses
$\Delta_j b_j=\Delta_j z=v_j$. Equality of the two cross divergences
gives the exact pressure reflection
\[
 L_j=[a\cdot\nabla,\Delta_j]z+(I-2\mathbb P)\Delta_j(z\cdot\nabla a).
\]
With $C_K=\int|y||\nabla K_0(y)|\,dy$ and
$C_\phi=\|\phi_0\|_\infty$, the reflection has $L^2$ norm one, and
kernel, Bernstein and annular gradient estimates imply
\[
 \|L_j\|_2\le C_{\rm low}m^{5/2}\lambda^{-1}\sqrt{E_0Z},\qquad
 C_{\rm low}=4(C_K+C_\phi)\sqrt{4\pi2^5/5}.
\]
No sharp projection is used in $L^1$.

The far difference retains the unprojected high sweeping:
\[
 F_j=w\cdot\nabla v_j-\mathbb P\Delta_j\operatorname{div}T_{\rm hi}.
\]
If $v_{\ge M/2}$ is the sharp high part of $v$, then
$T_{\rm hi}=w\otimes w+w\otimes v_{\ge M/2}+v_{\ge M/2}\otimes w$.
Only the projected very-low/far products vanish by support; all smooth
transition values remain. Since
$\|w\|_2\le\sqrt Z/(2\pi M)$ and each other factor is energy bounded,
\[
 \|F_j\|_2\le C_{\rm far}\lambda^{5/2}M^{-1}\sqrt{E_0Z},\qquad
 C_{\rm far}=6C_\phi\sqrt{32\pi/3}+C_\phi\sqrt{4\pi2^5/5}.
\]
Output differentiation prices the projected tensor by $\lambda$;
sweeping is paid separately. Its output need not stay in the shell.

An independent causal time law supplies
\[
 \left\|\int_0^\infty r^{-3/4}|f(t-r)|\,dr\right\|_4\le C_H\|f\|_2.
\]
Here the explicit constants are
\[
 C_H=(C_N+\sqrt2)(4\sqrt6)^{1/2},\qquad
 C_N=\frac{2^{3/4}}{1-2^{-1/4}}.
\]
Indeed centered maximal weak $(1,1)$ constant three and truncation at
half the level give its $L^2$ norm at most $\sqrt{24}$. The one-sided
maximal norm is at most $4\sqrt6$. Dyadic integration up to $R$ costs
$C_NR^{1/4}Mf$, while the tail costs
$\sqrt2R^{-1/4}\|f\|_2$. Choosing
$R^{1/2}=\|f\|_2/Mf$ and using the maximal $L^2$ bound proves the law.
Zero cases and infinite kernels follow by positive truncation.

For $n_{{\rm low},j}=\int_0^tU(t,s)L_j(s)\,ds$ and the analogous
far response, actual Nash plus $f=\sqrt Z$ and native energy imply
\[
 \left\|\sup_{j\ge0}\lambda^{-1/2}\|n_{{\rm low},j}\|_\infty\right\|_4
 \le\frac{4^{3/4}C_{\rm low}16^{-5/2}C_HE_0}{\sqrt2\nu^{5/4}},
\]
\[
 \left\|\sup_{j\ge0}\lambda^{-1/2}\|n_{{\rm far},j}\|_\infty\right\|_4
 \le\frac{4^{3/4}(C_{\rm far}/8)C_HE_0}{\sqrt2\nu^{5/4}}.
\]
A common positive envelope dominates the supremum before the time norm.
These bounds are uniform through every existing preterminal horizon;
no shell preservation by $U$ is assumed. The spectral coefficients do
not assert contraction after profile constants. The opening is treated
at positive delay and by original local existence at time zero.

The exact stress rewrite is also retained:
\[
 \tau_j=\Delta_j(u\otimes u)-u\otimes v_j-v_j\otimes u,\qquad
 R_j=-v_j\cdot\nabla u-\operatorname{div}(\tau_j+p_jI).
\]
It gives $\|\tau_j\|_2\le C_\tau\lambda^{-1/2}Z$ and exact adjoint
work with a spatial derivative of $U(t,s)^*q$. Scalar Nash does not pay
that derivative; even the zero-drift point-test norm scales as
$(\nu(t-s))^{-5/4}$. This is a calibration of an absolute estimate,
not a prohibition of signed cancellation. The complete middle response,
its original critical continuation consumer, and the independent adjoint
price remain open. Paying these two sectors does not establish Clay(A).

\section{Deficit mass, records and actual adjoint controls}
Set $c_\eta=(4\pi/3)\int_0^2\eta(s)s\,ds$ and
$c_4=(\pi/15)\int_0^2\eta(s)s^3\,ds$. The exact Fourier deficit
\[
 \delta(z)=c_\eta z^2-\Lambda(z),\qquad 0\le\delta(z)\le c_4z^4
\]
is globally nonnegative. Indeed $1-\operatorname{sinc}v
=\int_0^1(1-\cos(tv))\,dt$, and twice integrating $1-\cos x$ gives
$0\le x^2/2-1+\cos x\le x^4/24$. Its moving symbol has a global sign:
\[
 K(z)=z\delta'(z)-\delta(z)=c_\eta z^2+J(z)\ge0,
\]
because its scalar integrand is
\[
 v^2/6+1+\cos v-2\operatorname{sinc}v
 =v^2\int_0^1 t(1-t)(1-\cos(tv))\,dt\ge0.
\]
These are global identities, rather than a low-frequency truncation.

For $b=\max_{s\le t}Z(s)>0$ and $\rho=\alpha\nu^2/b$, put
$W=b\int\delta(k\rho)|\widehat u|^2\,d\xi$ and
$D^\omega_\delta=\int\delta(k\rho)|\widehat\omega|^2\,d\xi$.
Subtracting the original velocity increment law from its exact quadratic
enstrophy counterpart, with all parameter derivatives, gives
\[
 W'+2\nu bD^\omega_\delta+b'\int K(k\rho)|\widehat u|^2\,d\xi
 =b(2c_\eta\rho^2P-I_\rho),
\]
\[
 M+W=c_\eta\alpha^2\nu^4 Z/b,
 \qquad0\le W\le c_\eta\alpha^2\nu^4.
\]
The complete signed deficit production is unpaid. All gaps are included;
bounded mass does not upper-bound cumulative positive work.

The actual whole-path quotient separately obeys
\[
 \log(b/b_0)+2\nu\int_0^t Q/b
 =2\int_0^t P/b+1-Z(t)/b(t).
\]
This uses $Zb'/b^2=b'/b$ only where $b'$ is supported. At any actual
record endpoint $Z(t)=b(t)=B$, every $0<h\le t$ satisfies
\[
 \int_{t-h}^t Q\le\frac{C_0^4B^3h}{\nu^4}.
\]
Indeed the integrated enstrophy law has nonnegative endpoint difference
$B-Z(t-h)$, and H\"older bounds production by
$C_0B^{3/4}h^{1/4}(\int Q)^{3/4}$. Interior gaps are allowed. This
does not provide an unweighted terminal $Q$ budget.

Retaining viscosity in absolute production gives
$2|P|\le\nu Q+64\nu^{-3}Z^3$.
For $F=Z/b^2$, the exact record term integrates to
$2(1/b_0-1/b(t))\ge0$. Native energy therefore supplies
\[
 \int Q/b^2\le\frac1{\nu b_0}+\frac{32E_0}{\nu^5},\qquad
 \int |P|/b^2\le\frac1{2b_0}+\frac{32E_0}{\nu^4}.
\]
This refines the earlier reciprocal currency; it does not pay the missing
factor of $b$. The global quartic deficit then genuinely pays
\[
 \int_0^t bW\le c_4\alpha^4\left(\frac{\nu^7}{b_0}+32\nu^3E_0\right).
\]
Precisely this mass-time quantity is paid. Its derivative, viscosity,
parameter work, signed source and quadratic low-contact log return need
separate adapters.

On contact let $s=\nu^3b'/(8b^3)$ and $q=\nu^4Q/b^3$.
The necessary envelope $s\le(C_0q^{3/4}-q)/4\le1$ does not give small
$q$ from small $s$. This also has an actual smooth-data control:
for the compact seed $P_U>0$, choose
$u_0=a_\epsilon U$, $a_\epsilon=\nu Q_U(1+\epsilon)/P_U$.
Actual initial growth is positive, and on sufficiently nearby positive
contacts
\[
 q_\epsilon\longrightarrow\frac{P_U^4}{Q_U^3Z_U^3}>0,
 \qquad s_\epsilon=\epsilon q_\epsilon/4\longrightarrow0.
\]
Each datum has its own local contact interval. This rejects only the
universal inference of small dissipation from slow normalized growth.

An actual full-drift adjoint can also lose solenoidality. For
$\psi(s)=U(t,s)^*q$ and $d=\operatorname{div}\psi$,
\[
 -d_s-u\cdot\nabla d-\nu\Delta d
 =\sum_{i,k}\partial_i u_k\,\partial_k\psi_i.
\]
Compact-curl data and terminal tests with local forms
$u_0=(x_2,0,0)$ and $q=(0,x_1,0)$ give coefficient one near the origin.
At sufficiently small actual terminal time the coefficient stays nonzero,
so a solenoidal terminal test develops divergence backward. For the middle
velocity $h$, its complete source pairing retains
\[
 \langle R[h],\psi\rangle
 =\langle h\otimes h,\nabla\Delta_j\mathbb P\psi\rangle
 -\langle v_j\otimes h,\nabla\psi\rangle.
\]
Thus terminal solenoidality cannot delete the gradient test sector or
pressure work automatically. The alternative solenoidal projected
propagator has $L^2$ contraction, but its nonlocal generator has no proved
scalar Kato/Nash adapter here.

The weighted source criterion has the strength of a bounded record.
Its log return implies bounded $b$. Conversely, if $b\le B$,
the original drift estimate gives
\[
 \int b|I_\rho|\le K_\eta\sqrt\alpha\,BE_0/2.
\]
This converse assumes the bound being sought. It supplies no independent
estimate. The signed sources and the mass-time consumer remain open,
as does original Clay(A). Zero datum is treated separately before divisions.


\section{Direct resolving-shell consumer and same-path return}

Choose a nonnegative radial nonincreasing smooth cutoff $\chi$, equal to
one on the unit ball and zero outside the radius-two ball, and set
$\phi(\xi)=\chi(\xi)-\chi(2\xi)$, $\lambda_j=2^j$, and
$u_j=\phi(2^{-j}D_\xi)u$. This explicit homogeneous profile obeys
$\phi_j\ge0$, $\sum_j\phi_j=1$ away from zero, and
$\sum_j\phi_j^2\le1$. A merely arbitrary annular probe has no such
resolution assumption. Put
\[
 B(t)=\sup_{j\in\mathbb Z}\lambda_j^{-1/2}\|u_j(t)\|_\infty.
\]
Let $C_D$ be the $L^1$ norm of the gradient of a fixed reproducing kernel
for the radius-two Fourier ball. Then
$\|\nabla u_j\|_\infty\le C_D B\lambda_j^{3/2}$ and
$\|\nabla u_j\|_2\le4\pi\lambda_j\|u_j\|_2$ in the original
physical $2\pi$ convention. Original convection gives the exact cubic
\[
 P=-\sum_{i,j,k}\int\partial_k u_i\,\partial_k u_j\,\partial_j u_i.
\]
Sort its three shell indices as $\ell\le p\le q$. Nonzero interactions
have $q-p\le3$; six permutations overcount all orderings. The low
geometric sum is $\lambda_p^{3/2}/(1-2^{-3/2})$. Writing
$e_j=\lambda_j^{7/4}\|u_j\|_2$, the absolute cubic sum is at most
\[
 K_0B\sum_p\sum_{d=0}^3 2^{-3d/4}e_pe_{p+d},
 \qquad K_0=\frac{6C_D(4\pi)^2}{1-2^{-3/2}}.
\]
Cauchy for each shifted sequence and the resolving profile give
\[
 \sum_j e_j^2\le\pi^{-7/2}\int |k|^{7/2}|\widehat u|^2
 \le\pi^{-7/2}Z^{1/4}Q^{3/4}.
\]
Here $\lambda_j\le |k|/\pi$ on each shell. Thus
\[
 |P|\le C_* BZ^{1/4}Q^{3/4},\qquad
 C_*=K_0\pi^{-7/2}\sum_{d=0}^3 2^{-3d/4},
\]
\[
 Z'+\nu Q\le\frac{27C_*^4}{16\nu^3}B^4 Z.
\]
All cubic limits follow from homogeneous partial sums converging in the
existing compact $H^\gamma$ path, $\gamma>5/2$, hence in gradient
$L^2$ and $L^\infty$. No endpoint higher norm enters the constants.

A supplied finite horizon-uniform $L^4$ dose for
$B_{\rm high}=\sup_{j\ge0}\lambda_j^{-1/2}\|u_j\|_\infty$ therefore
has an exact original return. For $j<0$, native energy and Fourier volume
bound the normalized shell by
$C_\phi\sqrt{32\pi/3}\,\lambda_j\sqrt{E_0}$, so their full $B^4$
time cost is finite on a finite horizon. The original local initial
interval covers the opening Nash singularity. Gronwall gives
\[
 \sup Z\le Z_0\exp\left(\frac{27C_*^4}{16\nu^3}\int B^4\right).
\]
The retained inequality pays $\int Q$; alternatively native energy pays
$\int Z^2\le(\sup Z)E_0/(2\nu)$. The written strong-path continuation
and closed initial-face adapters then continue the \emph{same} maximal
path and return original Clay(A), conditionally. No weak endpoint
extension is assumed.

Reinstating the earlier source split with this admissible fixed profile,
the low and far responses are independently paid in the required time
currency. The full-middle critical response is the remaining open
applicability supplier for this route. The direct consumer closes an
analytic bridge; it does not pay that supplier or prove unconditional
Clay(A).


\section{Native normalized response and helical source interface}

The exact deficit from the preceding carrier calculations gives a
natively paid scalar response. For $b=\max_{s\le t}Z(s)>0$,
$\rho=\alpha\nu^2/b$, and $\delta\ge0$, put
$V=\int\delta(k\rho)|\widehat u|^2=W/b$ and
$H_\delta(z)=z\delta'(z)=K(z)+\delta(z)\ge0$. The full quotient law is
\[
 V'+2\nu D_\delta^\omega+
 \frac{b'}b\int H_\delta(k\rho)|\widehat u|^2
 =S_\delta=\frac{2c_\eta\alpha^2\nu^4P}{b^2}-I_\rho.
\]
Native record estimates and the original drift price yield
\[
 \int|S_\delta|\le L_\delta
 =\frac{c_\eta\alpha^2\nu^4}{b_0}
 +64c_\eta\alpha^2E_0+\frac{K_\eta\sqrt\alpha E_0}{2}.
\]
Every work term after $V'$ is nonnegative. Hence its total work is at
most $V_0+L_\delta$ and $\int|V'|\le V_0+2L_\delta$ uniformly in
preterminal cutoffs. This is a scalar bounded-variation statement.
Since $0\le V\le c_\eta\alpha^2\nu^4/b$, it permits $b\to\infty$
and $V\to0$. The quartic low-frequency symbol supplies no quadratic
contact charge. A critical or growing-weight return remains open.

A second original interface uses real helical projectors. For physical
$k=2\pi\xi$, let $D=|k|$, $C(k)v=ik\times v$, and
$\mathcal H_\pm=(\mathbb P\pm C/D)/2$ away from zero. These are
Hermitian transverse projections preserving real fields. Define
\[
 A_\pm=\int D|\widehat u_\pm|^2,\qquad
 Q_{3,\pm}=\int D^3|\widehat u_\pm|^2.
\]
Original helicity $\mathcal H=A_+-A_-$ obeys
\[
 \mathcal H'=-2\nu(Q_{3,+}-Q_{3,-}),\qquad
 A_+(t)+2\nu\int_0^tQ_{3,+}
 =A_-(t)+\mathcal H_0+2\nu\int_0^tQ_{3,-}.
\]
Thus supplied horizon-uniform $\sup A_-\le a_T$ and
$\int Q_{3,-}\le q_T$ imply
\[
 \sup(A_++A_-)\le2a_T+|\mathcal H_0|+2\nu q_T,
 \qquad \int(Q_{3,+}+Q_{3,-})
 \le\frac{a_T+|\mathcal H_0|}{2\nu}+2q_T.
\]
Fourier volume and H\"older give
$B^4\le [32/(3\pi)]^2(A_++A_-)(Q_{3,+}+Q_{3,-})$,
so this one-sector premise has the preceding same-path conditional return.
It remains unpaid. The two sector productions are equal and signed;
helicity does not make them nonpositive. Unequal-length positive-helical
inputs generate a negative-helical output. A Hermitian compact Fourier
packet lift realizes this on admitted real solenoidal Schwartz data and
on the actual initial time derivative. Ordinary dynamics does not preserve
the positive sector automatically.

The independently derived $\beta=5/4$ hyperviscous enstrophy control is
a changed-equation calibration: its native energy pays
$\int\|D^{5/4}u\|_2^2$, the coefficient of its enstrophy estimate.
That additional quarter derivative is unpaid in original viscosity.
Likewise an averaged nonlinearity retains an explicit solenoidal
remainder in original transport, curl and gauge identities; the operator
interface alone supplies no original-equation source estimate.

\section{Positive impulse metric and its unpaid physical readout}

On each existing original compact, set
$L=\partial_t+u\cdot\nabla-\nu\Delta$ and $A_{ij}=\partial_i u_j$.
The canonical scalar gauge $L\varphi=p-|u|^2/2$, $\varphi_0=0$,
gives $m=u+\nabla\varphi$, $Lm=-Am$ and $\mathbb Pm=u$.
Construct the specific symmetric matrix
\[
 LG=-AG-GA^T,\qquad G_0=I.
\]
Heat--Duhamel construction for $G-I$ and matrix cone comparison give
$e^{-2\int\|S\|_{\infty,\mathrm{op}}}I\le G
\le e^{2\int\|S\|_{\infty,\mathrm{op}}}I$ on each existing compact.
These local coefficient integrals are finite; their endpoint bounds are
unpaid. For $H=G^{-1}$ and $q=m^THm$, complete inverse/product calculus gives
\[
 LH=A^TH+HA-2\nu\sum_k(\partial_kH)G(\partial_kH).
\]
\[
 Lq=-2\nu\sum_k
 |\partial_km+G(\partial_kH)m|_H^2\le0.
\]
Consequently the native initial constants pay
\[
 q\le\|u_0\|_\infty^2,\qquad
 \int q(t)+2\nu\int_0^t\!\int\sum_k
 |\partial_km+G(\partial_kH)m|_H^2\le E_0.
\]
The spatial metric derivatives are essential. These estimates concern
inverse-weighted impulse; they give no unweighted endpoint bound without
an independent physical readout. Tracefree $A$ gives
\[
 L\log\det G=\nu\sum_k
 \|G^{-1/2}(\partial_kG)G^{-1/2}\|_F^2\ge0,
 \qquad\det G\ge1.
\]
A determinant lower bound does not control the largest eigenvalue.
For $y=\|G-I\|_2^2$ the exact matrix energy and original Sobolev constant
instead give
\[
 y'+\nu\|\nabla(G-I)\|_2^2
 \le(1+1024\nu^{-3}Z^2)y+2Z.
\]
Its $\int Z^2$ coefficient is unpaid by native energy. An actual compact
curl datum locally $(-ax_1,ax_2,0)$ has
$G_t(0)=\mathrm{diag}(2a,-2a,0)$; this is an initial deformation control,
not an example of singularity.

Let $K_j$ be the Schwartz matrix kernel of the final-time annular operator
$\Delta_j\mathbb P$, with uniform $L^1$ price, and define
\[
 D_{\rm ann}(t)=\sup_{j\ge0}\lambda_j^{-1/2}
 \big\|\|K_j\|_{\rm op}*\sqrt{\lambda_{\max}G(t)}\big\|_\infty.
\]
Then $B_{\rm high}\le\|u_0\|_\infty D_{\rm ann}$.
An independently paid horizon-uniform $\int D_{\rm ann}^4$ would give
full original Clay(A) through the preceding consumer and return.
A uniform bound $\lambda_{\max}G\le M_T$ is another sufficient premise,
yielding $\int B_{\rm high}^4\le
T C_{\rm ann}^4M_T^2\|u_0\|_\infty^4$. Both premises remain open.
No global Leray $L^\infty$ bound or drift/shell commutation is used.

Mass alone gives, for unit $v$,
$|v\cdot u_j(x)|^2\le E_0v^TR_j(x)v$ with
$R_j=\int K_j(z)G(x-z)K_j(z)^T\,dz$.
Even $G=I$ makes $\lambda_j^{-1}R_j$ grow as $\lambda_j^2$.
Thus this raw readout does not supply critical regularity. The exact
scalar $\theta=\log(\operatorname{tr}G/3)-\log\det G$ has favorable
diffusion, but retains signed strain and needs an upper determinant price
to control physical trace. Native covariant dissipation, correlated
annular readouts and the actual coupling $\mathbb Pm=u$ remain constructive
leads. Original unconditional closure remains open.


\section{Derived record limits and the remaining source}

Suppose the original maximal path had a finite endpoint $T_*$.
The preceding continuation law forces unbounded enstrophy $Z$.
Choose first hitting times $Z(t_k)=B_k\to\infty$ and physical centers $x_k$.
With unchanged viscosity set
\[
 \begin{aligned}
 U_k(y,s)&=B_k^{-1}u(x_k+y/B_k,t_k+s/B_k^2),\\
 \Pi_k(y,s)&=B_k^{-2}p(x_k+y/B_k,t_k+s/B_k^2).
 \end{aligned}
\]
On every fixed backward window $[-R,0]$, for sufficiently large $k$,
\[
 \begin{gathered}
 \sup Z_{U_k}\le1,\quad Z_{U_k}(0)=1,\quad
 \int_{-R}^0 Q_{U_k}\le KR,\quad Q_{U_k}(0)\le K,\\
 K=\frac{1024}{27\nu^4}.
 \end{gathered}
\]
The last bound follows from nonnegative first-hit derivative and the cubic
production bound. Scaled kinetic energy is $E_{U_k}=B_kE_u$ and can diverge.
Homogeneous Sobolev gives $\|U_k\|_6\le S_0=4/\sqrt3$.

These are sufficient to \emph{derive} local compactness. For canonical
pressure split $U_k\otimes U_k$ into a compact near source, whose double
Riesz multiplier is bounded on $L^2$, and a far source. The off-origin
kernel is
\[
 K_{ij}(z)=\frac{3z_iz_j-\delta_{ij}|z|^2}{4\pi|z|^5}.
\]
Its $L^{3/2}$ tail norm is $O(L^{-1})$, while the tensor source has bounded
$L^3$ norm. This gives uniform local pressure $L^2$ bounds. The equation
then bounds time derivatives in local $H^{-1}$; cube averages and
$\|f\|_2^2\le\|f\|_{H^{-1}}\|f\|_{H^1}$ give strong local
$C_tL^2$ compactness. The $Q$ budget gives strong spacetime local gradients,
and its terminal bound gives strong terminal local $H^1$ traces.
Local $L^4$ product convergence, the fixed near cutoff and then the far-tail
limit identify canonical pressure and the original ancient distributional
Navier--Stokes equation. No free harmonic force is added.

The ancient limit may be zero. Neither global normalization nor a terminal
high-frequency tail bound proves spatial concentration. A correct terminal
core condition uses a fixed ball \emph{after} recentering:
\[
 \int_{B_L(0)}|\nabla_yU_k(0)|^2\ge\epsilon>0,
 \quad\text{equivalently}\quad
 \int_{B_{L/B_k}(x_k)}|\nabla_xu(t_k)|^2\ge\epsilon B_k.
\]
A ball drifting in rescaled coordinates without recentering is insufficient.
For instance the profiles $W_L=L^{-1/2}W(\cdot/L)$ have $Z=1$, $Q=O(L^{-2})$
and vanish locally even after translations. These are norm controls,
not actual singular solutions.

A precise conditional nonzero bridge is available. If actual records obey
$Z_{U_k}(0)-Z_{U_k}(-R)\ge d$ for fixed $R,d>0$, signed production satisfies
$\int P_{U_k}\ge d/2$. Choose
$L=\max(1,(4C_0KR/d)^4)$. The bad-time price is
$\int_{Q>L}|P|\le C_0L^{-1/4}KR\le d/4$.
On unit cubes the direct cutoff Sobolev bound gives
\[
 |P|\le C_{\rm loc}\sqrt{\eta}(1+Q)^{3/4},\qquad
 \eta=\sup_j\int_{C_j}|\nabla U|^2.
\]
Consequently some good time and cube have a uniform positive gradient mass.
Cutoff integration by parts with $Q\le L$ gives a positive local velocity
norm. Moreover, the original projected equation gives
\[
 \int_{-R}^0\|\partial_sU_k\|_2^2\le
 (2\nu^2K+2S_0^3\sqrt K)R.
\]
Thus translation-invariant compactness preserves the velocity mass at a
subsequential time and gives a nonzero ancient limit. The fixed growth
supplier and an exact applicable rigidity theorem remain open.

\section{Covariant readout and source-local prices}

The native inverse-metric charge can retain correlation in the final
annular readout. Put $k_j=\|K_j\|_{\rm op}$, $C_K=\int k_0$, and
\[
 V=\sup_{j\ge0,x}C_K\lambda_j^{-1}
       \int k_j(x-y)\|G(y)\|_{\rm op}q(y)\,dy.
\]
Scalar Cauchy gives $B_{\rm high}^2\le V$.
Directional positive-weight optimization gives $(\int\sqrt{qr})^2$,
with $r=v^TK_jGK_j^Tv$. Zero-safe weights
$\sqrt{(r+\varepsilon q)/(q+\varepsilon r)}$ prove the infimum.
For the raw weight $k_j\|G\|_{\rm op}$ at its zero set, arbitrary positive
assignment can add an unwanted first-factor cost; adding $\varepsilon h$
with positive integrable $h$ and then taking the limit completes that proof.
The scalar Cauchy estimate itself never needs division by this weight.

Define the Schwartz annular antiderivative by
\[
 \widehat J_{jk}(\xi)=\frac{\xi_k}{2\pi i|\xi|^2}
               \mathbb P(\xi)\phi_j(\xi),\qquad
 \sum_k\partial_kJ_{jk}=K_j.
\]
For $n=Hm$ and $d_k=m_k-(\partial_kG)n=G\partial_kn$, the exact readout is
\[
 u_j=\sum_kJ_{jk}*d_k+\sum_kJ_{jk}*((\partial_kG)n).
\]
Let $d=\int\sum_kd_k^THd_k$, and define
\[
 \begin{aligned}
 W&=\sup_{j\ge0,x}\lambda_j^{-1}
   \left\|\int\sum_kJ_{jk}(x-y)G(y)J_{jk}(x-y)^T\,dy\right\|_{\rm op},\\
 C&=\sup_{j\ge0}\lambda_j^{-1/2}
       \left\|\sum_kJ_{jk}*((\partial_kG)n)\right\|_\infty.
 \end{aligned}
\]
Joint inverse-metric Cauchy gives
\[
 B_{\rm high}\le\sqrt{Wd}+C,\qquad
 B_{\rm high}^4\le2VWd+2VC^2.
\]
Supplied horizon-uniform $VW\le L_T$ and $\int VC^2\le J_T$ give
\[
 \int B_{\rm high}^4\le L_TE_0/\nu+2J_T.
\]
This gives the earlier full original conditional return.
Its independent source remains open. Here
$J_{jk}=\lambda_j^2J_{0k}(\lambda_j\cdot)$, so its squared kernel integral
scales as $\lambda_j$, repairing the raw mass-readout loss. The identity
heat ambient has $C=0$ and pays this calibration; it is a changed ambient,
since the actual specific metric $G=I$ would require $S=0$.

A local original curl split supplies another native component price:
\[
 V_r=(-\Delta)^{-1}\operatorname{curl}(\psi\omega),\qquad
 H_r=U-V_r,
\]
with $\psi=1$ on $B_{2r}$ and supported in $B_{3r}$.
Fourier contraction gives $\|\nabla V_r\|_2\le\|\nabla U\|_2$,
and $H_r$ is harmonic inside $B_{2r}$. Its ball mean value gives
\[
 \|\nabla H_r\|_{L^\infty(B_r)}
 \le2\sqrt{\frac3{4\pi}}r^{-3/2}\|\nabla U\|_2.
\]
This pays harmonic strain independently of diverging scaled energy.
The exact localized vorticity balance retains local stretching, transport
and diffusion collars. Native record drift has an integrated fourth-power
price, which does not itself supply small pointwise coefficients in a
backward weighted theorem. Local core, collar, coefficient and rigidity
returns remain open. These intermediate laws do not establish
unconditional Clay(A).


\section{Actual-history heat and localized sources}

The reviewed record object satisfies exact local energy, not merely its
inequality. For $e=|U|^2/2$ and compact smooth $\chi$,
\[
 \frac{d}{ds}\int\chi e+\nu\int\chi|\nabla U|^2
 =\int e(\chi_s+\nu\Delta\chi)+\int(e+\Pi)U\cdot\nabla\chi.
\]
Local strong products, strong spacetime gradients and canonical pressure
pass this identity from the actual finite histories. Write $\theta\ge0$
for a smooth compact terminal test of mass $M$, $b\le0$, and
$\Phi(s)=h_{\nu(b-s)}*\theta$. The near Fourier pressure multiplier and
far kernel give $\|\Pi\|_{L^2(B_L)}\le C_pS_0^2L^{1/2}$.
Gaussian pressure-shell estimates then prove, with
$F(s)=\int(e+\Pi)U\cdot\nabla\Phi$,
\[
 |F(b-\tau)|\le C_FMS_0^3(\nu\tau)^{-5/4}
 \quad\text{when }\tau\ge r^2/\nu.
\]
The remote energy is at most $S_0^2M/(6\sqrt{\pi\nu(b-a)})$.
Exhausting space first and then the remote past gives
\[
 \int\theta e(b)+\nu\int_{-\infty}^b\!\int\Phi|\nabla U|^2
   =\int_{-\infty}^b F(s)\,ds.
\]
For finite first-record histories the initial boundary is bounded by
\[
 \frac{ME_0}{2(4\pi\nu(t_k+b/B_k^2))^{3/2}B_k^2}.
\]
This is $O(B_k^{-2})$ for every fixed $b\le0$; its denominator must retain
$b$ before taking the limit.

The \emph{same} backward test supplies a stronger signed tail. For every
$R>0$,
\[
 \begin{aligned}
 T_R&=\int_{-\infty}^{b-R}F(s)\,ds\\
 &=\int\Phi(b-R)e(b-R)
   +\nu\int_{-\infty}^{b-R}\!\int\Phi|\nabla U|^2,\\
 0&\le T_R\le\frac{C_sM}{\sqrt{\nu R}},\qquad
 C_s=\frac{S_0^2}{6\sqrt\pi}+\frac1{4\pi^{3/2}}.
 \end{aligned}
\]
The absolute tail remains $O(R^{-1/4})$. The signed estimate does not
assert pointwise flux positivity. A supplied positive local energy forces
recent signed influx; no independent law excluding that influx is supplied.

The local curl split also pays a small quadratic stretch form. With
$a=\|\psi\omega\|_2\le1$, $f=\chi\omega$ and
$\beta=(4\pi/3)^{1/6}$,
\[
 \left|\int f_if_j\partial_jV_i\right|
 \le S_0^2\beta r^{1/2}a\|\nabla f\|_2^2.
\]
Thus $r\le\nu^2/(4S_0^4\beta^2)$ absorbs this term. Normalized cutoffs
$0\le\chi\le1$ retain both collars. Their transport price on
$I\subset[-R,0]$ uses $\int_I\|U\|_\infty\le
c_W|I|^{3/4}(KR)^{1/4}$, not an unproved $K|I|$ budget.
For reverse time $L_\nu=\partial_\tau+\nu\Delta$ and $w=\chi v$,
\[
 L_\nu w=b\cdot\nabla w-Jw+R_\chi,\qquad
 R_\chi=(\nu\Delta\chi-b\cdot\nabla\chi)v
              +2\nu\nabla\chi\cdot\nabla v.
\]
Native bounds pay absolute collars for each fixed weight. The quadratic
stretch estimate does not pay $|Jw|^2$ in a backward weighted argument;
seed, parameter-uniform collars and exact rigidity remain open.

\section{Rotational source and event calibration}

For the specific original metric put
\[
 n=Hm,\qquad d_k=G\partial_kn,\qquad
 \rho=\sum_k E_k(\partial_kG)n,
\]
where $E_kv=e_k\times v$. Then
$\omega=\sum_kE_kd_k+\rho$ and, with $A_{ij}=\partial_i u_j$,
\[
 \begin{aligned}
 L\rho={}&A^T\rho-\sum_kE_kG(\partial_kA)^Tn\\
 &+2\nu\sum_{k,l}E_k\big((\partial_kG)H(\partial_lG)
                         -\partial_l\partial_kG\big)\partial_ln.
 \end{aligned}
\]
The drift commutator $-\sum_l A_{kl}\partial_lG$ is retained in its
underlying differentiated metric law. Actual compact curl data can create
$\rho$ from zero initial connection. No invariant zero source is assumed.
For the matching annular inverse-curl kernel $R_j$ and $F_{jk}=R_jE_k$,
\[
 B_{\rm high}^2\le2W_{\rm rot}d+2C_{\rm rot}^2,
 \qquad B_{\rm high}^2\le c_R Z.
\]
Writing $V_* =\min(V,c_RZ)$ gives
$B_{\rm high}^4\le2V_*(W_{\rm rot}d+C_{\rm rot}^2)$.
Supplied uniform $V_*W_{\rm rot}\le L_T$ and
$\int V_*C_{\rm rot}^2\le J_T$ return
$\int B_{\rm high}^4\le L_TE_0/\nu+2J_T$ on the same path.
Their independent source remains open.

An exact scalar source identity records the sign explicitly. Set
$D=\det G$, $b=\log D$, $e=G\omega$,
$C_k=G^{-1/2}(\partial_kG)G^{-1/2}$ and $v=G^{1/2}\omega$. Then
\[
 L\left(\frac{\omega^TG\omega}{D}\right)
 =-\frac{2\nu}{D}\sum_k|\partial_ke-(\partial_kb)e|_H^2
   +\frac{2\nu}{D}\sum_kv^T\operatorname{adj}(C_k)v.
\]
Arbitrary SPD metric jets allow either cofactor sign. This observation alone
makes no assertion about their attainment by the specific original metric.

For actual first hits $B_k=2h_k$ and last preceding half-level crossings,
let $\ell_k=B_k^2(t_k-a_k)$. A bounded subsequence of $\ell_k$ supplies
$\int_{-\ell_k}^0P\ge1/4$. The anchored $Q$ budget removes bad times and
the cube law produces a nonzero ancient core. This uses the variable event;
it does not imply fixed-$R$ half-growth, since earlier excursions can occur.
If $\ell_k\to\infty$, the event instead supplies only
\[
 \frac1{\ell_k}\int_{-\ell_k}^0(P-\nu Q)=\frac1{4\ell_k}\to0.
\]
This signed near-balance proves neither local concentration nor stationarity.

A direct stationary calibration is available under an extra premise:
a smooth stationary original field with $\nabla W\in L^2$, $W\in L^6$,
canonical pressure and $W\in L^{9/2}$ must vanish. For annular norm
$n_R=\|W\|_{L^{9/2}(B_{2R}\setminus B_R)}\to0$, cutoff diffusion costs
$CR^{-1/3}n_R^2$ and cubic flux $Cn_R^3$. Near pressure has $L^2$ price
$CR^{1/6}\|W\|_{9/2}^2$ and far pressure supremum price
$CR^{-4/3}\|W\|_{9/2}^2$; both fluxes cost
$C\|W\|_{9/2}^2n_R\to0$. Exact cutoff energy therefore gives
$\nabla W=0$ and $W=0$. The native $L^6$ estimates instead lose $R^{1/2}$.
The extra integrability and any actual stationary-record adapter remain
open. These reviewed intermediate laws do not establish unconditional Clay(A).


\section{Two ancient adapters and actual scalar sign controls}

A hypothetical finite endpoint forces unbounded velocity amplitude, since
\[
 Z'+\nu Q\le\nu^{-1}\|u\|_\infty^2Z
\]
would otherwise bound $Z$ and return original continuation. At first hits
$M_k$ of amplitude, attained maxima $x_k$ give
\[
 V_k(y,s)=M_k^{-1}u(x_k+y/M_k,t_k+s/M_k^2),\qquad
 |V_k|\le1,\quad |V_k(0,0)|=1.
\]
The past domains exhaust $(-\infty,0]$. The exact heat--Leray divergence
kernel is
\[
 K_{ijk}(t)=\partial_k\left(\delta_{ij}h_{\nu t}
       +\partial_i\partial_j\int_{\nu t}^\infty h_r\,dr\right).
\]
Gaussian derivative integration gives $|D^mK(t,x)|\le
C_m(|x|^2+\nu t)^{-(4+m)/2}$ and
$\|D^mK(t)\|_1\le C_m(\nu t)^{-(m+1)/2}$.
Spatial translations split at $\tau=|z|^2/\nu$ give a logarithmic
Lipschitz modulus; heat restart gives time continuity. Local uniform
compactness includes the terminal face. Kernel spatial tails $O(1/R)$
and the final-slice price $O(\sqrt{\varepsilon/\nu})$ pass every finite
mild identity. Thus the limit $V$ is a nonzero bounded ancient mild field
with $|V(0,0)|=1$. Its scaled energy may diverge and nonzero constants
remain admissible.

The all-order classical adapter uses the explicit repaired derivative
construction. At order $m$, place all derivatives on the kernel before
$t/2$, and all derivatives on the product after $t/2$. For
$A_m=\sup(\nu t)^{m/2}\|D^mv(t)\|_\infty$,
\[
 A_m\le\|D^mh_1\|_1A_0+C_m\sqrt{h/\nu}
 \left(A_0^2+A_0A_m+\sum_{1\le j<m}A_jA_{m-j}\right).
\]
Each finite order has its own absorption length. Regularized Picard iterates
and finite differences construct derivatives before passing the singular
pressure operator. Near pressure retains its local delta correction and
Taylor subtraction; far derivative kernels have integrable $r^{-4}$ tails
and a positive unit-sphere boundary term. This passes the actual canonical
pressure gradients and all mixed derivatives. The coherent gauge fixes
$p(0,s)=0$. These ancient terminal traces are distinct from the original
physical right initial-face return.

The separate enstrophy-record class has global $L^6$ bound $M=S_0$.
Its actual mild identity passes using weak $L^6$ heat tails, tensor $L^3$
spatial tails and a final-strip $L^6$ price $O(\varepsilon^{1/4})$.
In $X=\sup\|v\|_6+\sup(\nu t)^{1/4}\|v\|_\infty$, put
\[
 \begin{aligned}
 C_h&=\|h_1\|_{6/5},\qquad H_0=1+C_h,\\
 D_0&=\|K_{\text{heat }1}\|_{6/5}
                  \bigl(4+\mathrm B(1/4,3/4)\bigr),\\
 \delta&=\nu^3/(16D_0H_0M)^4.
 \end{aligned}
\]
The radius $2H_0M$ ball contracts. Uniqueness against the actual solution
of $L^6$ size $M$ uses both radii, with coefficient at most
$(1+2H_0)/(16H_0)<1$. Every ancient restart then gives
$\|U\|_\infty\le32D_0H_0^2M^2/\nu$.
Zero is handled separately. Nonzero constants are excluded by $L^6$,
but the actual extracted field can still vanish locally.

This class has a paid signed projected-source account. Define
\[
 N_R(b)=-\int_{b-R}^b K(b-s)*(U\otimes U)(s)\,ds.
\]
Exact mildness gives $U(b)-N_R(b)=h_{\nu R}*U(b-R)$, so
\[
 \|U(b)-N_R(b)\|_\infty\le C_hM(\nu R)^{-1/4}.
\]
The full signed remote source converges uniformly; absolute convergence
is not established. Differentiated remainders use the heat identity,
without integrating a differentiated kernel at its singular tip.
For pointwise records, the distinct actual remote contribution is
\[
 \begin{aligned}
 A_{k,R}&=-\int_{a_k}^{-R}[K(-s)*(V_k\otimes V_k)(s)](0)\,ds,\\
 \lim_k A_{k,R}&=[h_{\nu R}*V(-R)](0).
 \end{aligned}
\]
The original initial heat term vanishes as $M_k^{-1}$. If $V=c$, $|c|=1$,
its remote nonlinear memory is $c$ for every fixed $R$. No uniform native
absolute bound excludes this memory. The two record classes cannot exchange
these source conclusions without an independent bridge.

Finally, the cofactor sign obstruction can be realized on actual compact
original data. Let
\[
 \phi=-ax^3/3+ax(y^2+z^2)/2,\quad u_2=\nabla\phi,
 \quad u_1=(w\times X)/2.
\]
Write $S_k=\partial_kD^2\phi$ and $M_{ik}=(S_kw)_i$. This $M$ is symmetric
tracefree. Set
\[
 u_4=-|X|^2(X\times MX)/(42\nu),\qquad
 \nu\Delta\operatorname{curl}u_4=MX.
\]
For each homogeneous solenoidal $u_d$,
$\operatorname{curl}[-X\times u_d/(d+2)]=u_d$.
A compact cutoff of these potentials preserves all origin jets and gives
legal Schwartz data. On its actual local solution, $G_{k,t}=-2S_k$ and
the quartic diffusion makes $\omega_{k,t}=2S_kw$, so
$\partial_t\partial_k(G\omega)=0$ initially. The completed square is
$O(t^4)$, while the cofactor is quadratic. The full operator has signs
\[
 \begin{aligned}
 w=b_0e_1:&\quad L(r/D)(0,t)=8\nu a^2b_0^2t^2+O(t^3),\\
 w=b_0e_2:&\quad L(r/D)(0,t)=-24\nu a^2b_0^2t^2+O(t^3).
 \end{aligned}
\]
This refutes a universal pointwise subsolution law for this particular scalar.
It supplies no singular solution or endpoint source price.
For the restricted harmonic quadratic subclass with locally constant initial
vorticity, put $E_T=\sum_{ijk}T_{ijk}^2$ and
$N_T=\sum_k|T_kw|^2$ for its symmetric tracefree third tensor.
The sharp law is $N_T\le3|w|^2E_T/5$. Its full leading operator is
$2\nu(3N_T-2|w|^2E_T)t^2+O(t^3)$ and is negative for nonzero $T,w$.
The quartic operation exits precisely that neighborhood hypothesis.
Actual core, class-sensitive rigidity, constant exclusion and the mixed
rotational source remain independent open obligations. Unconditional
Clay(A) is not established.

\section{Dissipation cores and canonical head-tail carriers}

Round59 constructs five laws, checked by three independent full reviews
and a distinct parent pass. The actual source suppliers remain open.
For the original first enstrophy records of round56, write
$D_{k,R}=\int_{-R}^0Q_{U_k}(s)\,ds$. Then
\[
 \begin{aligned}
 &\text{some translated subsequence has a nonzero ancient limit}\\
 &\qquad\Longleftrightarrow
 \text{some fixed }R>0\text{ has }\limsup_kD_{k,R}>0.
 \end{aligned}
\]
Here the existential choice of centers is essential. At a record,
$\int P\ge\nu D_{k,R}$. If $D_{k,R}\ge c>0$, choose
$L=\max\{1,(2C_0KR/(\nu c))^4\}$ with $D_{k,R}\le KR$.
Production on $Q>L$ costs at most $\nu c/2$; the remaining signed
production, the original cube inequality and local time compactness
give a positive recentered velocity witness. Conversely, if all fixed
$D_{k,R}$ vanish, every translated limit is spatially harmonic in
distributions and has global $L^6$ norm bounded. The ball mean-value bound
$|U(x,s)|\le C\rho^{-1/2}\|U(s)\|_6$ forces zero, and local time
continuity includes the terminal trace. A nonzero limit therefore has a
nonzero compact Laplacian pairing, whose Cauchy bound gives the reverse
subsequential dissipation lower bound.

Starvation also pays an exact variation estimate:
\[
 \sup_{-R\le s\le0}|1-Z_{U_k}(s)|
 \le2C_0R^{1/4}D_{k,R}^{3/4}+2\nu D_{k,R}=:\varepsilon_k.
\]
Only after this lower bound on $Z$ may one divide to obtain
\[
 \frac{\nu^4D_{k,R}}R\le
 \mathop{\mathrm{avg}}_{[-R,0]}\frac{\nu^4Q}{Z^3}
 \le\frac{\nu^4D_{k,R}}{R(1-\varepsilon_k)^3}.
\]
Taylor length divided by enstrophy length then diverges in measure,
with no uniform time claim. In the viscosity-preserving normalization,
the physical dissipation on this same window is $B_kD_{k,R}$.
None of these facts supplies fixed-depth nonstarvation or ancient rigidity.

For any smooth original unforced field, $e=|u|^2/2$ and $H=e+p$ satisfy
\[
 \nabla H=u\times\omega-\nu\operatorname{curl}\omega-u_t,
 \qquad (\nu\Delta-u\cdot\nabla)H=\nu|\omega|^2+e_t.
\]
The last term has no sign. For a stationary field $W\in L^6$,
$\nabla W\in L^2$, with the bounded spatial derivatives supplied by the
qualified ancient adapter, the canonical pressure tends to zero at infinity.
One proof uses the longitudinal $L^2$ projection of $(W\cdot\nabla)W$,
the homogeneous Sobolev representative, and vanishing canonical ball means
to fix its gauge. Bounded pressure gradient then gives uniform continuity;
uniform continuity and $L^6$ give decay. A compact test
$(H-\epsilon)_+$ of the stationary identity yields
\[
 -\nu\int_{H>\epsilon}|\nabla H|^2
 =\nu\int(H-\epsilon)_+|\omega|^2\ge0.
\]
Thus $H\le0$ in this stationary class. This sign is not available for
arbitrary time-dependent ancient fields or their averaged heads.

A radial pressure carrier gives a weaker conditional return than finite
weighted energy. Suppose a smooth solenoidal $W\in L^6$ has its canonical
pressure $\Pi$ and $H=|W|^2/2+\Pi\le0$. Define
\[
 \begin{aligned}
 E(R)&=\int_{B_R}|W|^2,&
 T(R)&=\int_{|x|>R}\frac{|W|^2}{|x|^3},\\
 T_{\rm rad}(R)&=\int_{|x|>R}\frac{|W\cdot x/|x||^2}{|x|^3},&
 I(R)&=\int_R^\infty\frac{E(s)}{s^4}\,ds.
 \end{aligned}
\]
All these quantities are finite for $R>0$, since $E(R)=O(R^2)$ and the
tail weight belongs to $L^{3/2}$. The Newton potential of $1_{B_R}$ is
$R^2/2-r^2/6$ inside and $R^3/(3r)$ outside. Its value and first derivative
match at the sphere, so its Hessian has no surface delta. The interior
Hessian is $-I_3/3$, retaining the local pressure multiplier correction.
The exact canonical duality is
\[
 \int_{B_R}H=\frac{E(R)}6+\frac{R^3}3(3T_{\rm rad}(R)-T(R)).
\]
To justify it, first truncate $W\otimes W$ to an $L^2$ tensor and use
selfadjointness of the full Fourier multiplier. On the ball, the omitted
far pressure is bounded by $CM^{-1}\|W\otimes W\|_3$. The test Hessian is
in $L^{3/2}$, which passes tensor exhaustion. Sharp sphere matching or
smooth radial approximation pays the ball boundary. No global integral
of pressure is assumed.

Integration by parts gives $T=-E/R^3+3I$, because $E(s)/s^3\to0$.
Head sign, with the radial term retained, therefore gives
\[
 \begin{aligned}
 E(R)&\le2R^3(I(R)-T_{\rm rad}(R)),\qquad F(R)=R^2I(R),\\
 F'(R)&=2RI(R)-E(R)/R^2\ge2RT_{\rm rad}(R)\ge0.
 \end{aligned}
\]
If $W$ is nonzero, $F$ is positive at some radius and stays positive.
The extra premise $E(R)=o(R)$ instead gives $F(R)\to0$, hence forces
$W=0$. Native $L^6$ bounds supply only $O(R^2)$ and do not pay this source.
Finiteness of $J=\int|W|^2(1+|x|^2)^{-1/2}$ implies $o(R)$ by dominated
convergence. The converse fails in the norm ambient: disjoint solenoidal
annular bumps of radius $2^n$ and amplitude $2^{-n}/\sqrt n$ have
$E(R)=O(R/\log R)$ but a divergent harmonic sum for $J$. They are norm
controls, not stationary solutions or actual record limits.

An independent weighted proof uses $\theta=(1+r^2)^{-1/2}$ and
$A(r)=\int_0^r s^2\theta(s)\,ds$. The limiting dual matrix has eigenvalues
$\theta/2-A/r^3>0$ and $2A/r^3-\theta/2\ge\theta/6$.
Compact weight cutoffs have uniformly bounded Hessians by $2\theta$.
Under extra $J<\infty$, dominated convergence of their quadratic pairings
gives a nonnegative limiting form, while every head pairing is nonpositive.
Thus $W=0$, without a finite uncut Newton potential or an independently
absolute weighted pressure integral. Actual stationary attainment, head
applicability and either spatial source remain open.

Finally, long-event time occupation on the middle half interval has shift
error $4|h|\|f\|_\infty/\ell_k$. Probability weak convergence with continuous
shifts would yield invariance; existence, original-NS support and nonzero
local mass remain separate adapters. On any actual normalized interval of
length $L$, integrated enstrophy and Young give
$\int Q\le64L/\nu^4+1/\nu$; shifted endpoints are not treated as records.
For a supplied invariant Borel probability on smooth eternal original
paths with legal local energy identities, the extra moment
$\mathbb E\|U(0)\|_{9/2}^3<\infty$ gives a conditional zero law.
Invariance cancels expected energy endpoints. The original shell prices
$C\nu R^{-1/3}n_R^2+Cn_R^3+CN^2n_R$ vanish in expectation by dominated
convergence. Fatou yields zero gradient; $L^6$ excludes constants, and
rational-time continuity gives zero entire paths almost surely. Correlations
remain inside $\mathbb E(HU)$; individual stationarity is never inferred.
The moment and positive local mean are still independent original-source
obligations. Unconditional Clay(A) remains open.

\section{Source-sensitive force and rotational controls}
\label{sec:source-controls60}

Round60 adds reviewed component laws, with the original all-data unforced
target still open. The changed-equation force control uses only an
explicitly supplied primary statement; its substantive proof and
formalization have not been audited here. No forced conclusion is imported
as an unconditional theorem about the original target.

For a smooth forced path with legal whole-space energy integrations, put
$G=-\langle\Delta u,f\rangle$. Direct momentum and divergence give
\[
 \begin{aligned}
 E'&=-2\nu Z+2\langle u,f\rangle,&
 Z'&=2P-2\nu Q+2G,\\
 -\Delta p&=\partial_i\partial_j(u_i u_j)-\operatorname{div}f.
 \end{aligned}
\]
The force correction to canonical pressure is
$-(-\Delta)^{-1}\operatorname{div}f$. At genuine diverging enstrophy
records, the repository normalization $U=B^{-1}u$, $F=B^{-3}f$ retains
viscosity and gives
\[
 \begin{gathered}
 \|D_y^m F\|_2=B^{-m-3/2}\|D_x^m f\|_2,\\
 E_U=BE,\qquad Z_U=B^{-1}Z,\qquad Q_U=B^{-3}Q.
 \end{gathered}
\]
Assume the required physical force norms are uniformly bounded. Cubic
Young and force Young leave $\nu Q_U/4$ after absorption. On a fixed
anchored record window of length $R$,
\[
 \int Q_U\le\frac{128R}{\nu^4}
       +\frac{4M_f^2R}{\nu^2B^3},\qquad
 \left|\int G_U\right|
       \le B^{-3/2}M_f\left(R\int Q_U\right)^{1/2}.
\]
An arbitrary shifted subinterval adds $2/\nu$ to the dissipation bound.
No shifted endpoint is treated as a record. The coefficient128 is coarse;
it does not replace the sharper unforced coefficient. Heat-Leray has the
direct Fourier bound $L^2\to L^\infty$ with norm
$C(\nu\tau)^{-3/4}$. Hence on a fixed time pair
\[
 \left\|\int_a^b e^{\nu(b-s)\Delta}\mathbb P F(s)\,ds\right\|_\infty
 \le4C\nu^{-3/4}(b-a)^{1/4}B^{-3/2}M_f\longrightarrow0.
\]
These estimates pay the direct finite-window force passage. Nonlinear
compactness, nonzero extraction, limiting pressure class, entire history
and original continuation remain separate obligations.

The type of work matters: on matched physical and normalized intervals,
\[
 \int G_U\,ds=B^{-1}\int G\,dt,
 \qquad \int\langle U,F\rangle\,ds=B\int\langle u,f\rangle\,dt.
\]
Vanishing scaled enstrophy work does not erase kinetic injection. If a
global smooth energy-class path has already been returned, define
$M=\int_0^\infty\|f(t)\|_2\,dt$ and assume it is finite. The regularized
square-root energy calculation gives
\[
 \frac{d}{dt}\sqrt{E+\epsilon^2}\le\|f(t)\|_2,
 \qquad \sup_{t\ge0}E(t)\le(\sqrt{E(0)}+M)^2.
\]
This includes zero energy and supplies a conditional uniform bound.
It constructs no global path. A contradiction with a supplied forced
statement would require separately proved adapters for every substantive
step, the same force and datum, source and inherited class, complete
history, coherent pressure and global smooth energy-class return.
Local force disappearance alone proves no complete transfer or blanket
rejection of an unforced strategy.

For the genuine vector rotation mean $v$, let $w=u-v$,
$R_{ij}=\langle w_iw_j\rangle$ in rotating cylindrical components, and
$\bar\Gamma=rv_\theta$. Direct angular momentum gives
\[
 \begin{aligned}
 (\partial_t+v_r\partial_r+v_z\partial_z)\bar\Gamma
 &=\nu(\partial_{rr}-r^{-1}\partial_r+\partial_{zz})\bar\Gamma+T
       +r\langle f_\theta\rangle,\\
 T&=-r^{-1}\partial_r(r^2R_{r\theta})-r\partial_zR_{z\theta}.
 \end{aligned}
\]
Angular pressure derivatives vanish by periodicity, without discarding
the other pressure components. For compact tests away from the axis,
physical measure $2\pi r\,dr\,dz$ gives
\[
 \int\psi T=\int r(R_{r\theta}\psi_r+R_{z\theta}\psi_z),
 \qquad |\int\psi T|\le\tfrac12\int r|\nabla\psi|\langle|w|^2\rangle.
\]
Under its actual source, growth and axis hypotheses, the circulation
maximum is charged by $\int\|T+r\langle f_\theta\rangle\|_\infty$.
The weak price is not this maximum-principle norm. Zero external force
does not remove Reynolds torque; mean circulation alone supplies neither
swirl shear nor the full fluctuation continuation criterion.

An exact nonlinear affine calibration retains viscous seed history.
For $\operatorname{tr}A=0$, $A'+A^2$ symmetric, $k'=-A^Tk\ne0$, and
$a\cdot k=0$, set
\[
 \begin{aligned}
 a'&=-Aa+2k(k\cdot Aa)/|k|^2-\nu|k|^2a,\\
 u&=Ax+a\cos(k\cdot x),\\
 p&=-\tfrac12x\cdot(A'+A^2)x
           -2(k\cdot Aa)\sin(k\cdot x)/|k|^2.
 \end{aligned}
\]
The full unforced equation holds: phase transport cancels and the wave
has zero self-advection. Its amplitude satisfies
\[
 (|a|^2)'=-2a\cdot Sa-2\nu|k|^2|a|^2.
\]
Reversed diagonal strain can amplify despite viscosity. These backgrounds
and waves generally have infinite energy. Localization, pressure residuals,
accumulated nonlinear error and original initial-seed comparison remain open.

Finally $f(x,t)=e^{-1/t^2}e_1$ before zero, extended by zero, is spatially
analytic on every slice and smooth with every endpoint jet zero, yet is
nonzero at every earlier time. Its regular local solution
$u=e_1\int_{-1}^t e^{-1/s^2}\,ds$, $p=0$ has infinite whole-space energy.
It refutes only the inference from spatial analyticity and endpoint time
flatness to locally zero force. A separate direct identity law propagates
exact symmetry from a nonempty open core for an actually analytic field
on a connected invariant ball. Actual analytic/core production and the
substantive endpoint regularity return remain unpaid. Unconditional
Clay(A) remains open.

\section{Native source charges and the anchored starvation alternative}
\label{sec:native-anchor61}

Round61 is independently reviewed and integrated at the following scopes.
Schwartz initial data supply $Y_0=\int|x|^2|u_0|^2$, while native energy
supplies $E\le E_0$ and $\int Z\le E_0/(2\nu)$. Evolved moment finiteness
is derived using bounded increasing weights
$w_L=|x|^2/(1+|x|^2/L^2)$, rather than assumed. Their gradients satisfy
$|\nabla w_L|\le2\sqrt{w_L}$ and $\Delta w_L\le6$.
Canonical pressure has the direct tensor $L^2$ estimate
$\|p\|_2\le\|u\|_4^2$, retaining the pressure energy flux. Thus
\[
 Y_L'\le6\nu E_0+6S_0^{3/2}E_0^{1/4}Z^{3/4}\sqrt{Y_L}.
\]
For every $m>0$, regularization by $\sqrt{Y_L+m^2}$ and native dissipation
give, at every existing $t\le T$,
\[
 \begin{aligned}
 \sqrt{Y(t)}\le C_Y(T,m)&=\sqrt{Y_0+m^2}+3\nu E_0T/m\\
 &\quad+3S_0^{3/2}E_0^{1/4}T^{1/4}(E_0/(2\nu))^{3/4}.
 \end{aligned}
\]
Monotone weight exhaustion constructs $Y(t)$ without a terminal higher
norm or evolved Schwartz assumption. This physical moment scales by $B^3$
under a record zoom and is not a uniform moment of an ancient limit.

Let $U=\mathcal A u$ be vector rotation projection, $v=u-U$ and
$R=\mathcal A(v\otimes v)$. Solenoidality before absolute values gives
the exact unprojected mean force
\[
 F_R=-\operatorname{div}R=-\mathcal A[(v\cdot\nabla)v],
 \qquad U_t+(U\cdot\nabla)U=\nu\Delta U-\nabla(\mathcal A p)+F_R.
\]
The canonical mean pressure depends on $U\otimes U+R$. Rotation
orthogonality, Jensen and spatial Cauchy give the actual source charges
\[
 \int_0^t\|F_R\|_1\le E_0\sqrt{T/(2\nu)},
 \qquad\int_0^t\|rF_R\|_1\le C_Y(T,m)\sqrt{TE_0/(2\nu)}.
\]
The torque is $r(F_R)_\theta$. Convex Kato in physical measure
$2\pi r\,dr\,dz$ has diffusion adjoint
$\partial_{rr}+3r^{-1}\partial_r+\partial_{zz}$.
Actual circulation $\Gamma=rU_\theta=O(r^2)$ removes its axis flux.
An exponential exhaustion with radius at least1, followed by zero-safe
regularizer exhaustion, returns
\[
 \|\Gamma(t)\|_{L^1(dx)}\le\|\Gamma(0)\|_{L^1(dx)}
                         +C_Y(T,m)\sqrt{TE_0/(2\nu)}.
\]
Compact drift bounds justify passage and disappear from the final constant.
This supplies mean circulation mass, not swirl shear or fluctuation return.
No bare Leray $L^1$ bound is claimed. With angular Fourier frequency,
the projected source instead satisfies
\[
 \int_0^t\|\mathbb P F_R\|_{\dot H^{-1}}^{4/3}
 \le\frac{S_0^2E_0^{4/3}}{2\nu},
\]
and curl force in $\dot H^{-2}$ has the same price. These are homogeneous
$L^2$ multiplier budgets, with no pointwise division at zero frequency.

For an actual forced energy-class control, supplied global norms
$M_1=\int_0^{T_0}\|f\|_2$, $M_{\nabla}^2=\int_0^{T_0}\|\nabla f\|_2^2$
and $M_0=\sup\|f\|_2$ give
\[
 \begin{aligned}
 E_{\max}&=(\sqrt{E_0}+M_1)^2,\\
 D_Z&=(E_0/2+\sqrt{E_{\max}}M_1)/\nu,\\
 \int_0^t|G|&\le\sqrt{D_ZM_{\nabla}^2}.
 \end{aligned}
\]
Here $G=\langle\nabla u,\nabla f\rangle$ by actual spatial integration.
Complete normalized history enstrophy work is therefore $O(B^{-1})$.
The entire explicit linear heat-force response, uniformly over its prefixes,
is at most $4C_H\nu^{-3/4}M_0T_0^{1/4}B^{-1}$,
$C_H=(8\pi)^{-3/4}$. Kinetic work still scales by $B$; neither result
estimates nonlinear forced/unforced amplification or verifies a supplied
forced-singularity theorem.

The actual record source has another paid spatial price. Write raw
$z=\|\nabla U_k\|_2^2\le1$, $q=\|\Delta U_k\|_2^2$ for the FULL normalized
velocity, distinct from rotation mean $U$ above. Cartesian Sobolev and
interpolation give
\[
 \|F_R\|_2\le S_0^{3/2}z^{3/4}q^{1/4},
 \qquad\int_J\|\mathbb P F_R\|_2^4\le S_0^6\int_Jq.
\]
The same product bound holds for full convection. A generic
$L^4_tL^2_x$ norm does not bound terminal heat response in $L^\infty$.
An explicit solenoidal Gaussian curl, scaled over logarithmic time shells,
has uniformly bounded source norm and response growing as $\sqrt N$.
This is a smooth linear force control, with no actual Reynolds realization.

The actual record has additional structure omitted from that generic test:
EVERY backward length ending at the same record satisfies
\[
 A(H)=\int_{-H}^0q\le cH,
 \qquad c=K/\nu^4,\quad K=1024/27.
\]
Thus its anchored quarter-root shell sum is already paid. More directly,
$M(H)=\int_0^Hq(-\tau)^{1/4}\,d\tau\le c^{1/4}H$ and integration
against the decreasing heat weight give
\[
 \int_0^R\tau^{-3/4}q(-\tau)^{1/4}\,d\tau
 \le4c^{1/4}R^{1/4}.
\]
The lower boundary is $O(\delta^{1/4})$ and vanishes. If
$D=A(R)>0$, split using $M(H)\le H^{3/4}\min(cH,D)^{1/4}$ at $D/c$:
\[
 \int_0^R\tau^{-3/4}q(-\tau)^{1/4}\,d\tau
 \le D^{1/4}\bigl[4+\tfrac34\log(cR/D)\bigr].
\]
At $D=0$ the integral is zero before any logarithm. This corrects the
source diagnosis while preserving the generic endpoint countercontrol.
Arbitrary earlier endpoints may have an enstrophy deficit and do not
inherit the same nested law.

Suppose an exhausting actual record sequence starves at EVERY fixed depth:
$D_k(R)\to0$ for each fixed $R$. Its exact finite-pair mild momentum has
full nonlinear response tending to zero globally by the logarithmic bound.
The earlier slice has $L^6$ norm at most $S_0$, so heat gives
$\limsup_k\|U_k(0)\|_\infty\le C_6S_0(\nu R)^{-1/4}$.
Only then let $R\to\infty$, obtaining global terminal velocity vanishing.
At a positive smooth record the enstrophy derivative is nonnegative;
the original pairing gives
\[
 q_k(0)\le\|U_k(0)\|_\infty^2/\nu^2\longrightarrow0,
 \qquad z_k(0)=1.
\]
For every fixed angular normalized cutoff $\kappa_0>0$, gradient energy
above it is at most $q_k(0)/\kappa_0^2$. Enstrophy therefore escapes
toward zero normalized frequency. Physical frequency is $B_k$ times that
frequency; no physical fixed-cutoff conclusion follows. Scaled energy
allows $q_k(0)\ge1/(B_kE_0)$, consistently tending to zero. Physical
velocity is $B_k\|U_k(0)\|_\infty$, with no rate supplied. This identifies
the actual starved alternative more precisely and excludes none of its
possible occurrences. Independent source/scale exclusion or another
original continuation return remains open. Unconditional Clay(A) is open.

\section{Native analytic and weighted sources, with exact symmetry provenance}
\label{sec:native-analytic-symmetry}

All conclusions in this section concern the original existing strong path.
They establish scoped laws and conditional returns; unconditional Clay(A)
remains open. The five round62 bodies retain the earlier displays and apply
an explicit divergent-history and fixed-frame qualification.

Write $k=2\pi\xi$, $L_0=e^{1/4}$,
$c_\gamma=2^{\max(\gamma-1,0)}$, and at an existing restart $a$ set
\[
 \lambda(t)=\sqrt{\nu(t-a)},\quad
 X_\lambda(u)=\int e^{\lambda|k|}|\widehat u(\xi)|\,d\xi,
 \qquad H(u)=\|u\|_{H^\gamma}.
\]
For $\tau=t-s$, the moving weight has increment at most
$\sqrt{\nu\tau}$, and the actual heat multipliers obey
\[
 \sup_{r\ge0}e^{-r^2+r}=L_0,
 \qquad \sup_{r\ge0}r e^{-r^2+r}=1.
\]
Tensor convolution and the full original heat--Leray divergence therefore
have price $(\nu\tau)^{-1/2}$. With $W_0=\|\widehat u(a)\|_1$,
$H_0=\|u(a)\|_{H^\gamma}$, $R=2L_0W_0$ and $S=2H_0$, the nonlinear map
on $X\le R,H\le S$ satisfies
\[
 X_{\rm new}\le L_0W_0+2\sqrt{h/\nu}\,R^2,
 \qquad H_{\rm new}\le H_0+4c_\gamma\sqrt{h/\nu}\,RS.
\]
Its scaled difference norm $\max(\Delta X/R,\Delta H/S)$ has factors
$4\sqrt{h/\nu}R$ and $8c_\gamma\sqrt{h/\nu}R$. Thus
\[
 h\le\frac{\nu}{256c_\gamma^2R^2}
\]
gives contraction at most $1/2$ and invariant image bounds $5R/8,3S/4$.
The zero state is treated by uniqueness before division. Weighted Fourier
continuity follows from dominated initial heat and a Volterra tip of size
$O(\sqrt\delta)$, with continuous dominated multipliers away from the tip.
Weighted $L^1$ and $H^\gamma$ limits identify the same Fourier distribution
on compact frequency tests, proving completeness. Original overlap uniqueness
identifies this genuine mild solution with the SAME maximal path, including
any lawful local extension. The possibly large finite $S$ cancels from the
time constraint; it has not become a uniform endpoint high norm.

Fourier inversion gives a spatial holomorphic tube $|\operatorname{Im}x|<\lambda$
with amplitude $X$ and
\[
 \|D^\alpha u(t)\|_\infty
 \le |\alpha|!\lambda(t)^{-|\alpha|}X_{\lambda(t)}(u(t)).
\]
For a positive compact $J=[t_{\min},t_{\max}]$, first bound $W$ on the
EARLIER existing compact $K=[t_{\min}/2,t_{\max}]$, then choose
$h\le t_{\min}$ and restart at $t-h/2$. This pays common spatial radius
and amplitude on $J$. It proves neither time analyticity nor physical
terminal-uniform analyticity.

At an actual normalized enstrophy record, $z\le1$ and every available
backward length obeys $\int_{-H}^0q\le cH$, $c=(1024/27)\nu^{-4}$.
The outer half contains an actual earlier $a$ with $q(a)\le2c$; native
Fourier interpolation gives $W(a)\le M=C_W(2c)^{1/4}$,
$C_W=\sqrt2/\pi$. Define
\[
 R_*=2L_0M,\quad h_*=\frac{\nu}{256c_\gamma^2R_*^2},
 \quad H=h_*/2,\quad \lambda_*=\sqrt{\nu H/2}.
\]
The SAME local map, with the finite actual $H^\gamma$ ball at $a$, reaches
record time0 and yields $X_{\lambda_*}(U_B(0))\le R_*$. This requires
$H\le B^2t_B$. It eventually holds for DIVERGENT records
$B_n\to\infty,t_n\to T>0$; approaching a regular time alone is insufficient.

The independent homogeneous analytic norm
$J_\lambda=\||k|e^{\lambda|k|}\widehat u\|_2$ satisfies
$J(a\otimes b)\le J(a)X(b)+X(a)J(b)$. The separate heat-divergence
factor gives $J_{\rm new}\le L_0J_0+\eta J_{\rm old}$,
$\eta=1/(4c_\gamma)\le1/4$, $J_0\le1$. Differences also contain the
summable $X$-difference source, so the Picard limit genuinely has
\[
 J_{\lambda_*}(U_B(0))\le2L_0,\qquad
 \int_{|k|>\kappa}|k|^2|\widehat U_B(\xi,0)|^2d\xi
       \le4L_0^2e^{-2\lambda_*\kappa}.
\]
This is NORMALIZED gradient high-frequency tightness, at physical cutoff
$B\kappa/(2\pi)$. It retains neither low-frequency nor spatial gradient mass.
If $D_B=\int_{-H}^0q_B>0$, the same window and same $h_*$ work with
$R_D=2L_0C_W(2D_B/H)^{1/4}\le R_*$. Therefore
\[
 \|D^\alpha U_B(0)\|_\infty
 \le R_D|\alpha|!\lambda_*^{-|\alpha|}.
\]
ONE fixed-window starvation $D_B\to0$ makes every fixed global spatial
derivative sup vanish. At $D_B=0$, continuity, $L^2$ harmonic zero and
forward uniqueness contradict a nonzero record before any division.
Global gradient $L^2$ may still equal1 on an expanding region.

The native weighted kinetic identity also constructs
\[
 \int_0^t\!\int |x|^2|\nabla u|^2\le D_Y,
 \qquad
 2\nu D_Y=Y_0+6\nu E_0T+
 6S_0^{3/2}E_0^{1/4}C_YT^{1/4}D_Z^{3/4}.
\]
Here $D_Z=E_0/(2\nu)$ and $C_Y$ is the previously constructed second-moment
price. Retain bounded-weight dissipation before nonnegative exhaustion;
canonical pressure is included. For $a(x)=\sqrt{\ell^2+|x|^2}$, bounded
$a_L=a/(1+a/L)$, and zero-safe curl $g_\epsilon=\sqrt{|\omega|^2+\epsilon^2}-\epsilon$,
the independent native magnitude mass
$M_*=\|\omega_0\|_1+E_0/(2\sqrt2\nu)$ pays exhaustion. Pointwise
$|S|\le|\nabla u|$, drift Cauchy and $\Delta a_L\le3/\ell$ give
\begin{align*}
 \int a|\omega(t)|+\nu\int_0^t\!\int a|\omega||\nabla n|^2
 &\le V_0+\frac{3\nu T}{\ell}M_*+\sqrt{E_0TD_Z}\\
 &\hspace{8mm}+\sqrt{(\ell^2D_Z+D_Y)D_Z}=:V_*.
\end{align*}
Direction density is zero at curl zeros; Fatou retains the nonzero-set
charge and may discard extra nonnegative zero-set dissipation. No weighted
$1/\sqrt2$ strain identity or evolved Schwartz persistence is assumed.
These are physical finite-horizon curl moment/tail and weighted direction
budgets, with no critical continuation consumer supplied.

The first curl moment gives continuous $\widehat\omega$. Distributional
$\operatorname{div}\omega=0$ along all radial directions implies
$\widehat\omega(0)=0$. Its first-moment increment and transverse inverse
curl yield $|\widehat u(\xi)|\le V_*$ almost everywhere off0. Exact radial
integration, with $d\xi$ retained, gives
\[
 \int_{|\xi|\le a}(2\pi|\xi|)^2|\widehat u|^2d\xi
 \le\frac{16\pi^3}{5}V_*^2a^5,
 \qquad
 \int_{|k|\le\kappa}|k|^2|\widehat u|^2d\xi
 \le\frac{V_*^2\kappa^5}{10\pi^2}.
\]
The earlier scout angular coefficient $4\pi/5$ was a valid loose upper
bound but an inaccurate conversion; the separate qualification preserves
it and uses $1/(10\pi^2)$. Centered zoom costs
$N_B=B^2(V_*+|x_B|M_*)$. At $z_B=1$, a low/high frequency split yields
\[
 q_B(0)\ge\tfrac12(5\pi^2)^{2/5}N_B^{-4/5}.
\]
For bounded centers this loses $B^{-8/5}$, asymptotically weaker than
$q_B\ge1/(BE_0)$. Both are compatible with $q_B\to0$ without a rate.

Finally, two original strong paths on an EXISTING compact have difference
$w_t+Aw=F$, $A=-\nu\Delta$,
$F=-\mathbb P[(u\cdot\nabla)w+(w\cdot\nabla)v]$.
Full $u=v+w$ retains the nonlinear term. Finite compact coefficients give
$\|F\|_2\le\alpha\|A^{1/2}w\|_2+\beta\|w\|_2$.
On $E=\|w\|_2^2>0$, put $N=\langle Aw,w\rangle/E$ and
$d=\|Aw-Nw\|_2/\sqrt E$. Direct differentiation gives
\[
 N'\le-2d^2+2d(\alpha\sqrt N+\beta)
 \le\alpha^2N+\beta^2.
\]
Its finite upper bound $N_*$ and
$(\log E)'\ge-2(N_*+\alpha\sqrt{N_*}+\beta)$ preclude a first later zero.
This proves backward uniqueness only on existing strong compacts.
The scalar pressure difference precedes projection; Leray kills its gradient.
No endpoint-uniform inverse stability follows.

For any fixed affine axis use
$v(x,t)=Q u(x_0+Q^{-1}(x-x_0),t)$ and scalar pressure composition.
Backward and forward uniqueness give global symmetry at an existing
positive time iff initial symmetry about the SAME fixed axis.
Actual positive-time spatial analyticity extends exact symmetry on a
nonempty invariant open core to connected whole space. Consequently exact
core plus outer asymmetry is excluded on original slices. The explicit
allowed datum $u_0=\operatorname{curl}(0,0,x_1e^{-|x|^2})$, with
$u_0(0)=-e_2$, refutes universal exact-core production about the prescribed
$z$-axis through0. It does not by that jet alone exclude every adaptive axis,
approximate cores or symmetric limiting classes after mass loss.
Those scopes require separate sources. The root, weighted critical consumer,
quantitative approximate-core leakage and normalized low-frequency/spatial
retention obligations remain open.

\section{Native finite horizon, initial Fourier source and global defect persistence}
\label{sec:native-finite-horizon}
The round63 five-body package has parent reconstruction, a distinct
adversarial pass and three independent entire-five written reviews.
Its original source and localized-retention duties remain open.
Put $S_0=4/\sqrt3$, $K=\nu^4/(16S_0^6)$ and $E=\|u\|_2^2$.
For $E_0>0$ native kinetic energy supplies the exact ledger
\[
 (E^2)'=-4\nu EZ,\qquad
 \int_0^t EZ=\frac{E_0^2-E(t)^2}{4\nu}.
\]
Set
\[
 H_c=\frac{E_0^2}{4\nu K}
     =\frac{16384E_0^2}{27\nu^5}.
\]
If $T_*\ge H_c$, an actual $s<H_c$ has $E(s)Z(s)\le K$.
Otherwise a positive surplus
$\delta=\int_0^b(EZ-K)>0$ on an earlier compact gives
at an EXISTING $t<H_c$ the ledger would give
\[
 E(t)^2\le E_0^2-4\nu Kt-4\nu\delta<0.
\]
This covers $T_*=H_c$ without evaluating its endpoint.

At that strong slice write $E_a=E(s)$, $Z_a=Z(s)$.
H\"older, Sobolev and interpolation give
\[
 |\langle\Delta u,(u\cdot\nabla)u\rangle|
 \le S_0^{3/2}(E_a Z)^{1/4}Q.
\]
While $Z\le2Z_a$ the coefficient is at most $2^{1/4}\nu/2<\nu$.
A first hit of $2Z_a$ is impossible since $Z'\le0$ beforehand.
Thus $Z\le Z_a$ and $Z(t)+\nu\int_s^tQ\le Z_a$ on the existing future.
The physical Wiener law then gives
\[
 \int_s^{T_*}W^2\le\frac{2Z_a}{\pi^2}
            \sqrt{\frac{T_*-s}{\nu}}<\infty
\]
if $T_*$ were finite. Original strong continuation excludes that endpoint.
Zero $E_a$ or $Z_a$ gives the zero future before division; zero datum is
separate. This restart uses finite-energy strong states, not evolved
Schwartz decay. Original initial Schwartz data and the all-order same-path
pressure adapter return every mixed closed-face trace and $E\le E_0$.
Consequently
\[
 T_*\ge H_c\quad\Longleftrightarrow\quad T_*=\infty.
\]
Universal survival through $H_c$ is UNPAID. The earlier coarse
$T_c=2H_c$ law follows directly by averaging $Z$ against the native energy
budget and remains a valid weaker sufficient interface.

For an independently priced finite consumer take $v=e^{\nu t\Delta}u_0$,
$r=\mathbb P[(v\cdot\nabla)v]$ and canonical $p_v$.
Since $\nabla p_v=-(I-\mathbb P)((v\cdot\nabla)v)$, the reference has
residual $+r$ and $w=u-v$, $w(0)=0$, satisfies
\[
 w_t-\nu\Delta w+
 \mathbb P[(v\cdot\nabla)w+(w\cdot\nabla)v+(w\cdot\nabla)w]=-r.
\]
For $z=\|\nabla w\|_2^2$, $q=\|\Delta w\|_2^2$,
$L=\|v\|_\infty+S_0\|\nabla v\|_3$, three Young allocations
$\nu q/4$ from $2\nu q$ yield
\[
 z'+\nu q\le az+bz^3+h,\qquad
 a=\frac{4L^2}{\nu},\quad b=\frac{108S_0^6}{\nu^3},\quad
 h=\frac{4\|r\|_2^2}{\nu}.
\]
Define reference-only finite charges
\[
 A(t)=\int_0^t a,\qquad H_T=\int_0^T e^{-A}h,\qquad
 K_T=b\int_0^T e^{2A}.
\]
If $8K_TH_T^2<1$, a first crossing of $y=e^{-A}z$ through $2H_T$
contradicts $y\le H_T+8K_TH_T^3<2H_T$.
The case $H_T=0$ uses zero-initial scalar uniqueness.
Thus $Z_u\le2Z_0+4e^{A(T)}H_T$ uniformly at existing times through $T$;
original squared-enstrophy continuation gives $T_*>T$.
At the SINGLE $T=H_c$ this returns globality with all original properties.
With $U_0=\|u_0\|_\infty$, $J_0=\|\nabla u_0\|_3$,
$L_0=U_0+S_0J_0$, a stronger coarse sufficient test is
\[
 13824S_0^6U_0^4Z_0^2\nu^{-5}T^3
       e^{8L_0^2T/\nu}<1.
\]
Same initial data removes only the initial difference; global residual,
pressure and nonlinear amplification remain charged. Nonzero residual makes
this particular test eventually fail as $T\to\infty$, which does not
refute its finite use. An arbitrary-data certificate at $H_c$ is open.

ONE accepted nonlinear analytic/Hgamma interval from original data has
\[
 h_0=\frac{\nu}{1024c_\gamma^2e^{1/2}W_0^2},\qquad
 c_\gamma=2^{\max(\gamma-1,0)}.
\]
Its constructed continuous strong endpoint is locally extendible, so
$T_*>h_0$. The proper initial class
\[
 E_0W_0\le\frac{\nu^3}{64c_\gamma e^{1/4}S_0^3}
\]
INCLUDING equality gives $H_c\le h_0$ and hence full global original return.
Zero $W_0$ is zero data before division. Amplitude multiplication grows
$E_0W_0$ cubically; NS scaling preserves it. No all-data inference follows.
To compare sufficient classes take a nonzero real even nonnegative smooth
unit-ball Fourier bump $\phi$, $I_j=\int\phi^j$ for $j=1,2$, and
\[
 \widehat u_{0,L}(\xi)=AL^{3/2}
 [\phi(L(\xi-k_0e_3))+\phi(L(\xi+k_0e_3))]\mathbb P(\xi)e_1,
 \qquad k_0L\ge4.
\]
Disjoint smooth transverse Hermitian supports away from zero give a real
solenoidal Schwartz inverse. Direct change of variables gives
\[
 \begin{gathered}
 \tfrac{16}{9}A^2I_2\le E_L\le2A^2I_2,\qquad
 W_L\le2AL^{-3/2}I_1,\\
 Z_L\ge4\pi^2(k_0-1/L)^2E_L,\qquad
 E_LW_L\le4A^3I_2I_1L^{-3/2}.
 \end{gathered}
\]
Choose $A,k_0$ first so that
$4\pi^2(3k_0/4)^2[(16/9)A^2I_2]^2>K$.
All eligible $L$ fail the earlier $E_LZ_L\le K$ class, while large $L$
satisfy the new class. This proves noncontainment in that direction,
without optimality, novelty or reverse containment.

The earlier exact-core counterdatum strengthens to every affine axis.
For $u_0=(-2x_1x_2,2x_1^2-1,0)e^{-|x|^2}$, coordinate-even kinetic density
has barycenter zero. An invariant density forces its barycenter onto the
axis. At zero $u_0=-e_2$, so that axis could only be the $e_2$ line.
The second component has Hessian entries $6$ and $2$ in the transverse
$x_1,x_3$ directions, violating quarter-turn invariance. Hence NO affine
axis works. Actual analytic identity and existing-compact backward
uniqueness exclude exact nonempty open full-vector axial cores at every
existing positive slice even with a slice-selected axis, fixed during the
comparison. Approximate/limiting cores and symmetric subclasses remain.

Finally, the actual uniform record analytic slab $J=[-H/4,0]$ pays
$\|U\|_\infty\le R_*$ and
$\|\nabla U\|_{\infty,F}\le\sqrt3R_*/\rho_J$,
$\rho_J=\sqrt{\nu H/4}$.
For any FIXED affine orthogonal polar action compare FULL $U$ to FULL
transformed solution $V$, with coherently transformed scalar pressure.
Put $\alpha=R_*/\sqrt\nu$, $\beta=\sqrt3R_*/\rho_J$, $\delta=H/4$.
Native $z\le1$ gives $\|\nabla(U-V)\|_2^2\le4$.
An independently supplied GLOBAL start defect energy at least $\eta>0$
therefore pays starting quotient $N_0\le4\nu/\eta$. With
\[
 N_\eta=(4\nu/\eta+\beta^2\delta)e^{\alpha^2\delta},
\]
the direct compact quotient law gives endpoint defect at least
\[
 \eta\exp[-2\delta(N_\eta+\alpha\sqrt{N_\eta}+\beta)].
\]
For genuine Haar projection $M$ about an affine axis and $v=U-M$,
$\int E_Q\,dQ=2\|v\|_2^2$, $E_Q\le4\|v\|_2^2$.
Strong continuity and compactness select one maximizing fixed rotation;
start $\|v\|_2^2\ge\eta$ yields
\[
 \|v(0)\|_2^2\ge\frac\eta2
   e^{-2\delta(N_{2\eta}+\alpha\sqrt{N_{2\eta}}+\beta)}.
\]
The Haar mean is not assumed unforced. Global endpoint fluctuation tending
to zero forces global starting fluctuation to zero; local weak symmetry
cannot supply the global hypothesis. The physical slab shrinks by $B^{-2}$.
Actual late witness production, spatial/low-frequency localization and
critical continuation returns remain open. No unconditional closure is
claimed by this finite package.

\section{Full reference consumers and the cost of local retention}
\label{sec:charged-reference-retention}
This reviewed package improves conditional consumers and actual record
accounting. It does not supply their arbitrary-data applicability.
Set $S_0=4/\sqrt3$. For continuous nonnegative $k,h$ and $p>1$, suppose
$y'\le ky^p+h$ on existing prefixes of a known interval $[0,T]$.
Put $H=y(0)+\int_0^T h$ and $K=\int_0^T k$. If $H>0$, then
\[
 KH^{p-1}\le\frac{(p-1)^{p-1}}{p^p}
 \quad\Longrightarrow\quad y(t)<\frac{pH}{p-1}.
\]
At a first crossing of $m=pH/(p-1)$ the integral bound is
$y(t)\le H+Km^p\le m$. At coefficient equality, either some coefficient
mass remains after the crossing or positive mass before it multiplies
$y(s)^p<m^p$; both give strict slack. Continuous coefficients have no
crossing atom. $K=0$ gives $y\le H<m$; $H=0$ gives the zero solution by
bounded-prefix comparison. Maximizing $(r-1)/r^p$ gives the stated
constant, optimal only for this constant-barrier packaging. For the
preceding cubic comparison this improves the old strict $1/8$ test to
the inclusive $K_{\rm old}H^2\le4/27$ test with barrier $3H/2$.

For any independently known smooth solenoidal reference $V$, retain its
full residual $r=V_t-\nu\Delta V+\mathbb P[(V\cdot\nabla)V]$ and mismatch.
Write $w=u-V$, $z=\|\nabla w\|_2^2$, $q=\|\Delta w\|_2^2$,
$L=\|V\|_\infty+S_0\|\nabla V\|_3$ and $R=\|r\|_2$. Then
\[
 z'+2\nu q\le2(L\sqrt z+R)\sqrt q+2S_0^{3/2}z^{3/4}q^{3/4}.
\]
Allocate $\nu q/2$ to each of the paired first term and the self term.
The first costs $4L^2z/\nu+4R^2/\nu$; maximizing $Ax^{3/4}-dx$ gives
$27A^4/(256d^3)$. Consequently
\[
 \begin{gathered}
 z'+\nu q\le az+b_{\rm new}z^3+h,\\
 a=4L^2/\nu,\quad h=4R^2/\nu,\quad
 b_{\rm new}=\frac{27S_0^6}{2\nu^3}=b_{\rm old}/8.
 \end{gathered}
\]
With $A(t)=\int_0^t a$, $y=e^{-A}z$,
$H=z(0)+\int_0^T e^{-A}h$ and $K_{\rm new}=\int_0^T b_{\rm new}e^{2A}$,
the function $g=H+\int_0^t b_{\rm new}e^{2A}y^3$ satisfies
$y\le g$ and $g'\le kg^3$. Thus
\[
 2K_{\rm new}H^2<1
 \quad\Longrightarrow\quad
 y(t)\le\frac{H}{\sqrt{1-2H^2\int_0^t k}}.
\]
$H=0$ is separate. Equality is not certified: the scalar equation
$y'=ky^3$ with constant $k>0$ can blow up at that boundary.
In old units the test is $K_{\rm old}H^2<4$. Its numerical threshold
is 32 times the old strict $1/8$, or 27 times the old inclusive $4/27$.
At the same coefficient, the Bihari threshold is four times the strict
constant-barrier threshold. These are comparisons of sufficient tests;
no strict enlargement of actual NS initial-data classes is proved.

An exact envelope keeps the simultaneous costs. For $C=S_0^{3/2}$ define
\[
 \begin{gathered}
 F_\nu(z,J)=\sup_{x\ge0}[-\nu x^4+2Cz^{3/4}x^3+2Jx^2],\\
 x_* =\frac{3Cz^{3/4}+\sqrt{9C^2z^{3/2}+16\nu J}}{4\nu}.
 \end{gathered}
\]
The positive critical point is the maximum, with the zero branch included.
$F_\nu(0,R)=R^2/\nu$ and
$F_\nu(z,0)=27S_0^6z^3/(16\nu^3)$. Integration by parts also gives
\[
 z'+\nu q\le2(2G+S_0D)z+F_\nu(z,R),\qquad
 G=\|\nabla V\|_{\infty,F},\quad D=\|\nabla^2V\|_3.
\]
The Hessian cost remains. The minimum of these valid upper envelopes is
valid. A positive finite smooth scalar supersolution bounds $z$ by local
comparison away from zero; no Lipschitz assertion at zero for fractional
powers is used. A scalar comparison ODE finite on the closed horizon is
another supplied certificate. Bounded $z$ and reference gradient give
original strong continuation through $T$; at $T=H_c$ the preceding full
same-path Clay return applies. Certificate production for all data is open.

An independently realizable reference is obtained by solving a known
linear problem. Let $v=e^{\nu t\Delta}u_0$ and construct
\[
 d_t-\nu\Delta d+\mathbb P[(v\cdot\nabla)d+(d\cdot\nabla)v]
      =-\mathbb P[(v\cdot\nabla)v],\qquad d(0)=0.
\]
For $m>5/2$, put $c_m=2^{\max(m-1,0)}$,
$C_m^2=\int\langle\xi\rangle^{-2m}$ and
$V_m=\sup\|v\|_{H^m}$. The heat-divergence norm is
$(2e\nu\tau)^{-1/2}$. Choose pieces with
\[
 \ell\le e\nu/(128c_m^2C_m^2V_m^2).
\]
The linear contraction factor is at most $1/2$.
With $F_m=4c_mC_mV_m^2\sqrt{\ell/(2e\nu)}$ and
$n=\lceil T/\ell\rceil$, the piecewise bound is $2F_m(2^n-1)$.
The zero case is split before division. Lower uniqueness identifies all
higher stages and spare orders give full time traces. For $V=v+d$ the
residual is exactly $\mathbb P[(d\cdot\nabla)d]$, not zero. Its norm
is at most $S_0^{3/2}\|\nabla d\|_2^{3/2}\|\Delta d\|_2^{1/2}$.
Fixed-time small-amplitude orders $d=O(\alpha^2)$ and $r=O(\alpha^4)$
do not control large amplitude at $H_c=O(\alpha^4)$.

A second consumer uses a full signed history. For known $v=ah$ with
$h=e^{\nu t\Delta}u_0$, keep $w_0=(1-a(0))u_0$ and
$R=a'h+a^2\mathbb P[(h\cdot\nabla)h]$. Let
$d=e^{\nu t\Delta}w_0-\int_0^t e^{\nu(t-s)\Delta}R(s)\,ds$.
Put $\kappa=2c_\gamma C_\gamma/\sqrt{2e\nu}$,
$M=\sup\|v\|_{H^\gamma}$, choose $\rho>0$, and set
\[
 \lambda=4\pi\kappa^2(2M+\rho)^2,\qquad
 \eta=\sup_{[0,T]} e^{-\lambda t}\|d(t)\|_{H^\gamma}.
\]
On a $\rho$ difference prefix, the discounted norm supremum satisfies
$Y\le\eta+Y/2$ since $\int_0^\infty\tau^{-1/2}e^{-\lambda\tau}
=\sqrt{\pi/\lambda}$. Hence the inclusive condition
$4e^{\lambda T}\eta\le\rho$ gives difference norm at most $\rho/2$;
the seed is at most $\rho/4$ and first exit is excluded. It returns
$T_*>T$, including equality in the test. The discount is outside the full
source history; seed, derivative, quadratic residual and nonlinear terms
remain. For $a=1$ it certifies the local interval
$1/(64e^2\kappa^2\|u_0\|_{H^\gamma}^2)$, not automatically $H_c$.

For a separate actual source, define $C(f)=\int ff^T$, $E=\operatorname{tr}C$.
Affine axial actions obey $C(T_Qf)=QC(f)Q^T$ without moments. If the least
adjacent eigenvalue gap is $d>0$, the axial Haar-average tensor has a
repeated eigenvalue and differs from $C$ by norm at least $d/2$.
Since tensor differences have norm at most $2\sqrt E\|f-T_Qf\|_2$,
every affine axis satisfies
\[
 \max_Q\|f-T_Qf\|_2^2\ge d^2/(16E).
\]
The real Schwartz curl of $(0,0,e^{-x_1^2-2x_2^2-3x_3^2})$ realizes
$C=\operatorname{diag}(2I,I,0)$, $E=3I$, $Z=28I$ and floor $I/48$,
$I=\pi^{3/2}/\sqrt{48}$. Existing-compact pair persistence propagates
this global witness, but gives no late uniform tensor gap or local mass.
For genuine Haar velocity $M$ and $v=u-M$, minmax and tensor factoring
similarly give $16E\|v\|_2^2\ge d^2$; $C(M)$ is not the average of $C(u)$.

At a fixed axis point $b$, bounded weights
$\phi_K=r^2/(1+r^2/K^2)$ have $|\nabla\phi_K|\le2\sqrt{\phi_K}$ and
$\Delta\phi_K\le6$. Actual canonical pressure obeys
$\|p\|_2\le\|u\|_\infty\sqrt E$. Weighted kinetic balance, followed by
remote cutoff removal at fixed $K$ and then monotone convergence, yields
\[
 J_b(t)\le E_0\left[\sqrt{J_b(0)/E_0+\ell^2}
              +3\nu t/\ell+3\int_0^t\|u\|_\infty\right]^2.
\]
Diffusion, convection and pressure have respective costs6,2,4.
Zero data is separate. Radial Haar orthogonality gives $J_v\le J_b$;
axis offsets are charged. For any $H^1$ field with positive energy and
finite $J=\int|x-b|^2|v|^2$, a cutoff at $R,2R$ gives
\[
 \|\nabla v\|_{L^2(B_{2R}(b))}^2
       \ge F(E,J,R)=\frac{(3E-7J/R^2)_+^2}{4J}.
\]
The loss7 is exterior mass3 plus collar4. Nonzero mass implies $J>0$.
$F$ increases in $E$ and decreases in $J$ on its positive branch, so
$E\ge\eta,J\le K$ gives the same bound with $\eta,K$.
$R^2\ge14K/(3\eta)$ gives $9\eta^2/(16K)$.
For a uniform limiting witness, every-index positive numerators are
insufficient: require a positive tail infimum of $F(\eta_B,K_B,R)$ at one
fixed $R$, bounded common balls and strong local derivative convergence.
A compact nonzero solenoidal $f$ with $v_n=\epsilon_n f$ has positive
numerators at suitable fixed $R$ but $F_n\to0$ and $v_n\to0$ strongly
in $H^1$. This substantive uniformity correction leaves actual late
source, scale-compatible moments, products and pressure return open.

Native records expose costs of other scales. Define $q_B=Q(t_B)/B^3$.
The native law $Z^2\le EQ$ gives $q_B\ge1/(BE_0)$.
The scale $\lambda=Q^{1/3}=Bq_B^{1/3}$ therefore tends to infinity
along divergent positive-time records, preserving growing histories.
It normalizes endpoint $Q$ to 1 and gives signed production at least
$\nu$, but gradient allowance $q_B^{-1/3}$ and fixed-window cost
$cR/q_B$ can diverge. Neither a product limit nor a local nonzero
witness follows.

For the growth-band scale keep the correct production
$P=+\langle\Delta u,(u\cdot\nabla)u\rangle$. At $Z>0$, $a=\sqrt{Q/Z}$,
$V=\|u\|_\infty$ and
\[
 D=2\nu(a-V/(2\nu))^2+2(Va-P/Z)\ge0,
 \qquad(\log Z)'=V^2/(2\nu)-D.
\]
For $0<\theta<1$ let $s_B$ be the last $\theta B$ crossing before an
actual record $t_B$. Put $L_B=E(s_B)-E(t_B)>0$,
$\delta_B=t_B-s_B$ and $M_B=\sup_{[s_B,t_B]}V$. Then
\[
 2\nu\theta B\delta_B\le L_B\le2\nu B\delta_B,
 \qquad M_B^2L_B/B\ge4\nu^2\theta\log(1/\theta).
\]
At hypothetical finite divergence $s_B\to T_*$ and $L_B\to0$, so
$M_B/\sqrt B\to\infty$. The scale $\lambda=\sqrt{B/L_B}$ has a
band duration between $1/(2\nu)$ and $1/(2\nu\theta)$ and positive
historical amplitude, with gradient/debit allowance $\sqrt{BL_B}$.
Writing $H_B=B^2\delta_B$ and $D_B=\int_{-H_B}^0q$ gives
\[
 \nu D_B+(1-\theta)/2\le C_0H_B^{1/4}D_B^{3/4}.
\]
All-fixed-depth starvation would contradict bounded $H_B$; hence in that
case $H_B\to\infty$, $BL_B\to\infty$, $H_B=o(B)$ and
$\sqrt B\ll\lambda\ll B$. The gradient allowance diverges and the
inherited analytic lower radius collapses. Fixed-depth limits cannot
bound the growing-depth cost. An independent growth-band amplitude upper
law or compatible source, strong topology, nonlinear products and pressure
return remains to be constructed. The unconditional target stays open.

\section{Canonical limits of bounded critical families}
\label{sec:critical-family-limit}

This section is conditional on an actual family, rather than a producer of
that family. Fix $\nu>0$, $T>0$ and real solenoidal strong seeds
$a_n\in L^2\cap L^3$, weakly converging to $a$ in $L^3$, with
$\|a_n\|_3\le M$. Supply canonical finite-energy strong NS paths $u_n$
on one common $[0,T)$ and their projected and local-energy identities.
The sizes of their initial kinetic energies may diverge. Neither the
common interval nor a bounded critical norm on late original slices is
derived here. The original unconditional Clay statement remains open.

For $h=G_{\nu t}*a$, put $\rho=|h|$ and $w=\rho^{3/2}$. The Hessian
of the $C^2$ function $|h|^3$ and cutoff exhaustion give
\[
 \frac{d}{dt}\int\rho^3
 =-3\nu\int\rho(|\nabla h|^2+|\nabla\rho|^2)
 \le-\frac{8\nu}{3}\int|\nabla w|^2.
\]
Zeros are treated by magnitude regularization and the weak chain rule.
Heat contraction, Sobolev $S_0=4/\sqrt3$ and interpolation yield
\[
 \int_0^\infty\!\int|h|^5
 \le\frac{3S_0^2}{8\nu}\|a\|_3^5=\frac{2}{\nu}\|a\|_3^5.
\]
Approximation in $L^3$, applying the same estimate to differences,
extends this law with $h\in C([0,\infty);L^3)$. Distributional divergence
commutes with heat; compact solenoidal approximation is unnecessary.
Weak $L^3$ seed convergence gives uniform convergence of every finite
heat derivative on positive-time compact boxes: the translated Gaussian
derivatives form a compact test set in the dual $L^{3/2}$. This does not
give strong convergence of the seeds or nonlinear smooth initial jets.

A separate scalar staging law keeps its PDE premise explicit. Suppose
$\alpha(0)=0$, $\alpha$ is continuous and nonnegative, $\beta,g\ge0$,
$\int g\le G$, and
\[
 \alpha(t)+\beta(t)\le C\left[t^\theta+
      \int_0^t\{\alpha(1+g)+\alpha^3\}\right].
\]
While $\alpha\le b$, an integrating factor for
$y=\int(1+b^2+g)\alpha$ gives
\[
 \alpha(t)+\beta(t)\le Ct^\theta
                 e^{C[(1+b^2)t+G]}.
\]
The horizon $S\le\min\{1,[b/(2Ce^{C(1+b^2+G)})]^{1/\theta}\}$
excludes a first crossing of $b$, including equality in this test.
$C=0$ is separate. For a heat-reference local estimate the coefficient
is $g=\|h\|_5^5$ and its integral has the preceding price. The full local
PDE inequality is still supplied in this scalar law; it is not imported
from a printed coefficient with power two.

There is also a direct global consumer in the stated strong energy class.
Set $v=u-h$, $Y=\|v\|_2^2$, $D=\|\nabla v\|_2$ and $W=\|h\|_5$.
All four difference tensors and canonical pressure remain in the equation.
Legal global transport and pressure pairings give
\[
 Y'+2\nu D^2=2\int[h\otimes(h+v)]:\nabla v.
\]
The mixed and heat self-work are retained. Interpolation and Young give
\begin{align*}
 2\left|\int(h\otimes v):\nabla v\right|
   &\le\tfrac\nu2D^2+64S_0^3\nu^{-4}W^5Y,\\
 2\left|\int(h\otimes h):\nabla v\right|
   &\le\tfrac\nu2D^2+2\nu^{-1}\|h\|_4^4.
\end{align*}
The exact first Young coefficient is
$(32/5)(8/5)^4S_0^3<64S_0^3$. Heat interpolation in space and time pays
$\int_0^T\|h\|_4^4\le\sqrt2 M^4\sqrt{T/\nu}$. Define
\[
 K=128S_0^3M^5/\nu^5,\qquad
 F_T=2\sqrt2 M^4\sqrt T\,\nu^{-3/2},\qquad
 D_T=F_Te^K.
\]
The integrating factor and the accumulated $kY$ charge prove, at every
common existing time,
\[
 Y(t)+\nu\int_0^tD(s)^2\,ds\le D_T.
\]
Each energy and dissipation cap is $D_T$. The sum of the maximum energy
and the full dissipation can exceed this cap.
The sharper individual bound is $Y(t)\le D_t=O(\sqrt t)$, and
$\int|v|^{10/3}\le S_0^2D_T^{5/3}/\nu$ on the existing interval.
These prices are independent of the magnitude of each finite seed energy,
but do not remove finite-energy applicability or create a strong path.

Uniform unit-ball energy $A$ prices a far pressure difference directly.
For $R\ge1$, $x\in B_R(b)$ and
$K_{ij}(z)=3z_iz_j|z|^{-5}-\delta_{ij}|z|^{-3}$, set
\[
 p_{\rm far}(x)=\frac1{4\pi}\int_{|y-b|\ge2R}
 [K(x-y)-K(b-y)]:(u\otimes u)(y)\,dy.
\]
The Frobenius derivative cost is $21+3\sqrt3\le27$; the segment price
is $432|x-b||y-b|^{-4}$. Cubes of side $1/\sqrt3$ and dyadic shells
give $\int_{|y-b|\ge s}|y-b|^{-4}|u|^2\le2C_{\rm cov}A/s$, where
$C_{\rm cov}=(4\sqrt3+3)^3$. Hence
\[
 |p_{\rm far}(x)|\le(108C_{\rm cov}/\pi)A|x-b|/R,
 \quad |\nabla p_{\rm far}(x)|\le(108C_{\rm cov}/\pi)A/R.
\]
This absolutely convergent field is harmonic. An outer source radius $L$
costs $216C_{\rm cov}RA/(\pi L)$ per tail. Strong local $L^2$ velocity
convergence identifies the full tensor in local $L^1$; pass a fixed
annulus before sending $L$ to infinity. This gives uniform spatial tail
passage, or $L^1$ in time of the spatial supremum in the spacetime version.
Weak velocities alone can leave an oscillatory tensor defect.

Near pressure is a separate law. Distributional differentiation gives
\[
 \partial_i\partial_j(4\pi|x|)^{-1}
       =(4\pi)^{-1}{\rm pv}\,K_{ij}-\delta_{ij}\delta_0/3.
\]
Its multiplier is bounded on $L^2$. Dyadic good/bad decomposition, the
mean-zero bad-cube cancellation and the $|x|^{-4}$ derivative bound prove
weak $(1,1)$. Splitting input at the output height and integrating its
distribution function gives strong $L^p$, $1<p<2$, with finite factors
$1/(p-1)$ and $1/(2-p)$. Density and finite bad sums followed by $L^2$
convergence justify the limits. This proves the needed $L^{3/2}$ estimate
for the full near tensor including its contact term.

The remainder and a local heat cutoff now give uniform unit-ball energy
$E$ and integrated gradients $D_0$, with spatial supremum outside the
gradient time integral. A bounded reflection extension and local Sobolev
give $\int_{B_R\times(0,T)}|u_n|^{10/3}\le C_RE^{2/3}(D_0+TE)$.
Near pressure and the far difference therefore have uniform local
$L^{3/2}$ spacetime prices modulo scalar time constants. Normalize every
pressure by one $B_1$ average, consistently on nested balls; define
$P_R=\sup_n\int_0^T\int_{B_R}|p_n-\operatorname{avg}_{B_1}p_n|^{3/2}$.
This retains physical pressure and excludes arbitrary affine additions.

For the projected identity the Oseen derivative obeys
\[
 \widehat D_{ij\ell,h}(k)=ik_\ell e^{-h|k|^2}P_{ij}(k),
 \qquad |D_h(x)|\le C(\sqrt h+|x|)^{-4}.
\]
Its local $L^r$ age exponent is $a_r=2-3/(2r)$. Recent ages below
$\eta$ cost $C_RE[\nu^{-a_r}\eta^{1-a_r}/(1-a_r)+\eta]$ for
$1\le r<3/2$. Smooth old-age cutoffs allow derivatives of the kernel and
cutoff only, never source-time derivatives. Arzel\`a--Ascoli plus recent
vanishing gives $C([\delta,T];L^r_{\rm loc})$ convergence, with terminal
traces from below. The endpoint $r=3/2$ is not supplied. Interpolation
with local $10/3$ and an initial-strip $L^3$ norm cost $\delta^{1/30}$
gives strong $L^3$ on whole bounded spacetime cylinders. Full tensors and
cubic flux pass strongly; pressure times velocity passes weak--strong.

For every compact smooth vector test $\psi$, the actual PDE bounds time
pairings by $C_\psi[(\nu\sqrt E+E)|t-s|+P_R^{2/3}|t-s|^{1/3}]$.
Dense-test compactness yields one every-slice weakly local-$L^2$ continuous
representative, including both endpoints. The remainder bound and heat
trace give strong initial local-$L^2$ attainment with remainder norm
$O(t^{1/4})$. Pass the projected identity by fixing old age and source
radius, passing local tensors, removing the $E/L$ tail and then the recent
age error. Pass near and far pressure with the same single normalization.
A local PDE alone would not identify these canonical objects.

Local-energy tests supported after zero pass by weak lower semicontinuity
on terminal energy/dissipation and the preceding products. To recover an
initial term, apply a rising time switch in the \emph{limit}; its derivative
has unit mass near zero. Proved strong initial attainment makes its term
converge to $\int|a|^2\phi(0)$. All other terms pass by integrability.
Weak initial energies of $a_n$ are never passed to the right-hand side.
This constructs the canonical local-energy limit with its correct initial
term. It supplies neither nonzero singular retention nor terminal zero,
exterior backward uniqueness, a uniform late critical seed source, original
$H^\gamma$ continuation or the full mixed original pressure/velocity face.

\section{A fixed-input test of critical pressure closure}
\label{sec:critical-pressure-source-control}

The preceding bounded-critical limit construction requires an independently
supplied late critical mass. We now test a specific original-source proposal,
with full canonical pressure and actual Schwartz initial data. Write
\[
\begin{aligned}
 &I(u)=\int |u|^3,\qquad Z(u)=\int |\nabla u|^2,\\
 &q(u)=\operatorname{div}(|u|u),\qquad P(u)=\int p_u q(u),\\
 &D(u)=\int |u|\bigl(|\nabla u|^2+|\nabla|u||^2\bigr).
\end{aligned}
\]
The vector $C^2$ chain rule for $|z|^3$, cutoffs and the original
unforced equation give $I'=3P-3\nu D$ on existing strong intervals.
A proposed bound $I'\le C(\nu,E_0,I_0)Z(1+I)$ would give
$I(t)+1\le(I_0+1)\exp(CE_0/(2\nu))$ using native energy.

\begin{proposition}
For each $\nu>0$ there are positive $E_*,I_*$ and real compact smooth
solenoidal data $U_\lambda$, $\lambda\ge1$, such that
$\|U_\lambda\|_2^2=E_*$, $\|U_\lambda\|_3^3=I_*$ and their actual
local original strong solutions satisfy
\[
 \frac{I_\lambda'(0)}{Z(U_\lambda)(1+I_*)}\longrightarrow+\infty.
\]
Consequently no finite constant depending only on $\nu,E_0,I_0$
can supply the stated universal differential closure.
\end{proposition}
\begin{proof}
First construct a compact solenoidal field with positive pressure work.
On the $2\pi$ torus, put $k=(1,0,0)$, $\ell=(0,1,1)$, $h=k+\ell$ and
\[
 v=(\cos(\ell\cdot x)+\sin(h\cdot x),\ \cos(k\cdot x),\ -\sin(h\cdot x)).
\]
Its mean-zero pressure has coefficients
$1/3,-1/3,-1/6,1/2,1/9,-1$ at, respectively,
$\cos((k-\ell)\cdot x),\cos(h\cdot x),\sin((2k+\ell)\cdot x),
\sin(\ell\cdot x),\sin((k+2\ell)\cdot x),\sin(k\cdot x)$.
Direct Fourier multiplication gives
$\langle p_v\partial_3v_3\rangle=1/6$ and $\|p_v\|_\infty\le22/9$.
Since $|v|\le3$, $|\nabla v|_F\le4$, for $C\ge6$ and $0\le a\le1$,
\[
 |q(Ce_3+av)-Ca\partial_3v_3|\le84a^2,\qquad
 \langle a^2p_v q(Ce_3+av)\rangle
 \ge Ca^3/6-(616/3)a^4.
\]
With $C=2048$ this is at least $136a^3$; the pressure envelope $a^2$
already appears in this work.

Let $A$ be the bounded periodic vector potential of $v$. Choose
$0\le\psi\le1$, nonzero and smooth with support in $B_1$, and radial
$\chi$ equal to one on $B_2$ and zero outside $B_3$. Set
\[
 b_R=\nabla\times[(e_3\times x)\chi(x/R)/2],\quad
 v_R=\nabla\times[\psi(x/R)A(x)],\quad W_R=Cb_R+v_R.
\]
Both fields are compact solenoidal. The boost is exactly $e_3$ on the
support of $v_R$, so the double divergence of their cross tensor is
zero. Thus $p_{W_R}=C^2p_b(x/R)+p_{v_R}$. Inversion parity gives
$P(b_R)=0$. Pair the slow boost pressure with the compact difference
flux by integration by parts; its gradient supplies $R^{-1}$, giving
$O_C(R^2)$. The remaining slow-boost pairing is also $O_C(R^2)$ by
$L^2$ bounds.

For each nonzero fast tensor mode, the full pressure multiplier
$H_{ij}(\xi)=-\xi_i\xi_j/|\xi|^2$ is bounded and continuous there.
Fourier shift and dilation therefore give normalized $L^2$ convergence
of its modulated envelope to its value at that mode. The zero tensor
mode gives a slow pressure $H_{ij}[T_{0,ij}\psi^2](x/R)$, which must be
retained. It pairs to zero after normalization with the fast periodic
$q(Ce_3+\psi(X)v)$, whose mean is zero for every $X$; Fourier averaging,
compact envelope support, local $L^2$ slow pressure and summable fast
Fourier tails justify this statement. The curl-cutoff error is
$O(R^{-1})$ in the normalized product bounds. Hence
\[
 \lim_{R\to\infty}R^{-3}P(W_R)
 =\int \psi(X)^2\langle p_vq(Ce_3+\psi(X)v)\rangle\,dX
 \ge136\int\psi^3>0.
\]
Choose a finite such $R$. Amplitude scaling gives
\[
 P(A_0W_R)=A_0^4P(W_R),\qquad D(A_0W_R)=A_0^3D(W_R).
\]
It supplies a compact field $V$ with $\delta_V=P(V)-\nu D(V)>0$.

Take any nonzero compact smooth solenoidal $g$. Write
$E_*=E_V+E_g$, $I_*=I_V+I_g$, and
\[
 L_\lambda=\frac{E_*-E_V/\lambda}{E_g},\qquad
 U_\lambda(x)=\lambda V(\lambda x)
       +L_\lambda^{-1}g((x-d)/L_\lambda).
\]
One fixed distant $d$ separates the supports for every $\lambda\ge1$,
since $1\le L_\lambda\le1+E_V/E_g$. Direct integration gives the exact
fixed norms, and
$Z(U_\lambda)=\lambda Z(V)+L_\lambda^{-1}Z(g)$,
$D(U_\lambda)=\lambda^2D(V)+L_\lambda^{-2}D(g)$.
Velocity cross stress is zero, but both pressure cross pairings remain.
Subtract the background pressure at the core origin using $\int q_V=0$;
its bounded gradient and the core first moment give $O(\lambda^{-1})$.
For the reverse pairing, the off-support Newton Hessian is bounded by
$\sqrt6/(4\pi r^3)$, the core tensor mass is $E_V/\lambda$ and the
background $q$ has invariant $L^1$ norm. Fixed separation gives another
$O(\lambda^{-1})$. Consequently
\[
 P(U_\lambda)=\lambda^2P(V)+L_\lambda^{-2}P(g)+O(\lambda^{-1}),\qquad
 I_\lambda'(0)=3\lambda^2\delta_V+O(1).
\]
Each datum has its own actual local original strong path; no common
lifespan is needed. Divide by $[\lambda Z(V)+O(1)](1+I_*)$ to conclude.
Continuity also gives a positive-time violation if exactly $t=0$ is
omitted from the proposed bound.
\end{proof}

This removes the particular norm-only supplier from the route. It proves
no late critical cap and no singular trajectory. Constants using the
full datum or additional native source information, and independently
proved signed temporal compensation, remain possible research targets.
Their source and original continuation returns remain open. Novelty of
this exact construction has not been checked against matched prior work.

\end{document}
